\providecommand{\phio}{\phi(0)}
\providecommand{\cR}{\mathcal R}
\providecommand{\Prb}{\mathbb P}
\section*{How to read these estimates}
The main proof uses one attained ratio maximizer and the contact inequality
at coefficient $C=C_n$. The first section bounds its support before any
higher-moment calculation. The resource and Fourier sections then acquire
simultaneous information about that same law. The local comparison section
uses only the acquired information and the actual iid binomial count.
The final refinements retain a joint resource, correct the sign of the cubic
term, and use the signed Gaussian phase. They provide the analytic licenses
for the finite-index implementation. Numerical domain coverage and arithmetic
acceptance remain separate from these analytic proofs.

The contact interface used below is precisely the normalized fixed-target
contact inequality at $c=C_n$, supplied by the main manuscript's ratio
attainment corollary. It is stated explicitly before its use. The Bernoulli
theorem and the finite moment caps are the only external inequality inputs
to this selected local analytic argument.


% Source: support.tex
\section{One support estimate, uniform in the threshold}\label{support:theorem}
\subsection{The scale correction}
Write $\phi_0=\phi(0)$.
For $n>1$ put $s_n=\sqrt{n/(n-1)}$ and
$e_n=n(s_n-1)-1/2$.
\begin{lemma}\label{support:scale}
For every integer $n\ge36$, one has $0<e_n<.384/n$.
\end{lemma}
\begin{proof}
Put $x=1/n$ and $c=.384=48/125$. Direct expansion gives
\[
 (1-x)(1+x/2+cx^2)^2-1
 =x^2\{(2c-3/4)-(c+1/4)x+(c^2-c)x^2-c^2x^3\}.
\]
The bracket is decreasing for $x\ge0$, since $0<c<1$.
At $x=1/36$ its exact rational value is $823/4050000>0$.
Squaring positive quantities gives
$s_n<1+x/2+cx^2$ for $0<x\le1/36$. Similarly
$(1-x)(1+x/2)^2-1=-3x^2/4-x^3/4<0$, which gives the lower bound.
\end{proof}


\subsection{A signed Gaussian envelope}
Define, for every real $H,t$ and $n>1$,
\[
 \mathcal H_n(H,t)=n^{3/2}\{\Phi((H-t)/\sqrt{n-1})-
             \Phi(H/\sqrt n)\}
       +nt\phi(H/\sqrt n)+\tfrac H2(t^2-1)\phi(H/\sqrt n).
\]
Write $v=n-1$, $s=\sqrt{n/v}$, $u=|t|$, and
\[
 A=\frac{u^3-3u}{6},\quad B=\frac{u^4-6u^2+3}{24\sqrt v},
 \quad D=\frac{u^5-15u}{120v}.
\]
\begin{lemma}[Fifth-order identity with a uniform finite remainder]
If $u>\sqrt3$, then, with $z=H/\sqrt v$,
\[
 \frac{\mathcal H_n(H,\pm u)}{s^3}
 =\pm A(1-z^2)\phi(z)+B\phi^{(3)}(z)
                  \mp D\phi^{(4)}(z)+R,
\]
where the sign of the remainder need not be fixed and
\[
 |R|\le E^{(5)}_v(u):=
 \frac{51u^6}{14400v^{3/2}}
 +\frac{51(3u^2-2)}{960v^{3/2}}+\frac{u}{8v^2}.
\]
This statement is uniform over every real $H$.
\end{lemma}
\begin{proof}
Let $F_w(x)=\Phi(x/\sqrt w)$, with primes denoting $x$ derivatives.
The heat identity $\partial_w F_w=F_w''/2$ and ordinary Taylor's theorem
give, with unlabelled derivatives at $(v,H)$,
\begin{align*}
 F_{v+1}&=F_v+F_v''/2+F_v^{(4)}/8+R_0,
 &|R_0|&\le\sup_{w\in[v,v+1]}\|F_w^{(6)}\|_\infty/48,\\
 F_{v+1}'&=F_v'+F_v'''/2+F_v^{(5)}/8+R_1,
 &|R_1|&\le\sup_{w\in[v,v+1]}\|F_w^{(7)}\|_\infty/48,\\
 F_{v+1}''&=F_v''+F_v^{(4)}/2+R_2,
 &|R_2|&\le\sup_{w\in[v,v+1]}\|F_w^{(6)}\|_\infty/8.
\end{align*}
Expand $F_v(H-t)$ through degree five, with remainder at most
$u^6\|F_v^{(6)}\|_\infty/720$. In
\[
 \frac{\mathcal H_n}{n^{3/2}}=
 F_v(H-t)-F_{v+1}(H)+tF_{v+1}'(H)
             -\frac{t^2-1}{2}F_{v+1}''(H),
\]
the first three derivative coefficients cancel. The next three are
\[
 \left(\frac t2-\frac{t^3}6\right)F_v'''
 +\frac{t^4-6t^2+3}{24}F_v^{(4)}
 +\left(\frac t8-\frac{t^5}{120}\right)F_v^{(5)}.
\]
The remaining absolute value is bounded by
\[
 \left\{\frac{u^6}{720}+\frac1{48}+\frac{u^2-1}{16}\right\}
       \sup_{w\in[v,v+1]}\|F_w^{(6)}\|_\infty
 +\frac u{48}\sup_{w\in[v,v+1]}\|F_w^{(7)}\|_\infty.
\]
Fourier inversion of the Gaussian, differentiated under its absolutely
integrable Fourier transform, gives
\[
 \|\phi^{(k)}\|_\infty
 \le\frac1\pi\int_0^\infty y^k e^{-y^2/2}\,dy.
\]
Integration by parts yields
$\|\phi^{(5)}\|_\infty\le8/\pi<2.55$ and
$\|\phi^{(6)}\|_\infty\le15\phi_0<6$.
Since $F_w^{(j)}=w^{-j/2}\phi^{(j-1)}$, the displayed remainder
becomes $s^3E^{(5)}_v(u)$ after multiplication by $n^{3/2}$.
The three retained derivative coefficients give the identity.
\end{proof}

\begin{lemma}[A lower tangent bound for the fourth Gaussian derivative]
For every real $z$,
\[
 \phi^{(4)}(z)\ge\phi_0(3-15z^2/2).
\]
\end{lemma}
\begin{proof}
The same Gaussian Fourier integral gives
$\phi^{(4)}(z)=\pi^{-1}\int_0^\infty
y^4e^{-y^2/2}\cos(yz)\,dy$.
Use the globally valid inequality $\cos w\ge1-w^2/2$ inside
this integral. Its fourth and sixth Gaussian moments are $3\phi_0$
and $15\phi_0$, respectively.
\end{proof}

\begin{lemma}[Uniform signed fifth-order envelope]\label{support:envelope}
Suppose $u>\sqrt3$ and $A,B,D\ge0$. If
$W=A-15D/2>0$, put
\[
 T=A-3D+\frac{9B^2}{4W}.
\]
If $T\ge A/2+2B+3D$, then, for every real $H$ and either sign,
\[
 \mathcal H_n(H,\pm u)\le s^3\{\phi_0 T+E^{(5)}_v(u)\}.
\]
\end{lemma}
\begin{proof}
For positive $t=u$ and $|z|\le1$, divide the retained terms by
$\phi_0$ and use the preceding tangent bound to obtain
\[
 A-3D-(A-15D/2)z^2+3B|z|
 \le A-3D+9B^2/(4W).
\]
Here $(1-z^2)\phi(z)\le(1-z^2)\phi_0$ and
$|\phi^{(3)}(z)|\le3|z|\phi_0$ on this interval.
For $|z|\ge1$ the first term is nonpositive, and the Fourier bounds
$\|\phi^{(3)}\|_\infty<2\phi_0$ and
$\|\phi^{(4)}\|_\infty\le3\phi_0$ give $2B+3D$.
For negative $t=-u$, the first term has global upper bound
$2e^{-3/2}A\phi_0<A\phi_0/2$, and the other terms give
$(2B+3D)\phi_0$. The hypothesis on $T$ covers both exterior cases.
The inequality $2e^{-3/2}<1/2$ follows from four positive terms of
the exponential series at $3/2$. Finally add the uniform remainder.
\end{proof}

\begin{lemma}[A global Gaussian remainder bound]
\label{ext:global-gaussian-ray}
For every integer $n\ge2$ and all real $H,t$,
\[
 |\mathcal H_n(H,t)|
 <\frac{|t|^3}{15}+\frac{|t|}{5}+\frac14
 <\frac{|t|^3}{15}+\frac{|t|}{5}+\frac12.
\]
\end{lemma}
\begin{proof}
Write $F_w(x)=\Phi(x/\sqrt w)$, $v=n-1$ and $y=H-t$.
The identity defining $\mathcal H_n$ is
\[
 n^{-3/2}\mathcal H_n(H,t)
 =F_v(y)-F_n(H)+tF_n'(H)-\frac{t^2-1}{2}F_n''(H).
\]
Taylor expansion of $F_n(y)$ at $H$ through degree two has remainder
at most $\phi(0)|t|^3/(6n^{3/2})$, since
$\|\phi''\|_\infty=\phi(0)$.
After this expansion the remaining correction is
\[
 \frac12\{F_n''(H)-F_n''(y)\}
 +\frac12\int_v^n\{F_n''(y)-F_w''(y)\}\,dw.
\]
Here we used $\partial_wF_w=F_w''/2$. The first term is at most
$\phi(0)|t|/(2n^{3/2})$ in absolute value. Applying the heat identity
once more bounds the second term by
\[
 \frac14\int_v^n\int_w^n
       |F_r^{(4)}(y)|\,dr\,dw
 \le\frac{\|\phi'''\|_\infty}{8v^2}
 \le\frac1{4\pi v^2}.
\]
The last inequality follows by differentiating the absolutely
convergent Gaussian Fourier integral, which gives
$\|\phi'''\|_\infty\le2/\pi$. Multiply by $n^{3/2}$ and use
$\phi(0)<2/5$. The function $x^{3/2}/(x-1)^2$ decreases for $x>1$,
because its logarithmic derivative is
$3/(2x)-2/(x-1)=-(x+3)/(2x(x-1))<0$. Hence, for $n\ge2$,
\[
 \frac{n^{3/2}}{4\pi(n-1)^2}
 \le\frac{\sqrt2}{2\pi}<\frac14,
\]
using $\sqrt2<3/2$ and $\pi>3$. This proves the stronger bound and
therefore the exact ray majorant used by the original checkers.
\end{proof}

\subsection{Finite bands and the exact support polynomial}\label{support:polynomial}
For integers $36\le N\le L$, use the simultaneous endpoint bounds
\[
 B_+=\frac{u^4-6u^2+3}{24\sqrt{N-1}},\quad
 D_+=\frac{u^5-15u}{120(N-1)},\quad
 D_-=\frac{u^5-15u}{120(L-1)}.
\]
For an infinite index range, set $D_-=0$ instead. Put
\[
 W_*=A-15D_+/2,\qquad
 T_*=A-3D_-+\frac{9B_+^2}{4W_*}.
\]
The proof of Lemma~\ref{support:envelope} applies uniformly on the whole band
if $W_*>0$ and $T_*\ge A/2+2B_++3D_+$. In particular the favorable
term is frozen at the upper index endpoint, and the adverse terms at
the lower one. No monotonicity of a signed coefficient is assumed.

Let $(P,H)$ be an attained positive closed maximizing pair, standardized
to mean zero and variance one. Write $C=C_n\ge C_E$,
$\rho=\mathbb E|X|^3$, $\beta=\mathbb E[X|X|]$, and
$d=n(C_{n-1}-C_n)_+$. The finite contact/predecessor bound is
\[
 \rho I(t)\le\mathcal H_n(H,t)+
 C\{-|t|^3+3\rho t^2/2+3\beta t+\rho e_n\}+\rho s_nd,
\]
with $I=0$ on the entire closed support. Suppose also, for a finite
band, that $C_n-C_E\ge K/n$, where $K>0$. The finite moment bounds of~\cite[p.~125]{Shevtsova2013}, namely
$C\rho\le .3031(\rho+.646)$ and
$C\rho\le .3322(\rho+.429)$, give
\[
 c=.4095+K/L<C_n,\qquad
 \rho<r=\min\left\{\frac{.3031(.646)}{c-.3031},
                        \frac{.3322(.429)}{c-.3322}\right\}.
\]
For the unrestricted nonnegative case put $K=0$, $c=.4095$.
The sign-moment inequality
$\rho\ge(1+\beta^2)/\sqrt{1-\beta^2}$ supplies $|\beta|<b$ whenever
$(1+b^2)^2>r^2(1-b^2)$ and $0<b<1$. All these statements concern the
same current law. The lower endpoint cannot replace $L$ in $K/L$.

When a preceding resource enclosure supplies sharper actual bounds
$\rho\le r$ and $|\beta|\le b$, these may instead be used in the
same polynomials. Their validity is the only extra premise.

Use rational $\ell^2\le N-1$ and $S^2>N/(N-1)$. For $u=|t|$ define
\[
 Q(u)=-u^3+3ru^2/2+3bu+.384r/N.
\]
First check $Q<0$, before replacing $C$ by its lower bound $c$.
Replace $1/\sqrt{N-1}$ by $1/\ell$ in the adverse remainder and in
the term $2B_+$ of the domination gate. Write $E_*$ for that rational
remainder. Strict negative contact for every $n$ in the band follows
from the polynomial inequality
\[
 \frac25S^3\left\{(A-3D_-)W_*+
              \frac{(u^4-6u^2+3)^2}{256(N-1)}\right\}
 +S^3E_*W_*+\{cQ+rSq\}W_*<0.
\]
The second polynomial gate is $W_*>0$. The third is
\[
 \{A/2-3D_--3D_+-2B_+\}W_*+
         \frac{(u^4-6u^2+3)^2}{256(N-1)}>0.
\]
This is the cleared-denominator version of the domination hypothesis.
The bound on $T_*$ is positive, so the replacement $\phi_0<.4$ and
$s_n<S$ both have their correct direction.

The exact kernel tests these polynomials on all adjacent closed intervals
of length $.05$ from $R$ to $4.5$. Positive Bernstein coefficients prove
their signs on every point of each cell, including endpoints. On the
rays $u\ge4.5$ use Lemma~\ref{ext:global-gaussian-ray}, which gives the same bound
$|\mathcal H_n(H,t)|\le u^3/15+u/5+1/2$.
The resulting cubic and its first three derivatives are strictly negative
at $4.5$, as is the coefficient $Q$ before replacing $C$ by $c$.
This excludes the whole rays. Thus, if $d\le q$, actual support is in
$(-R,R)$, including exclusion of arbitrarily small exterior atoms and
exterior accumulation points.


\subsection{The support rows used in this proof}
For the all-index local argument, take $K=0$, $R=3.84$, $q=.19$ on
\[
[45,49],\ [50,59],\ [60,89],\ [90,134],\ [135,199],\ [200,\infty).
\]
Only indices at least 47 are used. The last row uses $D_-=0$;
there is no upper index restriction. The exact checker reconstructs
all closed $u$ cells and the two exterior rays. Missing an index band
or the infinite row is an acceptance failure.

\subsection{One bounded-threshold series}\label{support:series}
Put $c_k=(-1)^k/[2^kk!(2k+1)]$, and let $(x)_j$ denote the rising
factorial. Define the following rational polynomials:
\begin{align*}
 P_j(H,t)={}&\sum_{k=0}^{j+1}
 \frac{c_k(k+1/2)_{j+1-k}}{(j+1-k)!}(H-t)^{2k+1}
       -c_{j+1}H^{2j+3}\\
 &+\frac{(-1)^{j+1}tH^{2j+2}}{2^{j+1}(j+1)!}
       +\frac{(-1)^jH^{2j+1}(t^2-1)}{2^{j+1}j!}.
\end{align*}
In particular,
\[
 P_0=(t^3-3t)/6,\qquad
 P_1=\frac{3t-t^3}{4}H^2+\frac{t^4-6t^2+3}{8}H
                  +\frac{-t^5+10t^3-15t}{40}.
\]
Fix $h>0$, $E(x)=(193/71)^{\lceil x\rceil}$, and
\[
 B_r=\frac{7(h+r)}6E((h+r)^2/6)
       +\left(h+r+\frac{h(r^2+1)}8\right)E(h^2/8).
\]
\begin{lemma}\label{support:finite-series}
For every $n\ge4$, $|H|\le h$ and $|t|\le r$,
\[
 \left|\mathcal H_n/\phi(0)-\sum_{j=0}^5P_j/n^j\right|
       \le4^7B_r/n^6.
\]
\end{lemma}
\begin{proof}
Put $z=1/n$. Integration of the Gaussian exponential series gives
$\mathcal H_n/\phi(0)=S(z)/z$, where
\begin{align*}
 S(z)={}&\sum_{k\ge0}c_k\{(H-t)^{2k+1}(1-z)^{-k-1/2}
                                  -H^{2k+1}\}z^k\\
 &+te^{-H^2z/2}+\tfrac12zH(t^2-1)e^{-H^2z/2}.
\end{align*}
Its constant coefficient vanishes, and expansion of the binomial factor
gives $[z^{j+1}]S=P_j$. At $z_0=1/4$, the sum obtained by replacing
all coefficients by their absolute-value majorants is at most $B_r$:
use $(1-z_0)^{-1/2}<7/6$, $1/(2k+1)\le1$, and $e<193/71$.
The latter bound follows, for example, by summing the exponential series
through degree six and bounding the rest by a geometric series with
ratio $1/8$. All products and rearrangements are therefore absolutely
convergent. The discarded terms have degrees at least seven in $S$,
so their sum at $0\le z\le z_0$ is at most $(z/z_0)^7B_r$.
Divide by $z$ to get the result. This proves a finite uniform bound,
including every threshold in the stated interval and every $n\ge4$.
\end{proof}

\subsection{A common contact test for radius and strip refinement}
\label{support:panels}
Suppose $n\ge47$, $0\le C-C_E\le.7/n$, $\delta_n\le.19$ and
$|X|\le R_0$. A previously acquired row gives
\[
 \rho\le r,\quad 0\le b_l\le\beta\le b_u<1,\quad |H|\le h.
\]
Define $Q(t)=-|t|^3+3\rho t^2/2+3\beta t$ and
$L_P(t)=C_EQ(t)+\phi(0)(t^3-3t)/6$.
For $t>0$, upper polynomials for $L_P(t),Q(t)$ are
\[
 (1/15-.4095)t^3+1.5(.40974)rt^2+(3(.40974)b_u-.3988/2)t,
 \quad -t^3+1.5rt^2+3b_ut.
\]
For $t<0$, write $u=|t|$ and use
\[
 (-.3988/6-.4095)u^3+1.5(.40974)ru^2+(.2-3(.4095)b_l)u,
 \quad -u^3+1.5ru^2-3b_lu.
\]
These coefficient choices are valid because $b_l\ge0$ is checked first.
On any closed cell of one sign, let $L_J$ be an upper bound for the
first polynomial and let $Q_J\ge0$ bound the positive part of the second.
For $j=1,2,3$, take $m_j$ to be the positive part of the largest tensor
Bernstein coefficient of $P_j$ on $[-h,h]\times J$. For $j=4,5$, the
absolute monomial bound
$\sum_{a,b}|[H^at^b]P_j|h^aR_J^b$ suffices, where
$R_J=\max_{t\in J}|t|$. The support contact is then bounded above by
\begin{equation}\label{support:paneltest}
 L_J+\frac{.7Q_J}{47}
 +.4\left(\sum_{j=1}^5\frac{m_j}{47^j}
                +\frac{4^7B_{R_J}}{47^6}\right)
 +r\left(.19S+\frac{.384(.425)}{47}\right),
 \qquad S>\sqrt{47/46}.
\end{equation}
Indeed substitute the finite series in contact, split $C=C_E+(C-C_E)$,
and use $C<.40974+.7/47<.425$. Every error retained in
\eqref{support:paneltest} is nonnegative and decreases with $n$.
A strictly negative rational result therefore excludes the entire cell
for every integer $n\ge47$, including atoms and accumulation points.

For a coordinate strip, conditional centering gives
\[
 X=\ell(Y-p),\quad
 \ell=\sqrt{(1-\theta)/(p(1-p))},\quad
 \beta=(1-\theta)(1-2p)+\kappa,\quad |\kappa|\le\theta.
\]
If $p\in[p_l,p_u]\subset(0,1/2)$ and $\theta\le\Theta$, put
$\ell_l\le\sqrt{(1-\Theta)/(p_u(1-p_u))}$.
The forbidden strip $Y\notin(-.55,1.55)$ maps into
\[
 X\le-\ell_l(p_l+.55),\qquad X\ge\ell_l(1.55-p_u).
\]
Apply \eqref{support:paneltest} between these thresholds and $\pm R_0$.
An interval beyond the known radius is empty; an equal endpoint is a
singleton that is still checked. If only an upper $\beta$ bound is given,
use $b_l=(1-\Theta)(1-2p_u)-\Theta$, after checking its nonnegativity.

The selected consequences used in the resource schedule are:
\[
\begin{array}{l|c|c|r}
\text{input rows}&\text{prior radius}&\text{conclusion}&\text{closed cells}\\\hline
10\text{ central preparation rows}&3.84&|X|<2.5&1080\\
45\text{ central rows}&2.5&-.55<Y<1.55&1431\\
38\text{ left rows}&2.7&-.55<Y<1.55&3073
\end{array}
\]
The respective largest contact bounds are below $-.084$, $-.097$ and
$-.01658$. These implications acquire no row by themselves: each is
used only after its input row is proved by the preceding resource step.
The left radii are applications of the all-threshold support polynomial,
with $(N,L,K,q)=(47,\infty,0,.19)$ and the actual moment caps
\[
 (r,b,R)=(1.390,.430,3),\qquad (1.2563,.3926,2.7).
\]
Thus the same heat proof supplies the initial radius, both left
contractions, and the five finite-base rows. The bounded-threshold
series supplies the central contraction and both coordinate strips.


% Source: resource_entry.tex
\section{Entry as a resource enclosure}\label{re:entry}

The entry calculation has one job: turn an attained extremizer into a
short list of simultaneous hypotheses for the local comparison. Its many
successive tables are applications of one enclosure principle. A table
is a cover of possible laws; a refinement intersects that cover with
newly proved inequalities. There is no assertion that an iteration
converges, and no improved bound is used to justify itself.

Throughout, $X$ has mean zero and variance one, $n\ge N$, and $H$ is
the raw sum threshold. Put
\[\begin{gathered}
 \rho=\E|X|^3,\quad \beta=\E[X|X|],\quad\gamma=\E X^3,\quad M_j=\E|X|^j,\\
 C=C_n,\quad e=C-C_E\ge0,\qquad C_E=\frac{3+\sqrt{10}}{6\sqrt{2\pi}}.
\end{gathered}\]
The input is an attained positive closed maximum, together with the
contact inequality stated below, an actual support bound, and an
independently established score cap. The global proof supplies these
inputs. A conditional upper bound on moments is never used to acquire
the support needed to define those moments.

\subsection{The conserved quantity}
Write $\tau(t)=\E[X\mathbf1_{X\ge t}]$. It is nonnegative,
$\int\tau=1$, and $2\int|t|\tau(t)\,dt=\rho$. On each bounded gap
$g=(a,b)$ of the closed support it has the constant value $\tau_g$.
Define
\begin{align*}
 J&=\sum_g\tau_g(b-a)^2,&
 J_1&=\sum_g\tau_g\int_a^b(t-a)t\,dt,\\
 J_2&=\sum_g\tau_g\int_a^b(t-a)t|t|\,dt,&
 D&=(3+\sigma)\rho-\gamma-3J,\qquad\sigma=\sqrt{10},\\
 L&=\sum_g\tau_g\int_{[a,b)}(t-a)\,\nu(dt),&
 \nu&=\mathcal L(H-S_{n-1}).
\end{align*}
The half-open convention includes the lower endpoint with weight zero
and excludes the upper endpoint. In particular it treats atoms without
an implicit midpoint convention. Bounded support makes all sums
absolutely convergent. The gap calculation below proves $D\ge0$;
clearly $L\ge0$.

Put $\varphi=\phi(H/\sqrt n)$ and
\begin{align*}
 K_1&=\E[X|X|^3]-3\beta-\tfrac32\rho\gamma,\\
 K_q&=M_5/4-\rho M_4/2-\tfrac32\beta\gamma+\rho/2,
 &U&=J_1+(M_4-3)/6.
\end{align*}
The exact remainder identity is
\begin{equation}\label{re:identity}
 ne\rho+nC_E\rho(1-e^{-w})+
 n\phi(0)e^{-w}D/6+n^{3/2}L=\mathcal R,
 \qquad w=H^2/(2n),
\end{equation}
where
\begin{equation}\label{re:raw}
 \mathcal R=C\{3\beta J/2+3\rho J_1-3J_2-K_q\}
 +H(CK_1+\varphi U)-H^2(C\rho+\gamma\varphi/2)+E_n.
\end{equation}
The complete Fourier estimate supplies a \emph{coherent} pair
$|E_n|\le K_F+s_F|H|$: the constant from one estimate is not combined
with the slope from another unless both were separately proved.
Section~\ref{re:foundation} derives \eqref{re:identity} and the whole
Fourier bound from the contact interface.

\begin{lemma}[Removing the Gaussian weight]\label{re:weight}
Suppose a first pass proves $\mathcal R<r<Nc_0$, where $C\rho\ge c_0$.
Every later valid upper bound $\mathcal R\le r'$ then gives
\begin{equation}\label{re:unweighted}
 ne\rho+n\phi(0)D/6+n^{3/2}L\le\mathcal R,
 \quad e^{-w}\ge1-\frac{5r'}{2n},
 \quad nD\le\frac{15000r'}{997}.
\end{equation}
The last two bounds use $C_E\rho>2/5$ and $\phi(0)>997/2500$.
\end{lemma}
\begin{proof}
Rearrange the exact identity before changing its Gaussian weight:
$\varphi(\gamma+3J)/6=C\rho+\sqrt nL-\mathcal R/n>0$.
Thus $C_E\rho-\phi(0)D/6>0$. Replacing $e^{-w}$ by one in
\eqref{re:identity} lowers its left side, proving the first assertion.
The original nonnegative phase term proves
$1-e^{-w}\le r'/(nC_E\rho)<5r'/(2n)$. The defect bound follows from
the first assertion. The initial sign license is proved before any
feedback is used.
\end{proof}

The same identity also gives a useful score update. If this cell forces
$D\ge d_l\ge0$ and a preceding pass gives raw ceiling $r$, then
\[
 C\le .40974+
 \frac{r/N-.3988(1-5r/(2N))d_l/6}{\rho_l}.
\]
A negative numerator excludes the cell because $e\ge0$. Otherwise its
numerator increases with $1/n$, and division by the positive lower
$\rho_l$ is outward. Thus the update is valid for every $n\ge N$ and
uses only the previous raw bound.

For later local comparisons we retain the sharper phase resource.
If previous bounds give $\mathcal R\le r$, $nD\le Nd$, then
\begin{equation}\label{re:phaseprice}
 ne\rho+\kappa nD+n^{3/2}L
 \le \mathcal B(H),\qquad \kappa=997/15000,
\end{equation}
where an upper bound for $\mathcal B$ is obtained from the raw
quadratic by increasing its curvature by
\begin{equation}\label{re:extra}
 q_{\rm extra}(\rho)=\frac{1-5r/(2N)}2
                         \left(.4095\rho-\frac d{15}\right)>0.
\end{equation}
Indeed the subtracted exact phase term is
$n(1-e^{-w})(C_E\rho-\phi(0)D/6)$.
Use $1-e^{-w}\ge we^{-w}$, $D\le d$, $\phi(0)<.4$, and the
previous density floor. The output field \texttt{Qbudget} includes
\eqref{re:extra}; it must not be charged a second time.

\subsection{One positive-measure rule}
Suppose zero lies in the gap $(-A,B)$, $A,B>0$. Let
\[\begin{gathered}
 p=\Pp(X>0),\quad q=1-p,\quad m=\E X_+=\E X_-,\quad V=\E[X\mid\mathbf1_{X>0}],\\
 W=X-V,\quad\theta_\pm=\E[W^2;\pm X>0],\quad\theta=\theta_++\theta_-.
\end{gathered}\]
Then $m^2=(1-\theta)pq$, and the conditional means have magnitudes
$\bar A=m/q$, $\bar B=m/p$. Set
$\Delta_A=\bar A-A\ge0$, $\Delta_B=\bar B-B\ge0$.
We also use the unnormalised endpoint displacements
\[
 M_+=\E[(X-B)\mathbf1_{X>0}]=p\Delta_B,\qquad
 M_-=\E[(-X-A)\mathbf1_{X<0}]=q\Delta_A.
\]

\begin{lemma}[Positive-measure pricing]\label{re:pricing}
For $x=|X|$, $a=|V|$, and a function $h$ constant on each sign class,
\begin{equation}\label{re:conditional}
 \E[h(|X|^j-|V|^j)]
 =\E\left[hW^2\sum_{k=0}^{j-2}(k+1)a^kx^{j-2-k}\right],\quad j\ge2.
\end{equation}
Also, for $d\mu_0=(|X|-1)^2(|X|+1/2)\,dP$,
\begin{align*}
 \mu_0(\R)&=\rho-1,\\
 M_4-1&=\int\frac{(x+1)^2}{x+1/2}\,d\mu_0,\\
 M_5-1&=\int\frac{x^3+2x^2+3x+3/2}{x+1/2}\,d\mu_0,\\
 K_1+3\rho\beta&=\int
  \frac{\operatorname{sgn}(X)x(x+2-3\rho/2)}{x+1/2}\,d\mu_0.
\end{align*}
For any nonnegative resources $z_i$ satisfying $\sum a_i z_i\le r$,
$a_i>0$, a signed linear cost obeys
\begin{equation}\label{re:simplex}
 \sum b_i z_i\le r\max\{0,b_1/a_1,\ldots,b_k/a_k\}.
\end{equation}
\end{lemma}
\begin{proof}
Divide $x^j-a^j-ja^{j-1}(x-a)$ by $(x-a)^2$; the linear term has
zero conditional mean. The radial identities follow by multiplying
out and using $\E X=0$, $\E X^2=1$. Finally each $b_i$ is at most
$a_i$ times the displayed maximum, and all $z_i$ are nonnegative.
\end{proof}

This one rule accounts for the moment, signed Taylor, endpoint movement,
outer-gap, and confinement-error estimates. In particular, no two costs
may each spend the full residual and then be added.

\subsection{The simultaneous gap budget}
Let $p_*=(4-\sigma)/2$ and
$T_+=\E X_+^3-mB^2$, $T_-=\E X_-^3-mA^2$.
If $J_o=J-m(A+B)^2$, signwise integration gives
\begin{equation}\label{re:gapexact}
 D=m(\sqrt{\sigma+1}A-\sqrt{\sigma-1}B)^2
   +(\sigma-1)T_++(\sigma+1)T_-+3(T_++T_--J_o).
\end{equation}
Moreover, writing $R_+=\E[X_+(X_+-B)]$ and similarly $R_-$,
\[
 T_++T_--J_o\ge2BR_++2AR_-\ge0.
\]
If zero belongs to the support, every gap has one sign and $J\le\rho$,
so $D\ge\sigma\rho-\gamma\ge(\sigma-1)\rho$.
Otherwise \eqref{re:gapexact} proves $D\ge0$.
Consequently $D<(\sigma-1)\rho$ licenses the actual central gap.
We also have $J\le\rho+2m\le\rho+1$; if $D<2$, then
$T_++T_-\le D/2$ and $A+B\ge\sqrt{2-D}$.

To verify the gap estimate, on a positive outer gap $(u,v)$,
$(v-u)^2\le2\int_u^v(t-B)\,dt$.
Summing with weights $\tau_g$ and applying Tonelli gives
$J_o^+\le T_+-2BR_+$; reflect for the other sign.
Expanding the definition of $D$ gives \eqref{re:gapexact}.
The inequality $AB\le1$ follows by integrating $(X+A)(X-B)\ge0$,
and $m\le1/2$ follows from $\E|X|\le1$.

Conditional expansion about $\bar A,\bar B$ now gives the stronger
single inequality
\begin{equation}\label{re:joint}
\begin{split}
 D\ge{}&\frac{(6+2\sigma)(1-\theta)^{3/2}}{\sqrt{pq}}(p-p_*)^2
 +6m\{(A+B)(\Delta_A+\Delta_B)+\Delta_A\Delta_B\}\\
 &+3(\sigma+1)B\theta_++3(\sigma+3)A\theta_-.
\end{split}
\end{equation}
For example
$\E X_+^3=m\bar B^2+\E[(X+2\bar B)(X-\bar B)^2;X>0]$
and $R_+=\theta_++m(\bar B-B)$.
The cubic residual is at least $3B\theta_+$.
The deterministic shift left after substituting both signs is
$6m\{\bar A\bar B-AB+B\Delta_B+A\Delta_A\}$, equal to the
second term above. The remaining deterministic expression, using
$m^2=(1-\theta)pq$, is exactly the displayed binary square.

\subsection{A state, its update, and its exported resource}
A state for a closed probability interval $I=[l,u]$ records
\[
 |X|\le R,\quad C\in[c_0,c_1],\quad\mathcal R\le r,
 \quad nD\le Nd,\quad\theta\le\Theta<1,
 \quad A\ge A_l>0,\quad B\ge B_l>0,
\]
together with simultaneous signed-moment and coefficient intervals.
The distinction between $nD\le Nd$ and $D\le d$ is essential.
Let $s_+=\max_Ip(1-p)$, $s_-=\min_Ip(1-p)$ and
\begin{gather*}
 b_I=\frac{(6+2\sigma)(1-\Theta)^{3/2}}{\sqrt{s_+}}
                 \operatorname{dist}(p_*,I)^2,
 \qquad r_I=d-b_I,\\
 K_+=3(\sigma+1)B_l,\quad K_-=3(\sigma+3)A_l,
 \quad K=\min(K_+,K_-),\quad m_l=\sqrt{(1-\Theta)s_-}.
\end{gather*}
If $r_I<0$, exclude the entire closed cell. Otherwise \eqref{re:joint}
gives, simultaneously,
\begin{equation}\label{re:residual}
 6m(A_l+B_l)(\Delta_A+\Delta_B)+K_+\theta_++K_-\theta_-\le r_I.
\end{equation}
Thus the actual law satisfies the intersected bounds
\begin{gather}
 \theta\le\min\{\Theta,r_I/K\},\notag\\
 A\ge\max\left\{A_l,\sqrt{(1-\Theta)l/(1-l)}-
                  \frac{r_I}{6m_l(A_l+B_l)}\right\},\notag\\
 B\ge\max\left\{B_l,\sqrt{(1-\Theta)(1-u)/u}-
                  \frac{r_I}{6m_l(A_l+B_l)}\right\}.\label{re:geometry}
\end{gather}
Each right side is evaluated from established bounds; sequential
improvements may then be used by the next expression. A nonpositive
endpoint estimate fails this update, rather than excluding a law.
Initialization uses the same formulas with the valid lower sum
$A+B\ge\sqrt{2-d}$.

In Bernoulli coordinates set $Y=B_0+U$, $\E[U\mid B_0]=0$,
$B_0\sim\mathrm{Bernoulli}(p)$, $X=\ell(Y-p)$ and
$d_Y=(\Delta_A+\Delta_B)/(\bar A+\bar B)$.
Then $\ell^2=(1-\theta)/(pq)$ and $\lambda=n\E U^2=npq\theta/(1-\theta)$.
The \emph{weighted} versions of \eqref{re:residual} prove
\begin{equation}\label{re:exports}
 n\theta\le\frac{Nr_I}{K},\quad
 nd_Y\le\frac{Nr_I}{6(1-\Theta)(A_l+B_l)},\quad
 \lambda\le\frac{Ns_+r_I}{(1-\Theta)K}.
\end{equation}
Indeed $nD\le Nd$ whereas $nb_I\ge Nb_I$.
These exports are derived from the weighted resource; multiplying an
unweighted variance ceiling by $N$ would not prove them.

For $j=3,4,5$ set
\[
 L_3(v)=R+2v,\quad L_4(v)=R^2+2Rv+3v^2,
 \quad L_5(v)=R^3+2R^2v+3Rv^2+4v^3.
\]
Put $v_+=\sqrt{(1-l)/l}$, $v_-=\sqrt{u/(1-u)}$ and
$k_j=[1-(1-\Theta)^{j/2}]/\Theta$, with $k_j=j/2$ at $\Theta=0$.
For $u<1/2$, valid Bernoulli moment ceilings are, with $s=l(1-l)$,
\[
 B_3^*=(1-2s)/\sqrt s,\quad B_4^*=1/s-3,
 \quad B_5^*=(1-4s+2s^2)/s^{3/2}.
\]
They decrease towards $p=1/2$. Convexity places $(1-\theta)^{j/2}$
below its endpoint chord on $[0,\Theta]$; positive-measure pricing gives
\begin{equation}\label{re:moments}
 M_j\le B_j^*+r_I\max\left\{0,
 \frac{L_j(v_+)-k_jB_j^*}{K_+},
 \frac{L_j(v_-)-k_jB_j^*}{K_-}\right\}.
\end{equation}
The same rule, using
$T_+\le2mv_+\Delta_B+L_3(v_+)\theta_+$ and its reflection, bounds
\begin{equation}\label{re:outerresource}
 W_o:=(\sigma-1)T_++(\sigma+1)T_-
 \le r_I\max\left\{\begin{gathered}
 \frac{(\sigma-1)v_+}{3(A_l+B_l)},
 \frac{(\sigma+1)v_-}{3(A_l+B_l)},\\
 \frac{(\sigma-1)L_3(v_+)}{K_+},
 \frac{(\sigma+1)L_3(v_-)}{K_-}\end{gathered}\right\}.
\end{equation}
Intersect every new bound with the previous one. The versions which
merge $\theta_+$ and $\theta_-$ are obtained by replacing both costs by
their maximum and both prices by their minimum; they need no new proof.

\subsection{The enclosure theorem}\label{re:operator}
On each cell, the formula atlas bounds the four coefficients in
\eqref{re:raw}. The result, for each actual sign $s=\operatorname{sgn}H$,
is a finite family of quadratics
\begin{equation}\label{re:quadratics}
 \mathcal R\le\max_j(A_j+F_{j,s}|H|-Q_j^{\rm raw}H^2),
 \qquad Q_j^{\rm raw}>0.
\end{equation}
The maxima retain all score--moment corners. In the initial signed
cover, intervals contain both signs of $\beta,\gamma$; no orientation is
discarded. In the probability cover, the central integral bound is a
maximum of affine functions of $\rho$, hence convex. The other
$\rho$-dependence is affine or the maximum of convex quadratics. For
fixed $H$, the score dependence is convex. Endpoints are therefore taken
\emph{before} completing the square.

\begin{theorem}[Finite resource enclosure]\label{re:enclosure}
Consider a finite ordered schedule of the following operations:
\begin{enumerate}
\item Split a cell into an explicitly covered family of closed cells.
\item Intersect its enclosures with an inequality proved in the atlas.
\item Exclude a cell by an empty necessary interval or a strictly
negative nonnegative resource.
\item Apply a separately proved support or confinement implication,
checking all its hypotheses against this cell's already acquired data.
\item Export a raw, weighted, or signed phase bound by
\eqref{re:quadratics}, \eqref{re:phaseprice}, and \eqref{re:exports}.
\end{enumerate}
If exact outward arithmetic checks every domain, corner and cover,
every output contains all attaining laws satisfying the input. The
conclusion holds for all integers $n\ge N$ wherever the formula's
monotonicity gates are checked. It includes zero conditional variance,
singleton intersections and arbitrarily small atoms.
\end{theorem}
\begin{proof}
Induct over the schedule. A split preserves the union, intersections use
necessary inequalities, and an excluded cell contains no feasible law.
Each callback consumes only earlier hypotheses for that same law.
For a scalar raw output, use
$fh-qh^2\le(f_+)^2/(4q)$ for $h\ge0$, $q>0$; using $f^2$ is a valid
weaker bound. For a resource output first add \eqref{re:extra} to each
curvature. A negative maximum excludes the cell; a nonnegative one may
be divided by $\kappa N$ to give its next defect parameter.
The phase bound follows from the positive root of
$A+fh-qh^2\ge0$, replacing $A$ by $A_+$ if necessary.
A strict output is exported only when the certificate retains a strict
upper gap, such as its strict Fourier or target enclosure; closed-domain
coverage does not change strictness. All constants are frozen before evaluating
corners. Thus no iteration
step relies on an unproved fixed point. The stated index uniformity is
that of the displayed finite formulas, not a fixed-law limit.
\end{proof}

For the Fourier cancellation estimate only, reflect $(X,H)$ if needed:
$I_0(-X,-H)=I_0(X,H)$. Both odd coefficients change sign, and
$\int e^{-x^2/2}|zx^2-cx^4|\,dx$ is unchanged by $(z,c)\mapsto(-z,-c)$.
Consequently a lower endpoint used there bounds $|\gamma|$, not necessarily
the original signed $\gamma$. The actual contact, gap and phase data keep
their original orientation. A useful additional signed lower bound is
\begin{equation}\label{re:gamma}
 \gamma\ge(1-\Theta)^{3/2}\frac{1-2u}{\sqrt{u(1-u)}}
       -\left(R+2\sqrt{\frac{u}{1-u}}\right)\Theta.
\end{equation}
To prove it expand $\E X^3=\E V^3+\E[W^2(X+2V)]$; the positive
residual is nonnegative and the negative residual has the displayed
lower bound. The binary term decreases in both $p<1/2$ and $\theta$.
Thus $\max(0,\text{any signed lower bound})$ is a valid lower bound for
$|\gamma|$ inside this auxiliary Fourier estimate.

\subsection{Support, confinement, and a short acyclic schedule}
Two different support implications are useful. Gaussian contact excludes
a point by a strict negative upper bound on its contact expression;
closed cells cover the finite exterior and a derivative-sign certificate
covers any remaining ray. The shared Gaussian support theorem supplies
this callback, including the bounded-threshold radius and coordinate-strip
contracts; it is a prior theorem, not a test on moments alone. Noise confinement proves $|U|\le a$ on an
already acquired row. It then implies, without a small-noise limit,
\begin{equation}\label{re:radiusfeedback}
 |X|\le R_*:=
 \frac{\max\{u+a,1-l+a\}}{\sqrt{\min_Ip(1-p)}}.
\end{equation}
The denominator uses $\E U^2\ge0$, so zero noise is included.
The Fourier/moment formulas remain valid for $R\ge1$; the paired Pearson
choice additionally checks $R>2g$. In particular there is no intrinsic
lower radius $2.3$: the radial ratios in Lemma~\ref{re:pricing} increase
on all $x\ge0$, and the other formulas use their explicit domains.

Here is the selected proof schedule. Each row names \emph{an application}
of Theorem~\ref{re:enclosure}, not a different analytic theorem.
\begin{center}\small
\begin{tabular}{p{.13\linewidth}p{.32\linewidth}p{.45\linewidth}}
\toprule
Node&Input and operation&Output consumed later\\\midrule
I&$N=47$, $R=3.84$, $C\le.4247$; 13 signed passes plus export&
100 closed probability cells; 35 survive; $\mathcal R<1.885$, $nD<47(.590)$\\
C&Central preparation, contact radius $2.5$, probability refinement&
45 cells $[.38,.47]$, width $.002$; canonical simultaneous input\\
L&Complete complementary cover, contact radius $3$, refinement&
39 cells $[.302,.38]$; all other complementary cells excluded\\
C1&First confinement from C; radius formula; enclosure&
First improved central table\\
C2&Antichain confinement from C; intersect C1; enclosure&
Selected central table: 45 cells, $[.38,.47]$\\
S&Separate sign-resource pricing on C2 and L&
One left cell excluded; 38 left cells remain\\
L2&Left contact radius $2.7$ from S; enclosure&
Selected left table: 38 cells, $[.304,.38]$\\\bottomrule
\end{tabular}
\end{center}
Every shared endpoint belongs to its adjacent closed cells. The
confinement producing C2 consumes C, not C2. The selected central local
comparison uses C2; it does not need the optional further central
refinement in S. The final cover is exactly $[.304,.47]$.

\begin{corollary}[Uniform acquired entry]\label{re:acquired}
Let an attained positive maximum have $n\ge47$, $C_E\le C\le.4247$,
$0\le n(C-C_E)\le.7$, decrement at most $.19$, and both deleted
score caps at most $.4253$.
Given the shared contact support and strip implications, the exact
schedule above places its law in one of the 38 selected left rows or
45 selected central rows, with every resource in that row holding
simultaneously. This assertion has no upper bound on $n$.
\end{corollary}
\begin{proof}
The shared all-threshold support theorem supplies radius $3.84$.
Apply the finite enclosure theorem in the displayed order, retaining
the weighted bounds at every edge. The initialization's stored upper
index 98 occurs only through $.4095+K/98$; here $K=0$, so it imposes
no upper-index restriction. All subsequent index-dependent inequalities
are checked for every $n\ge47$ by their stated monotonicity gates.
The accepted closed covers exclude the complementary cells and give
the stated 83 rows.
\end{proof}

\paragraph{The local-comparison interface.}
For each final row supply $p\in I$, $\theta\le\Theta$, $\lambda\le\Lambda$,
the actual support/amplitude, $d_Y\le d_0$, both signed phase limits,
and all four coherent quadratics bounding \eqref{re:phaseprice}.
On a phase/noise subcell, let $T_H$ be the maximum of each concave
quadratic at its clipped vertex, then the maximum over corners and the
minimum with the scalar bound. If a proved moment inequality is
\[
 D\ge\alpha d_Y+b\theta+P,\qquad\alpha>0,
 \quad \ell\ge\ell_-,\quad\lambda\ge L_0,
\]
then the exact strict resource supplied to the comparison is
\begin{equation}\label{re:localG}
 nd_Y<G:=\frac{T_H/\kappa-NP-b\ell_-^2L_0}{\alpha}.
\end{equation}
This uses $n\theta=\ell^2\lambda$ and $nP\ge NP$. If $G\le0$,
the entire closed subcell is excluded. Formula~\eqref{re:localG} is
the same resource as \eqref{re:phaseprice}, not an additional allowance.

\paragraph{Validation contract.}
The exact supplement must pin the input and output hashes of every
node, reconstruct all closed covers, and recompute its inequalities.
Every square-root enclosure is checked by squaring nonnegative endpoints;
every division has a positive-denominator gate. Exponentials and the
finite Gaussian integral use the stated rational series enclosures.
The first raw pass checks $r<.4N$ before exporting a defect; later passes
consume their predecessors' raw and weighted bounds. Excluded cells
retain an explicit mathematical reason and count towards coverage;
an empty final row is not a negative remainder bound. Reordered,
duplicated, missing, malformed and falsely strengthened objects reject.
The accompanying source map identifies the exact implementation of each
formula and each independently accepted callback. This consolidates the
proof; it does not replace the arithmetic replays by a hash check.




% Source: entry_formula_atlas.tex
\section{The minimal contact--Fourier foundation}\label{re:foundation}
The only contact input needed here is the following explicit interface.
At the attained positive maximum, with $F=\Pp(S_n\le H)$,
$z=H/\sqrt n$ and $\nu=\mathcal L(H-S_{n-1})$, one has
$\nu([y,\infty))\le\mathfrak q(y)$ for every real $y$, with equality
on the full closed support, where
\begin{align}
 n\mathfrak q(y)={}&nF+\tfrac12\{z\phi(z)+C\rho/\sqrt n\}
 -(\sqrt n\phi(z)+3C\beta/\sqrt n)y\notag\\
 &-\tfrac12\{z\phi(z)+3C\rho/\sqrt n\}y^2
 +C|y|^3/\sqrt n.\label{re:contact-q}
\end{align}

This contact statement is the normalized main-manuscript contact inequality at $c=C_n$. Attainment is supplied by its ratio-attainment corollary.
Everything that follows in this section is derived from it. We abbreviate
$M_4=\mu_4$, $M_5=\eta_5$ when comparing with the implementation.
\subsection{Compensated smooth inversion}
Write $Z_n=S_n/\sqrt n$, $f(t)=\E e^{itX}$, and let $N$ be standard normal.
The smooth payoff $q_z(x)=(x-z)|x-z|$ gives the exact error
\[
 \epsilon_n(z)=\E q_z(Z_n)-\E q_z(N)-\frac{2\gamma\phi(z)}{3\sqrt n}.
\]
\begin{lemma}\label{smooth:fourier}\label{smooth:inversion}
For every real $H$ and $n\ge1$,
\begin{equation}\label{smooth:smooth-integral}
 \epsilon_n(H/\sqrt n)=-\frac4{\pi n}\int_0^\infty
 \frac{\Im(e^{-itH}V_n(t))}{t^3}\,dt,\qquad
 V_n(t)=f(t)^n-e^{-nt^2/2}(1-in\gamma t^3/6).
\end{equation}
The integral is absolutely convergent. With the split at $t=1/2$,
\begin{equation}\label{smooth:split}
 I_{\rm low}=\int_0^{1/2}\frac{|\Im(e^{-itH}V_n(t))|}{t^3}\,dt,
 \qquad I_{\rm high}=\int_{1/2}^\infty\frac{|V_n(t)|}{t^3}\,dt,
\end{equation}
one has $|\epsilon_n(z)|\le4(I_{\rm low}+I_{\rm high})/(\pi n)$.
\end{lemma}
\begin{proof}
The scalar identities
\[
 \int_0^\infty\frac{ty-\sin(ty)}{t^3}\,dt=\frac\pi4y|y|,
 \qquad
 \int_0^\infty\frac{|ty-\sin(ty)|}{t^3}\,dt=\frac\pi4y^2
\]
follow by scaling and integration by parts. One elementary evaluation of
the constant uses
$J_a(b)=\int_0^\infty e^{-at}(1-\cos bt)t^{-2}dt$:
$J_a''(b)=a/(a^2+b^2)$, $J_a(0)=J_a'(0)=0$, so
$J_a(1)=\int_0^1\arctan(b/a)db\to\pi/2$.
Monotone convergence removes $a>0$.
For a finite signed measure $\mu$ with zero mass and first moment and
finite absolute second moment, absolute Fubini applied to the
\emph{compensated} scalar integrand therefore gives
\begin{equation}\label{smooth:signed-inversion}
 \int (s-H)|s-H|\,d\mu(s)
 =-\frac4\pi\int_0^\infty t^{-3}
       \Im(e^{-itH}\widehat\mu(t))\,dt.
\end{equation}
Apply this to the law of $S_n$ minus the pushforward, by $x\mapsto\sqrt n x$,
of the signed density
$\phi(x)[1+\gamma(x^3-3x)/(6\sqrt n)]$.
Gaussian integration by parts gives its transform
$e^{-nt^2/2}(1-in\gamma t^3/6)$, total mass one and first moment zero.
Since $(x^3-3x)\phi=-\phi'''$, integration twice by parts and then across
$z$ gives $-\int q_z\phi'''=4\phi(z)$. This proves the formula and its
scale factor. The high integral is finite from $|f|\le1$ and Gaussian
decay; the low integral is finite by the absolute Fubini argument.
\end{proof}


The payoff is smooth enough for compensated inversion. No discontinuous
CDF is inverted, so no atom midpoint convention is introduced.
\subsection{Derivation of the conserved identity}
\begin{proof}
Compute $\mathcal A=\int\tau(y)[\mathfrak q(y)-\nu([y,\infty))]dy$
in two ways. Put $r=-\mathfrak q'$. Contact at both endpoints of a gap gives
\[
 \nu([a,b))=\int_a^b r(t)dt,
 \quad
 \int_a^b[\mathfrak q(y)-\nu([y,\infty))]dy
 =\int_a^b(t-a)r(t)dt-\int_{[a,b)}(t-a)d\nu(t).
\]
The integrand vanishes on the support and $\tau=0$ on the exterior rays.
Summing, and differentiating \eqref{re:contact-q}, gives
\begin{equation}\label{smooth:A-gap}
 \mathcal A=\frac{\phi(z)J}{2\sqrt n}+\frac{G}{n^{3/2}}-L,
 \qquad
 G=\frac32C\beta J+(H\phi(z)+3C\rho)J_1-3CJ_2.
\end{equation}
All sums and integrals are absolutely convergent on the compact convex
hull of the support, since $|J_1|\le R_0J/2$, $|J_2|\le R_0^2J/2$
when $|X|\le R_0$.

For the second computation let $X^*$ have density $\tau$, independent of
$S_{n-1}$. Signwise Fubini gives
$\E g'(X^*)=\E[Xg(X)]$ for continuously differentiable $g$.
In particular,
$\E X^*=\gamma/2$, $\E(X^*)^2=\E X^4/3$ and
$\E|X^*|^3=\E|X|^5/4$.
Applying the identity to each summand shows that
$(S_{n-1}+X^*)/\sqrt n$ has density
$\tau_{Z_n}(t)=\E[Z_n\mathbf1_{Z_n\ge t}]$.
Its CDF at $z$ is
$T_z=\E[(Z_n^2-zZ_n)\mathbf1_{Z_n\le z}]$, by integrating that density;
the boundary integrand is zero. Substituting the three moments into
\eqref{re:contact-q} therefore gives
\begin{equation}\label{smooth:zero-bias-average}
 \mathcal A=F-\frac{\gamma\phi(z)}{2\sqrt n}
 +\frac{H\phi(z)(3-\E X^4)/6+CK_q}{n^{3/2}}-T_z.
\end{equation}
Multiply contact equality by $X$ and integrate. By independence and symmetry,
$n\E[X\nu([X,\infty))]$ equals $\E[S_n\mathbf1_{S_n\le H}]$.
Expanding the other side yields, for $M_z=\E[Z_n\mathbf1_{Z_n\le z}]$,
\begin{equation}\label{smooth:stationarity}
 M_z+\phi(z)=-\frac{\gamma z\phi(z)}{2\sqrt n}+\frac{CK_1}{n}.
\end{equation}
Finally $q_z(x)=(x-z)^2-2(x-z)^2\mathbf1_{x\le z}$, so Gaussian
integration gives
\begin{equation}\label{smooth:truncated-square}
 T_z-\Phi(z)=-\tfrac12\{\E q_z(Z_n)-\E q_z(N)\}
 +z(M_z+\phi(z))-z^2(F-\Phi(z)).
\end{equation}
Insert \eqref{smooth:stationarity} and the definition of $\epsilon_n$
into \eqref{smooth:zero-bias-average}. With
\[
 P_0=H^2(C\rho+\gamma\phi(z)/2)
       -H\{CK_1+\phi(z)(\E X^4-3)/6\}+CK_q,
\]
the result is $\sqrt n\mathcal A=C\rho-\gamma\phi(z)/6+
P_0/n+\sqrt n\epsilon_n(z)/2$. Equating this with
\eqref{smooth:A-gap} proves
\begin{equation}\label{smooth:signed-identity}
 C\rho+\sqrt nL=\frac{\phi(z)}6(\gamma+3J)
                 +\frac{G-P_0}{n}-\frac{\sqrt n}2\epsilon_n(z).
\end{equation}
Now substitute $\gamma+3J=(3+\sqrt{10})\rho-D$ and
$\phi(z)=\phio e^{-w}$, and expand $G-P_0$.
\end{proof}

\subsection{One complete low-frequency transfer}
Put $f(t)=\E e^{itX}$, $h(t)=e^{-t^2/2-i\gamma t^3/6}$,
$e_n(t)=e^{-nt^2/2}(1-in\gamma t^3/6)$, and
\[
 I_0(H)=\int_0^{1/2}|\Im(e^{-itH}(f^n-e_n))|\frac{dt}{t^3},\qquad
 I_\infty=\int_{1/2}^\infty|f^n-e_n|\frac{dt}{t^3}.
\]
The inversion bound~\eqref{smooth:split} charges
$2\sqrt n(I_0+I_\infty)/\pi$; both integrals below are complete.

\begin{lemma}[Signed Gaussian transfer]\label{fourier:early-low}
\label{fourier:paired-low}\label{entryleap:fourier-transfer}
Suppose $n\ge N>2$, $\mu_4\le M\le11.1$, $|\gamma|\le g\le1.85$,
$\mu_6\le M_6$, and the actual coefficients
\[
 a=(3-\mu_4)/24,\quad z=(10\gamma-\E X^5)/120,\quad c=a\gamma/6
\]
satisfy $|a|\le a_*$, $|z|\le Z$. Put
\begin{equation}\label{fourier:taylor}
 E_6=M_6/720+g^2/72+(1/2+g^2/144)^3/5,\quad
 B_6=a_*+Z/2+E_6/4.
\end{equation}
Let $a_0=(N-1)/(2N)$,
$0<b_0\le(45-M)(1-2/N)/96$, and
$\lambda_N=\sqrt{\pi/2}(N/(N-1))^{3/2}$.
Then, for every real $H$,
\begin{equation}\label{fourier:low-form}
 I_0(H)\le(L_0+L_H|H|)/\sqrt n+L_2/n,
 \qquad L_2=E_6/(2a_0^2)+B_6^2/(2b_0^3).
\end{equation}
Two admissible choices, each used as a complete triple, are as follows.
The triangle choice is
\begin{equation}\label{fourier:early}
 L_0=c_0Z+c_1a_*g+251g^3/17280,\quad
 L_H=c_0a_*+251g^2/4800,
\end{equation}
where $c_0\ge\lambda_N$, $c_1\ge\lambda_N/2$.
For the cancellation choice, take any rectangle containing the actual
$(z,c)$, $I_u\ge\sqrt{\pi/2}$ and $F_N\ge(N/(N-1))^{3/2}$. Define
\begin{gather*}
 U(z,c)=\begin{cases}
 I_u(|z|+3|c|),&zc\le0,\\
 I_u\bigl||z|-3|c|\bigr|+6\sqrt{|zc|}e^{-|z|/(2|c|)},&zc>0,
 \end{cases}\\
 J_u=\max_{\rm corners}U(z,c).
\end{gather*}
Then
\begin{align}\label{fourier:paired}
 L_0&=F_NJ_u+\frac{F_NI_u}{N-1}
 (35a_*g^3/48+35Zg^2/24)+251g^3/17280,\notag\\
 L_H&=F_NI_u\{a_*+5Zg/[2(N-1)]\}+251g^2/4800.
\end{align}
Neither the transfer nor its cancellation choice assumes signs for
$a,z,c$ or $\gamma$.
\end{lemma}
\begin{proof}
Taylor with integral remainder gives
$f-h=-at^4-izt^5+r_6$, $|r_6|\le E_6t^6$ for $0\le t\le1/2$.
Indeed for $w=-t^2/2-i\gamma t^3/6$,
$|w|\le t^2(1/2+g^2/144)<1/6$; the terms through degree five in
$1+w+w^2/2$ are exact, its degree-six term costs $g^2t^6/72$, and
the remaining exponential series costs $|w|^3e^{|w|}/6\le|w|^3/5$.
Add the true sixth-moment remainder.

Four integrations of $1-\cos u\ge0$ give
$\cos u\le1-u^2/2+u^4/24$. Symmetrization, using
$\E(X-X')^4=2\mu_4+6$, therefore gives
$|f(t)|,|h(t)|\le e^{-(45-M)t^2/96}$ on $[0,1/2]$.
Every point on their joining segment has the same bound, so polynomial
Taylor expansion yields
\[
 f^n-h^n=n(f-h)h^{n-1}+R_2,\qquad
 |R_2|\le\tfrac12n^2B_6^2t^8e^{-b_0nt^2}.
\]
The $r_6$ and $R_2$ integrals are at most
\begin{equation}\label{fourier:common-error}
 \frac1n\{E_6/(2a_0^2)+B_6^2/(2b_0^3)\}=L_2/n.
\end{equation}
For the remaining comparator put $u=n\gamma t^3/6$. Its imaginary part
is bounded by $e^{-nt^2/2}(|u|^3/6+|H|tu^2/2)$, so its divided integral
is at most
\begin{equation}\label{fourier:comparator-low}
 n^{-1/2}(251g^3/17280+251g^2|H|/4800).
\end{equation}
Here the Gaussian moments follow from integration by parts,
$\int_0^\infty x^{2j}e^{-x^2/2}dx=(2j-1)!!\sqrt{\pi/2}$,
and $\sqrt{2\pi}<251/100$.

It remains to price the leading quartic and fifth terms. Set $b=n-1$,
$v=b\gamma t^3/6$. The triangle estimate uses
$|\sin(Ht+v)|\le|H|t+bg t^3/6$, $|\cos|\le1$; the Gaussian
moments of orders two and four give respectively $\lambda_N$ and
$\lambda_N/2$. This proves~\eqref{fourier:early}.
For cancellation, after division by $t^3$, rotation by $Ht$ bounds the
leading imaginary part by
\[
 |zt^2-atv|+a_*t|\sin v-v|+Zt^2|\cos v-1|
       +|H|t\{a_*t+Zt^2|\sin v|\}.
\]
Multiply by $ne^{-bt^2/2}$ and use the global sine/cosine remainders.
The first integral equals $nJ(z,c)/b^{3/2}$, where
$J(z,c)=\int_0^\infty e^{-x^2/2}|zx^2-cx^4|dx$.
This norm is convex and $J(-z,-c)=J(z,c)$, so its rectangle maximum
is at a corner. For $zc\le0$ the integral is
$\sqrt{\pi/2}(|z|+3|c|)$. For $z,c>0$, split at $r=\sqrt{z/c}$;
integration by parts gives
\[
 J=\sqrt{\pi/2}(z-3c)+6cr e^{-r^2/2}
       +2(3c-z)\int_r^\infty e^{-x^2/2}dx.
\]
Discard the last term when nonpositive, and otherwise bound its integral
by $\sqrt{\pi/2}$. This proves $J\le U$; simultaneous sign reflection
covers the remaining quadrant. The other four phase terms use moments
of orders two, six, eight and ten and give the coefficients
$1,5/2,35/24,35/48$ in~\eqref{fourier:paired}. Finally $n/(n-1)$,
$1/(n-1)$, $a_0^{-1}$ and the allowable damping inverses decrease with
$n$, proving the all-index bounds at $N$.
\end{proof}

\begin{lemma}[Retaining the signed Gaussian tail]\label{entry189:mills}
For $r>0$, put
\[
 \mathcal M(r)=e^{r^2/2}\int_r^\infty e^{-t^2/2}\,dt,
 \qquad
 \mathcal M_k(r)=\cfrac1{r+\cfrac1{r+\cfrac2{\ddots+\cfrac{k}{r}}}}
 \quad(k\ge1),\qquad \mathcal M_0(r)=1/r.
\]
Odd truncations are lower bounds for $\mathcal M$; even truncations
are upper bounds. For $z,c>0$, set $x=z/c$ and
\[
 d_{k+1}=1,\qquad d_j=1+\frac j{xd_{j+1}}\quad(j=k,\ldots,1),
 \qquad m_k(x)=d_1^{-1}.
\]
Then the cancellation norm in Lemma~\ref{entryleap:fourier-transfer}
satisfies
\begin{equation}\label{entry189:mills-bound}
 J(z,c)\le \sqrt{\pi/2}(z-3c)+\sqrt{zc}\,e^{-x/2}
 \left(6+\frac{2(3-x)}x m_k(x)\right),
\end{equation}
where $k=11$ if $x\ge3$ and $k=12$ if $x<3$.
The factor in parentheses is positive. One may take the minimum with
the previous bound $U(z,c)$ at each corner of the coefficient rectangle.
\end{lemma}
\begin{proof}
The positive integrals
$I_j(r)=\int_0^\infty t^j e^{-rt-t^2/2}\,dt$ satisfy
$I_1=1-rI_0$ and $I_{j+1}=jI_{j-1}-rI_j$ by integration by parts.
Thus, writing $q_j=I_j/I_{j-1}>0$,
\[
 I_0=\frac1{r+q_1},\qquad q_j=\frac j{r+q_{j+1}}.
\]
Replacing the final positive $q_{k+1}$ by zero reverses the inequality
at each substitution. This proves the alternating finite bounds, since
translation gives $I_0=\mathcal M(r)$. No convergence of an infinite
continued fraction is required. The displayed rational recursion gives
$m_k(x)=r\mathcal M_k(r)$ when $r=\sqrt x$.

Splitting the norm at $r=\sqrt{z/c}$, as in the preceding proof, gives
the exact identity
\[
 J(z,c)=\sqrt{\pi/2}(z-3c)+\sqrt{zc}\,e^{-x/2}
 \left(6+\frac{2(3-x)}x r\mathcal M(r)\right).
\]
The multiplier of $\mathcal M$ is nonpositive when $x\ge3$ and
positive when $x<3$. Substitute the corresponding finite lower or
upper bound. Positivity of the remaining factor is immediate for
$x<3$; for $x\ge3$, use $0<m_{11}(x)<1$ to bound it below by
$6-2(x-3)/x>4$. Finally the exact norm $J$, rather than either
upper-bound formula, is convex, so its containing rectangle is still
exhausted by its four corners.
\end{proof}

At every positive coefficient corner use the displayed finite Mills bound.
A negative Gaussian coefficient uses a lower enclosure for
$\sqrt{\pi/2}$; a nonnegative coefficient uses an upper one. The remaining
factor is positive, so outward upper square-root and exponential enclosures
are valid. A simultaneous reflection of both odd coefficients leaves the
norm unchanged.

\paragraph{The two coefficient choices in the certificates.}
All coefficients come from the actual marginal. For $Q=X^2-\gamma X-1$,
multiplication gives
\begin{align}\label{fourier:pearson}
 \E Q^2&=\mu_4-1-\gamma^2,&
 \E X^5&=\gamma^3+2\gamma+\E[Q^2(X+2\gamma)],\notag\\
 \mu_6&=1+3\gamma^2+\gamma^4+
       \E[Q^2\{(X+\gamma)^2+2\gamma^2+2\}].
\end{align}
For any radius $R_0\ge|X|$ with $R_0\ge2g$, the two signs of
$\E X^5-10\gamma$ are bounded, at $t=|\gamma|$, by
\begin{align*}
 P_{R_0}(t)&=R_0(M-1)+(10-2M)t-R_0t^2+t^3,\\
 Q_{R_0}(t)&=R_0(M-1)+(2M-10)t-R_0t^2-t^3.
\end{align*}
The gate $R_0-2g\ge0$ permits substituting $M$ for $\mu_4$ in both
signs. The triangle choice uses the valid enclosing radius $R_0=6$, $a_* =\max(2,M-3)/24$,
$M_6=RE_5$, $b_0=(45-M)(1-2/N)/96$ and
\begin{equation}\label{fourier:z0}
\begin{gathered}
 120Z=\min\{E_5+10g,Z_0\},\\
 Z_0=\begin{cases}
 P_6(g),&M\le5, P'_6(g)\ge0,\\
 6(M-1)+(10-2M)^2/[4(6-g)],&M\le5, P'_6(g)<0,\\
 6(M-1)+(2M-10)^2/24,&M>5.
 \end{cases}
\end{gathered}
\end{equation}
If $M\le5$, $P_6\ge Q_6$ and $P_6''<0$ on $[0,g]$; use the
increasing endpoint or $t^3\le gt^2$ and complete the square.
If $M>5$, $P_6\le6(M-1)$ and the upper quadratic for $Q_6$ gives
the third row. One may also use the actual radius by setting
$r_0=\max(R,2g)>g$ and
\[
 Z_r=\begin{cases}
 r_0(M-1)+(10-2M)t_*-(r_0-g)t_*^2,
 & M\le5,\\
 & t_*=\min\{g,(10-2M)/[2(r_0-g)]\},\\
 r_0(M-1)+(2M-10)t_*-r_0t_*^2,
 & M>5,\\
 & t_*=\min\{g,(2M-10)/(2r_0)\}.
 \end{cases}
\]
Indeed for $M\le5$, $P_{r_0}\ge Q_{r_0}$ and $t^3\le gt^2$;
for $M>5$, $P_{r_0}\le r_0(M-1)$ and dropping $-t^3$ bounds
$Q_{r_0}$. Each displayed quadratic is maximized at its clipped vertex.
Thus the selected implementation takes
$120Z=\min(E_5+10g,Z_0,Z_r)$.
It uses $c_0\ge\sqrt{\pi/(16a_0^3)}=\lambda_N$ and $c_1=c_0/2$.
The weaker choice $c_1\ge\sqrt{\pi/(256a_0^5)}
=\lambda_N N/[2(N-1)]$ is also permitted by the same transfer lemma.

For the paired choice the already proved cell has
$0\le\gamma_l\le|\gamma|\le g\le1$, $M\le2$, and $R>2g$. After reflecting only this auxiliary Fourier integral if necessary, set
\begin{gather*}
 a_*=(2-\gamma_l^2)/24,\qquad b_0=(43/96)(1-2/N),\\
 M_6=\min\{RE_5,1+3g^2+g^4+(R^2+2Rg+3g^2+2)(M-1-\gamma_l^2)\},\\
 t_*=\min\{g,(10-2M)/[2(R-g)]\},\\
 120Z=\min\{E_5+10g,R(M-1)+(10-2M)t_*-(R-g)t_*^2\},\\
 z_l=\max\{-Z,[8\gamma_l-\gamma_l^3-(R+2\gamma_l)(M-1-\gamma_l^2)]/120\},
 \quad z_u=Z,\\
 c_l=(3-M)\gamma_l/144,\qquad c_u=(2-\gamma_l^2)g/144.
\end{gather*}
Check $M-1-\gamma_l^2\ge0$ and that the rectangle is ordered. Here
$P_R\ge Q_R$, and $P_R(t)\le R(M-1)+(10-2M)t-(R-g)t^2$,
whose maximum is at $t_*$. The numerator of the lower $z$ bound increases
in $\gamma\ge0$ because its derivative is
$10-2M+2R\gamma+3\gamma^2>0$. Pearson's first and third identities
give $a_*$ and $M_6$. If $3g^2\le2$, the increasing function
$2\gamma-\gamma^3$ further licenses $c_u\le(2g-g^3)/144$.
Thus this positive rectangle is one specialization of the signed transfer,
with exactly the coefficients used by the fixed acceptance programs.

\subsection{The complete high-frequency integral}
For $b,\kappa>0$ and $n>0$, put
\[
 T_n=\sqrt n\int_b^\infty e^{-n\kappa t^2}t^{-3}\,dt,
 \qquad x=n\kappa b^2.
\]
\begin{lemma}[HIGH\_LAGUERRE2]
For every $x>0$,
\[
 T_n\le\frac{\sqrt n e^{-x}}{2b^2}
                  \frac{x+2}{x^2+4x+2}.
\]
If $N\kappa b^2\ge1/4$, the right side decreases with $n\ge N$.
\end{lemma}
\begin{proof}
Substitute $t^2=b^2(1+s)$ and use positive Fubini to obtain
\[
 T_n=\frac{\sqrt n e^{-x}}{2b^2}
 \int_0^\infty\frac{e^{-xs}}{(1+s)^2}\,ds
 =\frac{\sqrt n e^{-x}}{2b^2}
 \int_0^\infty\frac{u}{x+u}e^{-u}\,du.
\]
Put $P(u)=u^2-4u+2$ and $D(x)=x^2+4x+2=P(-x)>0$.
Exact division gives
\[
 \frac1{x+u}-\frac{Q_x(u)}{D(x)^2}
       =\frac{P(u)^2}{(x+u)D(x)^2}\ge0,
\]
where
\[
 Q_x(u)=-u^3+(x+8)u^2-(x^2+8x+20)u
                  +x^3+8x^2+20x+16.
\]
Integration uses only $\int_0^\infty u^j e^{-u}du=j!$, proved by
integration by parts. It gives
$\int Q_x(u)e^{-u}du=(x+3)D(x)$ and hence
\[
 \int_0^\infty\frac{u}{x+u}e^{-u}du
 =1-x\int_0^\infty\frac{e^{-u}}{x+u}du
 \le1-\frac{x(x+3)}{D(x)}=\frac{x+2}{D(x)}.
\]
For the all-index statement, differentiate
$\sqrt{x}e^{-x}(x+2)/D(x)$. After multiplication by positive factors,
its derivative has numerator
$4-10x-22x^2-13x^3-2x^4$. This polynomial is negative at $x=1/4$
and strictly decreasing for $x>0$. Thus it is negative on the entire
required half-line.
\end{proof}


Let $EX=0$, $EX^2=1$, $EX^4\le M$, and write
$\gamma=EX^3$, $f(t)=Ee^{itX}$. For independent copies suppose an
actual event $E$ satisfies
\[
 P(E)\ge a>0,\qquad |X-X'|\ge d_0>0\text{ on }E,\qquad
 E[(X-X')^2\mid E]\le M_E.
\]
The event is measured under the product of the same marginal law.
Define the Taylor comparator
$e_n(t)=e^{-nt^2/2}(1-in\gamma t^3/6)$ and
\[
 I_\infty=\int_{1/2}^{\infty}|f(t)^n-e_n(t)|t^{-3}\,dt,
 \qquad L(x)=\frac{x+2}{x^2+4x+2}.
\]
The sign of the imaginary term is immaterial for its absolute-value bound.
Use $b=9/10$, $0<r\le1/4$,
\begin{gather*}
 h=\frac{333}{106}-\frac{(22/7)^3r^2}{6}>0,\quad
 \kappa=a(23/48)^2d_0^2,\\
 S_3=\sum_{j=1}^{20}j^{-3}+\frac1{800},\qquad
 S_5=\sum_{j=1}^{20}(j-r)^{-5}+\frac1{4(20-r)^4}.
\end{gather*}
Set
\[
 \begin{gathered}c=\frac{1-\max_{u\in\{1/4,81/100\}}
              [1-u+(M+3)u^2/12]}2,\\
 c_0=\frac18-\frac{M+3}{384},\quad
 s=\frac{1-(M+3)53/600}{2}.\end{gathered}
\]
The usual damping domain requires $c>0$.

\begin{lemma}[HIGH\_WINDOW50]
\label{ext:high-window50}
For every integer $n\ge5$, the complete upper bound is
\begin{align*}
 \sqrt n I_\infty\le{}& A_n
 +\frac{M_E}{4\pi^{3/2}h\sqrt a}
       \left(S_3+\frac{3S_5}{nah^2}\right)
 +\frac{M_E\sqrt n}{2\pi^2}e^{-nah^2r^2}\\
 &+Z_n+\frac{\sqrt n}{2b^2}e^{-na(23/48)^2}
 +\sqrt n e^{-n/8}\left(2L(n/8)+\frac{\sqrt{M-1}}3\right),
\end{align*}
where always $A_n=2\sqrt n e^{-nc}$ is permitted, and, if $s>0$,
it can be replaced by the smaller of itself and
\[
 2\sqrt n e^{-nc_0}L(ns/4).
\]
The zero-cell term $Z_n$ may be taken as the smaller of
\[
 \frac{e^{-n\kappa b^2}}{2b^4\kappa\sqrt n},\qquad
 \frac{\sqrt n e^{-n\kappa b^2}}{2b^2}L(n\kappa b^2).
\]
Each term may be evaluated at a common floor $N$ for all $n\ge N$
provided
\[
 2Nc'>1\quad(c'\in\{c,ah^2r^2,a(23/48)^2,1/8\}).
\]
For the new first term also require $2Nc_0>1$;
for the second zero-cell option require $N\kappa b^2\ge1/4$.
If either extra condition is unavailable, retain the corresponding old term.
Any fixed collection of permitted windows may be evaluated, and the minimum
of its complete six-term bounds remains valid.
\end{lemma}

\begin{proof}
We give the entire frequency partition to make clear what changes.
On $[1/2,b]$, symmetrization and
$\cos u\le1-u^2/2+u^4/24$ yield
\[
 |f(t)|^2\le1-t^2+(M+3)t^4/12,
 \qquad |f(t)|^n\le e^{-n c(t^2)},
 \quad c(u)=u/2-(M+3)u^2/24.
\]
Bounding $c(t^2)$ below by $c$ and using
$\int_{1/2}^{b}t^{-3}dt<2$ gives the old first term.
Concavity supplies the more precise chord bound
\[
 c(u)\ge c_0+s(u-1/4),\qquad 1/4\le u\le81/100,
\]
because the difference is exactly
$(M+3)(u-1/4)(81/100-u)/24$. If $s>0$, extend only the resulting
positive exponential majorant to the half-line. The preceding rational
Gaussian-tail lemma, with lower endpoint $1/2$, gives
\[
 \sqrt n\int_{1/2}^{b}e^{-nc(t^2)}t^{-3}dt
 \le e^{-n(c_0-s/4)}\sqrt n
       \int_{1/2}^{\infty}e^{-ns t^2}t^{-3}dt
 \le2\sqrt n e^{-nc_0}L(ns/4).
\]

For $t\ge b$, conditioning on the actual event and convex Jensen give
\[
 |f(t)|^n\le e^{-nE\sin^2(t(X-X')/2)}
 \le E[e^{-na\sin^2(td/2)}\mid E],\qquad d=|X-X'|.
\]
Positive Tonelli allows a separate partition for every $d>0$.
Put $\ell=2\pi/d$. The zero cell is $[0,\ell/2]$; the other cells are
$\ell[j-1/2,j+1/2]$, $j\ge1$. In each nonzero cell take the core
$t=\ell(j+v)$, $|v|\le r$. These cores are disjoint even for $r\le1/2$.
For our smaller range $r\le1/4$, the elementary sine bound gives
$|\sin(\pi v)|\ge h|v|$ throughout a core, and
$|\sin(\pi v)|\ge hr$ throughout its wings. The latter uses that
$\sin(\pi|v|)$ increases on $[0,1/2]$.

Since $j-r>0$, Taylor's formula on the two sides gives
\[
 \frac{(j-v)^{-3}+(j+v)^{-3}}2
 \le j^{-3}+6(j-r)^{-5}v^2.
\]
Integrate this majorant against $e^{-nah^2v^2}$ on the full real line.
The result is
$\sqrt\pi[j^{-3}/\sqrt{nah^2}
  +3(j-r)^{-5}/(nah^2)^{3/2}]$.
The Jacobian is
$t^{-3}dt=d^2(4\pi^2)^{-1}(j+v)^{-3}dv$.
Summing and averaging $d^2$ gives the second term. The two series tails
are bounded by their decreasing integrals after $20$, producing exactly
$S_3,S_5$. The total wing weight is at most
$\int_{\ell/2}^{\infty}t^{-3}dt=d^2/(2\pi^2)$, giving the third term.
Extending cores or wings below $b$ only adds a positive integral.

In the zero cell write $u=td/2\le\pi/2$.
For $u\le1/2$, use $\sin u\ge(23/24)u$ and the rational
Gaussian-tail lemma with rate $\kappa$; its two choices give $Z_n$.
For $u\ge1/2$, use $\sin u\ge23/48$ and total weight
$1/(2b^2)$, giving the fifth term. Shared endpoints have zero integration
measure, including when the conditional distance distribution has atoms.

Finally keep the comparator on the entire interval $[1/2,\infty)$.
Its absolute-value bound splits its integral into
\[
 \sqrt n\int_{1/2}^{\infty}e^{-nt^2/2}t^{-3}dt
 +\frac{n^{3/2}|\gamma|}{6}
                 \int_{1/2}^{\infty}e^{-nt^2/2}dt.
\]
The first part is at most $2\sqrt n e^{-n/8}L(n/8)$.
For the second, $t\ge1/2$ implies
\[
 \int_{1/2}^{\infty}e^{-nt^2/2}dt
 \le2\int_{1/2}^{\infty}t e^{-nt^2/2}dt
 =2e^{-n/8}/n.
\]
Pearson's elementary identity
$E[(X^2-1-\gamma X)^2]=EX^4-1-\gamma^2\ge0$
gives $|\gamma|\le\sqrt{M-1}$ and proves the sixth term.
Any separately proved upper bound $|\gamma|\le g$ may replace
$\sqrt{M-1}$ in that term. The selected implementation uses
$g=\min(1.85,\sqrt{M-1})$.

For uniformity, $\sqrt n e^{-nc'}$ decreases when $2nc'>1$,
whereas $n^{-1}$ and $n^{-1/2}e^{-nc'}$ always decrease.
The derivative of $L(x)$ is
$-(x^2+4x+6)/(x^2+4x+2)^2<0$.
The new first term is thus decreasing under its stated gate, and the
comparator decreases for every $n\ge5$. The earlier tail lemma proves
the separate $N\kappa b^2\ge1/4$ gate for its new zero-cell option.
These comparisons concern the bounding expressions, so they are uniform
even when the marginal law varies with $n$.
\end{proof}\paragraph{Events, not artificial kernels.}\label{fourier:events}
The entry uses only the following actual choices, with the direct rate
from the same fourth moment in each case:
\[
\begin{array}{c|c|c|c}
 \text{event}&a&d_0&M_E\\\hline
 (X-X')^2\ge9/25 &(2-9/25)^2/(2M+6)&3/5&2/a\\
 XX'<0, D\le d\le.8 &1/(2r_*^2)&\sqrt{2-d}&4r_*^2\\
 XX'<0, p\in[l,u] &2\min(l(1-l),u(1-u))
 &\max(A_l+B_l,\sqrt{2-d})&2/a
\end{array}
\]
The first row follows from Cauchy using
$\E[(X-X')^2;E]\ge2-d_0^2$ and
$\E(X-X')^4=2\mu_4+6$. The second is Lemma~\ref{entry201:event}; all its moment and $\alpha<3$ gates are checked at use.
In the third, $D\le d<2$ ensures a genuine central gap,
$P(XX'<0)=2pq$, and restriction of the squared distance gives $M_E=2/a$.
Taking a minimum of \emph{complete} high-frequency bounds is legitimate. No tail is
dropped and the third construction also allows a cell crossing $1/2$.


\section{Coefficient atlas for the enclosure}\label{re:atlas}
These formulas are inequalities for one actual law. They are the two
coordinate descriptions of the same enclosure: signed moments before a
probability gap is located, and conditional moments afterwards. A bound
is used only with its stated gates, and every refinement is an intersection.
\subsection{Signed moments with a common first moment}
Put $a_\pm=(1\pm\beta)/2$, $b_\pm=(\rho\pm\gamma)/2$.
Cauchy gives $m^2\le a_+a_-$ and $b_\pm m\ge a_\pm^2$.
Consequently
\[
 \frac{(1+\beta)^{3/2}}{\sqrt{1-\beta}}-\rho
 \le\gamma\le
 \rho-\frac{(1-\beta)^{3/2}}{\sqrt{1+\beta}},\qquad
 m\ge\max\left\{\frac{1+\beta^2}{2\rho},
                  \frac{a_+^2}{b_+},\frac{a_-^2}{b_-}\right\}.
\]
These statements also hold with a zero atom. Projection of $X|X|$ onto
the span of $X$ and $|X|$ gives
\[
 \mu_4\ge\rho^2+\frac{(\gamma-\rho\beta)^2}{1-\beta^2}.
\]
For completeness, this follows by minimizing the nonnegative expectation
$E(X|X|-uX-v|X|)^2$ over $u,v$; its Gram matrix has diagonal entries
one and off-diagonal entry $\beta$, and is positive definite.

Two support-localizing moment matrices give the additional upper bounds
as follows. If
$D_\pm=Ra_\pm-b_\pm>0$, then
\[
 2m+\sum_\pm\frac{(Rm-a_\pm)^2}{D_\pm}\le R,
 \qquad
 \mu_4\le R\rho-\sum_\pm\frac{D_\pm^2}{Rm-a_\pm}.
\]
Indeed the first follows by summing the determinant inequalities for
\[E[(R-|X|)(1,|X|)^{\mathsf T}(1,|X|);\pm X>0].\]
The second follows from the corresponding matrices with an additional
factor $|X|$. One more factor $|X|$ yields, sign by sign,
\[
 E[X_\pm^5]\le R E[X_\pm^4]
 -\frac{(Rb_\pm-E[X_\pm^4])^2}{D_\pm}.
\]
Zero denominators are treated by dropping their nonnegative correction;
they never imply exclusion without a separate contradictory numerator.
The common $m$ is retained throughout.

A second useful upper bound for $m$ follows by integrating
$(x-a)^2(R-x)\ge0$ for $x=|X|\in[0,R]$. For every fixed $a>0$,
\[
 2m\le\frac{R-\rho+2a+Ra^2}{a(2R+a)}.
\]
The cover uses only the eleven rational choices $a=2/10,\ldots,12/10$.

Write $\delta=\gamma-\rho\beta$. Cauchy and a square completion on
each sign give the signed constant bound
\begin{align*}
 \frac\rho2\mu_4-\frac14\eta_5
 &\le\frac{\rho^3}{4}
 -\frac{(\rho+\gamma)(\delta_+)^2}{8(1+\beta)^2}
 -\frac{(\rho-\gamma)((-\delta)_+)^2}{8(1-\beta)^2}.
\end{align*}
Indeed on either sign, write $a=\E X_s^2$, $b=\E X_s^3$,
$c=\E X_s^4$, $f=\E X_s^5$. Cauchy gives $c\ge b^2/a$,
$f\ge c^2/b$, and
$\rho c/2-f/4\le\rho^2b/4-(c-\rho b)^2/(4b)$.
Distance to the half-line $c\ge b^2/a$ proves the stated penalty
on summing both signs. Independently,
$\mu_4^2\le\rho\eta_5$ and the projection bound above give
\[
 \frac\rho2\mu_4-\frac14\eta_5
 \le\frac{\rho^3}{4}
 -\frac{(\gamma-\rho\beta)^4}{4\rho(1-\beta^2)^2}.
\]
The larger of the two nonnegative deductions is retained.

\subsection{An actual central gap without assuming \texorpdfstring{$D<2$}{D below 2}}
If zero belongs to the support, the signwise gap calculation gives $J\le\rho$.
Thus $D\ge\sqrt{10}\rho-\gamma\ge(\sqrt{10}-1)\rho$.
Whenever $D<(\sqrt{10}-1)\rho$, there is therefore a central gap
$(-A,B)$, $A,B>0$, with no zero atom. Put $p=P(X>0)$, $q=1-p$,
$T_+=b_+-mB^2$ and $T_-=b_--mA^2$. The exact gap decomposition is
\[
 D=mQ+(\sigma-1)T_++(\sigma+1)T_-+3(T-J_o),\qquad
 Q=(\sqrt{\sigma+1}A-\sqrt{\sigma-1}B)^2,
 \quad\sigma=\sqrt{10}.
\]
All terms are nonnegative. Hence $T\le D/(\sigma-1)$ and
\[
 (A+B)^2\ge A^2+B^2\ge\frac{\rho-D/(\sigma-1)}m.
\]
Moreover $A\le a_-/m$, $B\le a_+/m$ and $J\le\rho+2m$.
The last bound gives the compulsory defect
\[
 D\ge \sigma\rho-\gamma-6m.
\]
The inequalities $T-J_o\ge2BR_++2AR_-$ give the useful additional
endpoint constraint
\[
 D\ge(\sigma-1)b_++6a_+B-(\sigma+5)mB^2.
\]
If $D<(\sigma-1)b_+$, this downward quadratic forces $B$ to be at
least its positive root. The reflected constraint has $a_-,b_-,A$
and the coefficients $\sigma+1,\sigma+7$. In a closed moment box,
positive coefficients are rounded down, negative coefficients up, and
the lower square-root endpoint is used in the resulting root bound.

The actual opposite-sign event has probability $2pq$.
The first two moment constraints give lower bounds for both $p$ and $q$.
For independent copies one has exactly
\[
 E[(X-X')^2\mid XX'<0]
 =\frac{1+\beta(1-2p)+4m^2}{2pq}
 =\frac{1+4m^2+\beta}{2p}
  +\frac{1+4m^2-\beta}{2q}.
\]
For an upper bound, use the upper $m$ endpoint and both endpoints of
$p$ and $\beta$. The expression is convex in $p$ and affine in $\beta$.
This improves the former $2/P(XX'<0)$ bound inside the complete HIGH
estimate. Both copies still have the same actual marginal.

\subsection{Preserving central and outer integral costs}
Let
\[
 P_1(A,B)=B^3/3+AB^2/2-A^3/6,\quad
 P_2(A,B)=B^4/4+AB^3/3-A^4/12,\quad G=\rho P_1-P_2.
\]
Then $J_{1,c}=mP_1$, $J_{2,c}=mP_2$. Since
$\partial_BG=B(A+B)(\rho-B)$,
\[
 G(A,B;\rho)\le G(A,\rho;\rho)
 =\rho^4/12+\rho^3A/6-\rho A^3/6+A^4/12.
\]
The last expression increases with $A$ because its derivative is
$(\rho-A)^2(\rho+2A)/6$, and is convex in $\rho$ because its second
$\rho$ derivative is $\rho^2+\rho A$. Its closed-rectangle maximum is
therefore bounded using the upper $A$ endpoint and both $\rho$ endpoints.

Positive outer gaps satisfy
$(t-a)(\rho-t)\le(\rho-B)_+^2/4$. On negative outer gaps the positive
part is bounded by $(R-\rho)^2/4$ after extracting $|t|$.
Integration gives
\[
 3(\rho J_1-J_2)_o
 \le\frac38\{(\rho-B)_+^2T_++(R-\rho)^2T_-\}.
\]
The available weighted outer budget is $D-mQ$, not the entire $D$.
The cover intersects this budget with the individual bounds
$T_+\le b_{+,u}-m_lB_l^2$ and $T_-\le b_{-,u}-m_lA_l^2$.

For the signed phase term $U=J_1+(\mu_4-3)/6$, retain
\[
 -RT_-/2\le J_{1,o}\le RT_+/2.
\]
The central minimum of $P_1$ occurs at $B_l$ and an $A$ endpoint;
its maximum occurs at $B_u$ and $A$ clipped to the interval at $A=B_u$.
A nonnegative polynomial bound is multiplied by $m_u$ for an upper
bound and $m_l$ for a lower bound; a negative one uses the opposite
endpoint. For example, the binary endpoints $A=3/2$, $B=2/3$ make this
maximum negative; the sign switch is necessary even when $D<2$.
The fourth-moment sign intervals also sharpen
$K_1=E[X|X|^3]-3\beta-3\rho\gamma/2$.
The complete $K_1,U,\phi$ rectangles are used in the signed expression
$|CK_1+\phi U|$, without replacing it by two separate absolute values.


\paragraph{A complete coarse bound before a central gap is licensed.}
Always $-\mu_4/6\le J_1\le\mu_4/2$ and
$-3J_2-\eta_5/4\le0$. On a gap of length $h=b-a$, these follow from
\[
 \int_a^b(t-a)t\,dt+\tfrac12\int_a^bt^2dt=hb^2/2,\quad
 \tfrac32\int_a^bt^2dt-\int_a^b(t-a)t\,dt=h(2a+b)^2/6,
\]
and, on a negative gap $(-v,-u)$, from
$\tfrac13\int_u^v x^3dx-\int_u^v(v-x)x^2dx=u^3(v-u)/3$.
For $1\le r\le\rho$ and actual support radius $R\le6$, set
\[
 \Gamma(R,r)\ge\max_{r\le v\le R}
 \left\{-\frac{v^2}4+\frac{(R+r)v}3-\frac{Rr}2+
                      \frac{r^3(2R-r)}{12v^2}\right\}_+.
\]
Then
\[
 3(\rho J_1-J_2)\le\frac32\Gamma(R,r)(\rho-\gamma)
                              +\frac38\rho^2J.
\]
On negative gaps extract $x=|t|$ and discard the nonpositive part
$x<\rho$. Replacing $\rho$ by $r$ and the outer endpoint by $R$
gives $\int_r^v(R-x)x(x-r)\,dx$. Dividing this integral by $v^2$
gives exactly the displayed expression. Signwise Tonelli now charges
$\Gamma\E X_-^3=\Gamma(\rho-\gamma)/2$.
On positive gaps $(t-a)(\rho-t)\le\rho^2/4$, giving $\rho^2J/8$.
The implementation uses 64 closed $v$-intervals: on each, maximize the
concave quadratic at its clipped vertex $2(R+r)/3$ and the positive
reciprocal term at the left endpoint. This proves its interval bound.
The alternative bound $\Gamma\le5/2$ follows at $r=1,R=6$ by
$3v^4-28v^3+66v^2-11>0$ on $[1,6]$: the endpoint values are $30,205$,
and the interior critical values are $(85\mp35\sqrt5)/2>0$.
These monotone enclosing parameters cover all $r\ge1,R\le6$.
\subsection{Moment and gap coefficients for a whole cell}
\label{mom:conditional-boxes}

The following coefficient rules are collected here to make the finite
calculation transparent. Every rule is an inequality for an actual law;
there is no extremal-law assumption inside a cell.
Write $v=1-\theta$. For $p\in[l,u]\subset(0,1)$ set
\begin{gather*}
 s_0=\min\{l(1-l),u(1-u)\},\quad
 v_* =\max\{\sqrt{u/(1-u)},\sqrt{(1-l)/l}\},\\
 r_B=(1-2s_0)/\sqrt{s_0},\quad
 g_B=\max_{t=l,u}|1-2t|/\sqrt{t(1-t)},\quad b_B=\max_{t=l,u}|1-2t|,\\
 L_3=R+2v_*,\quad L_4=R^2+2Rv_*+3v_*^2,\\
 L_5=R^3+2R^2v_*+3Rv_*^2+4v_*^3,\quad
 e_B=r_B\frac{1-4s_0+2s_0^2}{s_0(1-2s_0)}.
\end{gather*}
The polynomials $T_3=x+2a$, $T_4=x^2+2ax+3a^2$ and
$T_5=x^3+2ax^2+3a^2x+4a^3$ are nonnegative. Their maxima at
$x=R$, $a=v_*$, in Lemma~\ref{re:pricing}, give, if $\theta\le\Theta<1$,
\begin{align}
 \rho&\le r_B+(L_3-r_B)\Theta,&
 |\gamma|&\le g_B+(L_3-g_B)\Theta,\notag\\
 |\beta|&\le b_B+(1-b_B)\Theta,&
 \mu_4&\le(1-\Theta)^2(1/s_0-3)+L_4\Theta,\notag\\
 \eta_5&\le(1-\Theta)^2e_B+L_5\Theta.&&\label{mom:boxes}
\end{align}
Replacing $\theta$ by $\Theta$ uses the explicitly checked inequalities
$L_3>r_B,g_B$, $L_4>2(1/s_0-3)$ and $L_5>2e_B$.
The exact binary factors are $v^{3/2},v,v^2,v^{5/2}$; the upper bounds
use $v^{3/2}\le v$ and $v^{5/2}\le v^2$. Thus cells crossing $1/2$ and
zero conditional variances are included. The same cubic expansion, in the
later domain $p\in[.37,.47]$, $A\ge.75$, $B\ge1.05$, gives
\begin{equation}\label{mom:rho-lower}
 \rho\ge v^{3/2}\rho_B(p)+3\min(A,B)\theta
       \ge\rho_B(p)+.45\theta,\qquad
 \rho_B(p)=\frac{1-2pq}{\sqrt{pq}}.
\end{equation}
Indeed $\rho_B<1.11$ there and $v^{3/2}\ge1-3\theta/2$.
Also $\beta\ge1-2p-2(1-p)\theta$.

For the signed coefficient
$K_1=\E[X|X|^3]-3\beta-3\rho\gamma/2$, define
\[
 K(p,v)=(1-2p)\{(3v^3/2-v^2)/(pq)-3v^3+v^2+3v\}.
\]
Direct substitution gives $K_1(V)=-K(p,v)$, with the coefficient $3$
unchanged. For $v\ge2/3$, $K$ decreases in $p$: both coefficients in
braces are nonnegative and $(1-2p)/(pq)$ decreases. Its derivative in $v$
has the sign of $1-2p$, since the derivative of the braces is at least
$9v^2-6v+3>0$. Hence both variance endpoints are retained. The error is
\begin{equation}\label{mom:K-error}
 |K_1(X)-K_1(V)|\le
 \Theta\{L_4+3+\tfrac32L_3(R_\rho+g_B)\}
\end{equation}
when $\rho\le R_\rho$. This follows from the conditional expansions and
$\rho\gamma-\rho_V\gamma_V=
\rho(\gamma-\gamma_V)+\gamma_V(\rho-\rho_V)$.

Here are the remaining signed charges. They also explain the two
alternative bounds retained in the initial, coarse moment cover:
\begin{align}
 -\mu_4/6&\le J_1\le\mu_4/2,&
 -3J_2-\eta_5/4&\le0,                                  \label{mom:gap-signs}\\
 \rho J_1-J_2&\le\tfrac54(\rho-\gamma)+\tfrac18\rho^2(\rho+1),
                                                              \label{mom:oldgap}\\
 -75D/104&\le J_{1,o}\le3D/2.                           \label{mom:outer}
\end{align}
To prove the first line, a gap of length $h=b-a$ satisfies
\[
 \int_a^b(t-a)t\,dt+\tfrac12\int_a^bt^2dt=hb^2/2,\quad
 \tfrac32\int_a^bt^2dt-\int_a^b(t-a)t\,dt=h(2a+b)^2/6.
\]
Sum with weight $\tau_g$. On a negative piece $(-v,-u)$,
$\frac13\int_u^v x^3dx-\int_u^v(v-x)x^2dx=u^3(v-u)/3\ge0$;
Tonelli proves the $J_2$ assertion. For~\eqref{mom:oldgap}, the positive
part is at most $\rho^2J/8$. The negative part is at most
$\E[V\int_1^V(6-x)x(x-1)dx;V\ge1]$, $V=X_-$, hence at most
$(5/2)\E X_-^3$. The last comparison is equivalent to
$3V^4-28V^3+66V^2-11>0$ on $[1,6]$: the endpoint values are $30,205$,
and at the possible interior critical points $(7\pm\sqrt5)/2$ its value
is $(85\mp35\sqrt5)/2>0$.
For~\eqref{mom:outer}, use $J_{1,o}\ge-3T_-$,
$J_{1,o}\le3J_o\le3T$ and~\eqref{re:gapexact}.

The central integrals are $J_{j,c}=mP_j(A,B)$, with
\[
 P_1=B^3/3+AB^2/2-A^3/6,\qquad
 P_2=B^4/4+AB^3/3-A^4/12.
\]
They increase in $B$; their $A$ derivative changes from positive to negative
at $A=B$. A rectangle minimum is therefore at $B_l$ and an $A$ endpoint.
Multiply a negative minimum by an upper bound for $m$, and a nonnegative
minimum by a lower bound. For the upper central charge, maximizing at
$(\bar a,\bar b)$ if $p\le1/2$ gives
$J_{1,c}\le v^2(2-3p)/(6pq)$, a decreasing function of $p$ on that interval.
If $p\ge1/2$, maximize first at $A=B$ to get
$J_{1,c}\le(2/3)v^2q^2/p\le1/3$. Thus every cell may use
\begin{equation}\label{mom:central-upper}
 J_{1,c}\le\max\{1/3,(2-3l)/(6l(1-l))\}.
\end{equation}

For $|X|\le R$, $R\ge1$, define
\[
 k_4(R)=\frac{(R+1)^2}{R+1/2},\qquad
 k_5(R)=\frac{R^3+2R^2+3R+3/2}{R+1/2}.
\]
Then
\begin{equation}\label{mom:universal-moments}
 \mu_4\le\min\{R\rho,1+k_4(R)(\rho-1)\},\quad
 \eta_5\le\min\{R\mu_4,1+k_5(R)(\rho-1)\},\quad
 \mu_6\le R\eta_5.
\end{equation}
Writing $u=|X|$ in~\eqref{re:pricing}, the first two constants
are the endpoint values of the ratios
$u+3/2+1/[4(u+1/2)]$ and
$u^2+3u/2+9/4+3/[8(u+1/2)]$, whose derivatives are nonnegative.
The other bounds are pointwise. The following coarse cancellation estimates
continue to use the valid radius-six majorants:
The last identity in~\eqref{re:pricing} gives
\begin{equation}\label{mom:K-cancellation}
 |K_1+3\rho\beta|\le(96-18\rho)(\rho-1)/13.
\end{equation}
Indeed $u(u+2-3r/2)/(u+1/2)$ is convex on $[0,6]$, its maximum is
$(96-18r)/13$, and its negative part is at most $31/40$, smaller than that
maximum for $1\le r\le37/20$.
The following bounds are used only for $|X|\le6$ and
$1\le\rho\le37/20<2$, with the signed moment hypotheses below.
The other bounds needed on this domain are
\begin{align}
 |K_1|&\le6\rho-3\rho^2/2+3|\beta|,\notag\\
 K_1&\le M+3|\beta|-3\rho g_-/2,&
 -K_1&\le3\rho+3|\beta|+(3\rho/2-3)g_-,\notag\\
 -K_q&\le(\rho/2-3/8)M-\rho/2+1/8+3|\beta|g/2,&
 -K_q&\le8\rho^3/27+(3|\beta|/2-1/2)\rho,               \label{mom:K-bounds}
\end{align}
where $\mu_4\le M$, $|\gamma|\le g$, $\gamma\ge g_-$, and
$K_q=\eta_5/4-\rho\mu_4/2-3\beta\gamma/2+\rho/2$.
For the first two lines use
$|\E[X|X|^3]-3\gamma|\le3\rho$, and note the negative coefficients
of $\gamma$. For the last line use
$\eta_5\ge\mu_4^{3/2}\ge3\mu_4/2-1/2$, or maximize
$\rho y^2/2-y^3/4$ over $y\ge0$.


\paragraph{Compatible weaker resource constants.}
The implementation also permits the weaker binary charge
$b_I=(616/25)(1-\Theta)^2\operatorname{dist}(p_*,I)^2$.
It follows from $2(6+2\sigma)>616/25$ and
$(1-\theta)^{3/2}\ge(1-\theta)^2$.
With $K=\min(K_+,K_-)$ and $v_* =\max(v_+,v_-)$, the merged charge
\eqref{re:residual} gives
\[
 T_++T_-\le T_*=(d-b_I)\max\{v_*/[3(A_l+B_l)],(R+2v_*)/K\},
 \qquad d_o=\min\{d,(\sigma+1)T_*\}.
\]
Thus $W_o\le d_o$ and, for $R\le6$,
\[
 -75d_o/104\le J_{1,o}\le3d_o/2,\qquad
 3(\rho J_{1,o}-J_{2,o})\le\frac38
 \max\{\rho^2/(\sigma-1),(R-\rho)^2/(\sigma+1)\}d_o.
\]
These are precisely the merged-cost versions of the sign-specific
resource rule; they do not introduce another available budget.

\subsubsection{Conditional lower moments improve two signed charges}

Let $s_u=\max_{p\in[l,u]}pq$ and $\theta\le\Theta$.
The unsigned cases of~\eqref{re:conditional} give
\begin{align}
 \mu_4&\ge M_l:=\max\{1,(1-\Theta)^2(1/s_u-3)\},\label{entry201:m4lower}\\
 \eta_5&\ge E_l:=
 (1-\Theta)^{5/2}\frac{1-4s_u+2s_u^2}{s_u^{3/2}}.\label{entry201:e5lower}
\end{align}
Both binary moment functions decrease in $s\in(0,1/4]$.
For the fifth this follows because the derivative has the sign of
$s^2+2s-3/2<0$; the fourth is immediate.
The fourth-moment lower bound sharpens the lower endpoint of
$J_1+(\mu_4-3)/6$; its explicit use appears in the cell map below.

The constant signed charge is
\[
 -K_q=\frac{\rho\mu_4}{2}-\frac{\eta_5}{4}
       +\frac32\beta\gamma-\frac{\rho}{2}.
\]
The old coherent estimate used $\eta_5\ge3\mu_4/2-1/2$ and
$\mu_4\le M$. A second estimate uses \eqref{entry201:e5lower}.
Their difference is independent of $\rho$. Consequently the improved
bound is exactly the old one minus
\begin{equation}\label{entry201:eta-improvement}
 \frac14\max\{0,E_l-3M/2+1/2\}.
\end{equation}
This is a constant shift on the cell and does not spoil convexity.
It does not assert the generally invalid implication
$\eta_5\ge3M/2-1/2$ from an upper bound on $\mu_4$.

\subsubsection{Signed conditional Taylor cancellation}

Let $K_1=\E[X|X|^3]-3\beta-\tfrac32\rho\gamma$.
On a sign class put $x=|X|$, $a=|V|$, $W=X-V$ and
$\varepsilon=1$ on the positive class, $-1$ on the negative class.
In the actual central-gap domain there is no atom at zero.
With $w_3=x+2a$, $w_4=x^2+2ax+3a^2$ and
$\gamma_V=\E V^3$, exact conditional expansion gives
\begin{equation}\label{entry201:Kexact}
 K_1(X)-K_1(V)
 =\E\!\left[W^2
 \left\{\varepsilon(w_4-3-\tfrac32\rho w_3)
                   -\tfrac32\gamma_Vw_3\right\}\right].
\end{equation}
Apply~\eqref{re:conditional} to $j=2,3,4$, with $h=1$
or $h=\varepsilon$, and use
$\rho\gamma-\rho_V\gamma_V
=\rho(\gamma-\gamma_V)+\gamma_V(\rho-\rho_V)$.

The positive coefficient in \eqref{entry201:Kexact} is
$F(x,a;t)=x^2+2ax+3a^2-3-\tfrac32t(x+2a)$
with $t=\rho+\gamma_V$; the negative coefficient is $-F$
with $t=\rho-\gamma_V$. Their ranges are bounded on complete
rectangles of actual conditional means and support magnitudes.
The positive-support magnitude is at most $R$.
The same formula permits separate actual bounds on the two signed support magnitudes.
On the late domain $p<1/2$, bounds for $\gamma_V$ are
\[
 \sqrt{(1-\Theta)^3(1-2u)^2/[u(1-u)]}
 \le\gamma_V\le(1-2l)/\sqrt{l(1-l)}.
\]
The positive mean lies in
$[\sqrt{(1-\Theta)(1-u)/u},\sqrt{(1-l)/l}]$, and the negative
magnitude mean lies in
$[\sqrt{(1-\Theta)l/(1-l)},\sqrt{u/(1-u)}]$.
Use the lower bound $\rho\ge\rho_B(u)$ only after checking
$3\min(A_l,B_l)\ge(3/2)\sup_I\rho_B$. Indeed conditional expansion
and $(1-\theta)^{3/2}\ge1-3\theta/2$ then imply
$\rho\ge(1-\theta)^{3/2}\rho_B(p)+3\min(A,B)\theta\ge\rho_B(p)$.
Otherwise retain the universal lower bound one.

Since $x+2a>0$, $F$ decreases with $t$.
For fixed $t$ its Hessian is
$\left(\begin{smallmatrix}2&2\\2&6\end{smallmatrix}\right)>0$.
Its rectangle maximum is at a corner. Its minimum is among the four
corners, the minimizer on each edge, and the stationary point
$(x,a)=(3t/8,3t/8)$ if that lies in the rectangle.
On fixed-$x$ edges use $a=t/2-x/3$, clipped to the edge; on
fixed-$a$ edges use $x=3t/4-a$, also clipped.
This is a complete finite extremum list.
Multiplication by $0\le\theta\le\Theta$ uses
$\Theta\min(0,L)$ and $\Theta\max(0,U)$ for the lower and upper
coefficient bounds. It intersects this new interval with the retained
absolute-error interval; neither need be discarded.

\subsection{Generic event and complete central rectangle}
\begin{lemma}[Generic small-defect event]\label{entry201:event}\label{mom:third}
Assume $|X|\le R\le6$ and $\rho\le37/20$. Let $0\le d<2$ and $D\le d$. Define
\[
 \begin{gathered}
 r_*=\min\{37/20,(3+d)/(\sigma-1)\},\qquad
 W=\sqrt{37d/10},\\
 c=\sqrt{\sigma+1},\qquad e=\sqrt{\sigma-1},\qquad C_0=c+e,
 \end{gathered}
\]
\[
 h_*=\frac{C_0}{6}\sqrt{12+W^2}+\frac{W}{3C_0},\qquad
 L_*=\max\left\{h_*,\frac{(2c-e)h_*+3W}{C_0}\right\}.
\]
Then $\rho\le r_*$, $A+B\le h_*$ and $|2B-A|\le L_*$.
If $\alpha=L_*/6\le3$, then
\[
 |J_1|\le\alpha(\rho+1)+(3-\alpha)d/2,\quad
 \mu_4\le h_*\rho+Rd/2,\quad
 \eta_5\le h_*^2\rho+R^2d/2.
\]
For independent copies, the actual event $E=\{XX'<0\}$ obeys
\[
 P(E)\ge a:=1/(2r_*^2),\quad
 |X-X'|\ge d_0:=\sqrt{2-d}\text{ on }E,\quad
 \E[(X-X')^2\mid E]\le2/a.
\]
\end{lemma}
\begin{proof}
The inequality $J\le\rho+1$ gives $D\ge(\sigma-1)\rho-3$.
Equation \eqref{re:gapexact} gives $w^2\le2\rho D\le37d/10$ for
$w=cA-eB$. Since $ce=3$,
\[
 A+B=\frac{C_0}{6}\sqrt{w^2+12AB}-\frac{w}{3C_0},
 \qquad 2B-A=\frac{(2c-e)(A+B)-3w}{C_0}.
\]
Use $AB\le1$, $|w|\le W$; also $2B-A\ge-(A+B)$.
Now $J_{1,c}=J_c(2B-A)/6$, $|J_{1,o}|\le3J_o$ and
$0\le J_o\le T\le d/2$ imply
$|J_1|\le\alpha J+(3-\alpha)J_o$.

For either inner endpoint $b_*$ and $b_*\le x\le R$,
\[
 x^4-h_*x^3\le R x(x^2-b_*^2),\qquad
 x^5-h_*^2x^3\le R^2x(x^2-b_*^2).
\]
Indeed $h_*\ge b_*$, and after factoring $x-b_*$ the remaining
comparisons follow from $x\le R$. Integrate and use $T\le d/2$.
Finally $m^2\le pq$ and $m\ge1/(2\rho)$ give $2pq\ge a$;
the conditional squared distance is bounded by its unconditional
expectation $2$, divided by the event probability.
\end{proof}

The central part of $\rho J_1-J_2$ is exactly
$mG(A,B;\rho)$, where $G=\rho P_1-P_2$.
Let a proved probability cell give
$A\in[A_l,A_u]$, $B\in[B_l,B_u]$, $m\in[m_l,m_u]$.
Here
\[
 A_u=\sqrt{u/(1-u)},\quad B_u=\sqrt{(1-l)/l},\quad
 m_l=\sqrt{(1-\Theta)s_0},\quad
 m_u=\sqrt{\max_{p\in[l,u]}pq}.
\]

\paragraph{The closed-box Bernstein identity.}\label{re:bernstein}
For a polynomial $P(x)=\sum_{k=0}^dt_kx^k$,
\[
 P(x)=\sum_{i=0}^d b_i\binom di x^i(1-x)^{d-i},\qquad
 b_i=\sum_{k\le i}t_k\binom ik/\binom dk.
\]
Indeed $\binom di\binom ik=\binom dk\binom{d-k}{i-k}$ and the
binomial theorem recover each monomial $x^k$. The basis terms are
nonnegative and sum to one on the closed unit interval. Tensor products
prove the two-variable version below, including all boundary points.

\begin{lemma}[Complete central rectangle]\label{entry201:central}
Write $G(A_l+(A_u-A_l)x,B_l+(B_u-B_l)y;\rho)
=\sum t_{kh}x^ky^h$. Define the 25 rational coefficients
\[
 b_{ij}(\rho)=\sum_{\substack{k\le i\\h\le j}}
 t_{kh}\frac{\binom{i}{k}}{\binom4k}
          \frac{\binom{j}{h}}{\binom4h},\qquad 0\le i,j\le4,
\]
and
\[
 \mathcal G(\rho)=\max_{i,j}
       \{m_l b_{ij}(\rho),\,m_u b_{ij}(\rho)\}.
\]
Then $\rho J_{1,c}-J_{2,c}\le\mathcal G(\rho)$.
Moreover $\mathcal G$ is convex in $\rho$.
\end{lemma}
\begin{proof}
Apply the Bernstein identity~\eqref{re:bernstein} separately to the
two coordinates after the affine substitution. The Bernstein basis is nonnegative and sums to one on the closed square,
so $G$ is a convex combination of its 25 coefficients.
For each coefficient retain both $m$ endpoints, or equivalently use
$m_u$ when it is nonnegative and $m_l$ otherwise. This gives a uniform
bound on the entire three-dimensional rectangle, including boundaries.
Each endpoint-multiplied coefficient is affine in $\rho$; their maximum
is convex. No interior stationary point is omitted.
\end{proof}

\subsection{One interval-feedback theorem}
\label{entry:feedback}\label{fourier:uniform}

The signed identity from the preceding section is
\begin{equation}\label{entry:remainder}
 \cR=C\{3\beta J/2+3\rho J_1-3J_2-K_q\}
 +H(CK_1+\varphi U)-H^2(C\rho+\gamma\varphi/2)+E_n,
\end{equation}
where $U=J_1+(\mu_4-3)/6$, $\varphi=\phi(H/\sqrt n)$.
For any one coherent low-frequency triple and complete high-frequency value $K_H$ at floor $N$,
\begin{equation}\label{entry:fourier-budget}
 |E_n|\le K_F+s_F|H|,\quad
 K_F=c_\pi(L_0+K_H+L_2/\sqrt N),\quad s_F=c_\pi L_H,\quad c_\pi\ge2/\pi.
\end{equation}
The implementation uses $16/25$ before the signed phase and $212/333$
afterward; both are valid since $\pi>333/106$.

\begin{lemma}[Continuous cells and feedback]\label{entry:cell-theorem}\label{entry:corners}
Suppose bounds for one actual cell give
\[
 r_0\le\rho\le r_1,\quad0<c_0\le C\le c_1,\quad
 |\beta|\le b,\quad |\gamma|\le g,\quad \mu_4\le M,\quad D\le d,
\]
$K_1\in[k_l,k_u]$, $U\in[u_l,u_u]$, and
$0\le\varphi_-\le\varphi\le2/5$. Freeze the shared outer charge $d_o$,
lower fifth moment $E_l$, upper $s_u=\max_{p\in[l,u]}p(1-p)$, and
central envelope $\mathcal G$ from the preceding bounds.
Define
\begin{align*}
 A(r)&=\tfrac32 b(r+2\sqrt{s_u})+3\mathcal G(r)
 +\tfrac38\max\{r^2/(\sigma-1),(R-r)^2/(\sigma+1)\}d_o\\
 &\quad +(r/2-3/8)M-r/2+1/8+\tfrac32bg
 -\tfrac14\max\{0,E_l-3M/2+1/2\},\\
 B(c)&=\max_{k\in\{k_l,k_u\},\,u\in\{u_l,u_u\},\,
                      \psi\in\{\varphi_-,2/5\}}|ck+\psi u|,\\
 Q(c,r)&=cr+\tfrac12\varphi_-(3.16r-3-d)
       -\tfrac12(2/5-\varphi_-)(d-.16)_+.
\end{align*}
If all four corner values of $Q$ are positive, then, for every real $H$,
\begin{align}
 \cR&\le\max_{(c,r)\in\{c_0,c_1\}\times\{r_0,r_1\}}
       \{K_F+cA(r)+(B(c)+s_F)|H|-Q(c,r)H^2\}\notag\\
 &\le\max_{c,r\text{ corners}}
       \{K_F+cA(r)+(B(c)+s_F)^2/(4Q(c,r))\}.           \label{entry:four-corners}
\end{align}
If $T$ is a strict upper bound for the final maximum in
\eqref{entry:four-corners}, the weight-removal step in
Lemma~\ref{re:weight} gives
\begin{equation}\label{entry:feedback-values}
 D<\frac{15000T}{997N},\quad
 \varphi>\frac{997}{2500}\left(1-\frac{5T}{2N}\right),\quad
 C<\CE+V/N\quad\text{if also }\cR/\rho<V.
\end{equation}
Only positive density floors are used. These statements remain valid at
$H=0$, on closed cell edges, and for zero conditional variance.
\end{lemma}
\begin{proof}
The constant coefficient follows from the shared outer budget,
Lemma~\ref{entry201:central}, and~\eqref{mom:K-bounds}, using
\[
 J\le\rho+2ABm\le\rho+2(1-\theta)^{3/2}\sqrt{pq}
       \le\rho+2\sqrt{s_u}.
\]
This follows from $A\le\bar a$, $B\le\bar b$, $m^2=(1-\theta)pq$;
it is used only in the charge $3|\beta|J/2$. Since $\gamma\ge3.16\rho-3-d$, its negative
part is at most $(d-.16)_+$. Using the upper density only for that negative
part gives $C\rho+\gamma\varphi/2\ge Q$. Absolute value is separately
convex in each coordinate of $ck+\psi u$, proving the formula for $B$.
For fixed $r$ and $|H|$, the resulting upper bound is convex in $c$.
Take those endpoints \emph{before} completing the square. At either
endpoint $c>0$, it is convex in $r$, since $\mathcal G$ and the maximum of the two outer quadratics are
convex, every other term of $A$ is affine, and $Q$ is affine in $r$. This proves the first line of~\eqref{entry:four-corners};
$fh-qh^2\le f^2/(4q)$ proves the second. Bilinearity of $Q$ proves
its positivity throughout the rectangle from its corners.

\end{proof}

\subsection{Signed phase curvature in the resource schedule}
\label{re:signed-phase-repair}
The auxiliary Fourier orientation does not change the signed third moment
in the actual transfer identity. The paired phase estimate must therefore
retain its sign. On a fixed resource cell put
\[
 L=\frac{997}{2500}\left(1-\frac{5r}{2N}\right),\qquad
 g_0(\rho)=3.16\rho-3-d,\qquad U=.4,
\]
where $r,d$ are the preceding raw and defect upper bounds, and the
checker requires $0\le L\le U$. The actual density lies in $[L,U]$
and the actual signed moment satisfies $\gamma\ge g_0(\rho)$. Consequently
\begin{equation}\label{re:sign-aware-price}
 C\rho+\tfrac12\gamma\phi
 \ge q_{\rm s}(C,\rho):=
 C\rho+\tfrac12\min\{L g_0(\rho),U g_0(\rho)\}.
\end{equation}
This follows by first multiplying the lower bound on $\gamma$ by the
nonnegative actual density, and then minimizing over its two endpoints.
The simpler expression $C\rho+Lg_0(\rho)/2$ is licensed if $g_0\ge0$,
or if nonnegativity of the actual $\gamma$ has been separately proved.
It is not licensed merely by reflecting the auxiliary Fourier integral.

The corrected price preserves the four-corner reduction. For fixed $C$,
$q_{\rm s}$ is concave in $\rho$, being affine plus the minimum of two
affine functions. For fixed $\rho$ it is affine in $C$. In the phase
envelope $A+f|H|-q_{\rm s}H^2$, the remaining terms are separately
convex in $C,\rho$, by the moment and coefficient bounds already proved.
Thus the entire envelope is separately convex, and two successive
one-dimensional endpoint bounds place its maximum at the four corners.
The same argument applies after adding the affine conserved-resource
curvature $e_\rho$ and subtracting a barrier price $K\rho$.
If the four corrected divisors are positive, concavity makes them
positive throughout the rectangle, and square completion gives the
valid raw and resource bounds
\[
 \max_{C,\rho\ {
m corners}}\left(A+\frac{f^2}{4q_{\rm s}}\right),
 \qquad
 \max_{C,\rho\ {
m corners}}
       \left(A+\frac{f^2}{4(q_{\rm s}+e_\rho)}-K\rho\right).
\]
Any threshold-root export likewise uses the corrected corner curvature.

For completeness, nonnegativity in the later acquired domains can be checked
without replaying their iterations. Write $B=\mathbf1_{\{X>0\}}$,
$V=\mathbb E(X\mid B)$,
$W=X-V$, $\theta=\mathbb E W^2$, and $p=\mathbb P(X>0)$. Conditional
centering and $W\ge-R-V$ imply
\[
 \gamma=\mathbb EV^3+3\mathbb E(VW^2)+\mathbb EW^3
 \ge\mathbb EV^3+2\mathbb E(VW^2)-R\theta.
\]
For $p\in[l,u]\subset(0,1/2)$ and $\theta\le\Theta<1$ this gives
\begin{equation}\label{re:actual-gamma-license}
 \gamma\ge
 \frac{(1-\Theta)^{3/2}(1-2u)}{\sqrt{u(1-u)}}
 -\left(R+2\sqrt{\frac{u}{1-u}}\right)\Theta.
\end{equation}
Here the binary cubic decreases in $p<1/2$, and the right side decreases
as its variance bound increases. These are bounds on the same actual law
throughout every later call. The fixed acquired input domains license
\eqref{re:actual-gamma-license} before those calls are made.
The exact domain check covers 302 input rows: 45 central-input rows,
45 left-input rows, 45 first-confinement feedback rows, 45 Sperner feedback
rows, 84 signed-resource rows, and 38 final left-radius rows. With the
common already acquired radius $R=3.84$, each right side is greater than
$.096$. Thus all later paired calls have an actual nonnegative signed
third moment, independently of any auxiliary Fourier orientation.

There is also a convenient cell-constant correction. Since $\rho\ge1$,
$(-g_0(\rho))_+\le(d-.16)_+$, so the weaker price
\[
 C\rho+\tfrac12L g_0(\rho)
       -\tfrac12(U-L)(d-.16)_+
\]
is valid. This is the correction selected for the additional index--46
calculation. Both forms keep the raw remainder separate from the
barrier-subtracted resource budget.

The broad preliminary and complementary resource passes are replayed with
\eqref{re:sign-aware-price}; their subsequent consumers retain the same
fixed target tables only when the complete corrected replay verifies every
exported resource. A diagnostic equality of two corner maxima alone does
not replace this check, since threshold bounds also consume the phase
coefficients. Later uses of the simpler formula require their own explicit
nonnegative signed-moment or small-defect license.


\subsection{The unsigned first-stage cover is the same enclosure}
\label{re:unsigned}
The first finite barrier uses a cover of $\rho$ alone before introducing
the signed cover. Here are its exact coefficient choices, so that the
first barrier needs no separate inherited entry theorem. Fix a closed
cell $\rho\in[l,u]\subset[1,1.85]$, actual radius $1\le R\le6$, and
write $k_j(R)$ for the radial ratios of Lemma~\ref{re:pricing}. Set
\begin{gather*}
 M=\min\{11.1,Ru,1+k_4(R)(u-1)\},\quad
 g=\min\{1.85,u,\sqrt{M-1}\},\\
 E_5=\min\{RM,1+k_5(R)(u-1)\},\qquad
 b=\sqrt{[u\sqrt{u^2+8}-u^2-2]/2}.
\end{gather*}
The last expression solves $\rho\ge(1+\beta^2)/\sqrt{1-\beta^2}$,
proved by $\rho m\ge a_+^2+a_-^2$ and $m^2\le a_+a_-$.
If a previous pass licenses $D\le d<2$, the generic-event lemma also
permits the ceilings $h_*u+Rd/2$ and $h_*^2u+R^2d/2$ for $M,E_5$;
all dependent quantities are recomputed after taking these minima.
Use the triangle LOW triple and the minimum of complete permitted HIGH
values, with the same moment cap. Put
\[
 K_F=\tfrac{16}{25}(L_0+K_H+L_2/\sqrt N),\qquad s_F=\tfrac{16}{25}L_H.
\]
The radial identity bounds $|K_1|$ by
\[
 k_{\rm rad}=3ub+
     \frac{R(R+2-3l/2)}{R+1/2}(u-1),
\]
provided the displayed radial coefficient is at least $31/40$.
Indeed its positive maximum is at $R$ and its negative part is at most
$31/40$ on $[0,R]$ for $1\le\rho\le1.85$; convexity of
$x(x+2-3\rho/2)/(x+1/2)$ proves this directly.
Two additional valid bounds are
\[
 k_{\rm old}=6u-3u^2/2+3b,\qquad
 k_R=Ru/2+\max_{t=l,u}|R/2-3t/2|g+3b.
\]
For $k_R$, bound the positive and negative fourth moments separately by
$R(\rho\pm\gamma)/2$. Before defect feedback use
$k=\min(k_{\rm rad},k_{\rm old},k_R)$.
After feedback one may replace $k_{\rm old}$ by $\max(k_+,k_-)$, where
\begin{gather*}
 g_l=3.16l-3-d,\quad g_u=3.16u-3-d,\\
 k_+=M+3b-\tfrac32\min_{r=l,u;,v=g_l,g_u}rv,\\
 k_-=3u+3b+\max_{r=l,u;,v=g_l,g_u}(3r/2-3)v.
\end{gather*}
These substitutions bound the lower affine graph for the actual $\gamma$,
not an independently chosen skewness.

Put $\Delta=2u$ initially and $\Delta=(1-3.16)l+3+d$ after feedback.
The constant coefficient of the raw identity is bounded by
\begin{align*}
 A_* =\min\bigl\{&\tfrac32b(2u+1)+\tfrac32\Gamma(R,l)\Delta
       +\tfrac38u^2(u+1)+\tfrac8{27}u^3-l/2,\\
      &\tfrac32b(2u+1)+2uM-l/2\bigr\}.
\end{align*}
This is the coarse gap bound above together with the two bounds for
$-K_q$. The density coefficient obeys
\[
 v_H=\min\{.4j+.4\max(2,M-3)/6,
                     .4\max(1/2,2M/3-1/2)\},
\]
where $j=3(u+1)$ initially; the generic-gap choice is
$j=\alpha(u+1)+(3-\alpha)d/2$, with its checked $\alpha<3$.
For a score endpoint enclosure $C\in[c_-,c_+]$, the resulting raw
bound is
\begin{equation}\label{re:unsignedbound}
 K_F+c_+A_*+\frac{(c_+k+v_H+s_F)^2}{4q_*},
 \qquad q_*>0,\quad k\ge0,\quad A_*>0,
\end{equation}
where initially $q_*=c_-l-g/5$. After feedback use the larger of that
quantity and
\[
 c_-l+\varphi_-(3.16l-3-d)/2
                  -(2/5-\varphi_-)(d-.16)_+/2.
\]
The score endpoints are the lower score floor and the upper allowed
score cap; if the latter depends on $\rho$, use its values at $u,l$
for $c_-,c_+$, respectively. Convexity in the actual score precedes the
completion of the square, and these outward endpoint intervals cover it.
For the phase resource add \eqref{re:extra} before maximizing. A barrier
$e\ge K/n$ subtracts $K l$ from the nonnegative resource. The first raw
target must pass the license of Lemma~\ref{re:weight} first.

\paragraph{Optional signed improvement inside an unsigned cell.}
If $g_0=\max(-g,\sigma_l l-3-d)>g$, exclude the cell. When $g_0\ge0$,
the sign-moment constraint forces
$(1-\beta)^3\le(1+\beta)(u-g_0)^2$.
This inequality is monotone on $(-1,1)$, so exact bisection yields a
lower $b_0\le\beta$; if it fails even at $b$, the cell is empty.
For $x\in[g_0,g]$, set
\begin{align*}
 f_-(x)&=(l+x)^2/[2(1+b)]-R(u-x)/2-3b-3ux/2,\\
 f_+(x)&=R(u+x)/2-(l-x)_+^2/[2(1-b_0)]-3b_0-3lx/2.
\end{align*}
The signwise fourth-moment Cauchy bounds prove
$\min f_-\le K_1\le\max f_+$. The first minimum is at its clipped
vertex $3u(1+b)/2-l-R(1+b)/2$. The second maximum is among its endpoints,
the clipped kink $l$, and the clipped vertex
$l+(R/2-3l/2)(1-b_0)$. Intersect with $[-k,k]$ and replace the linear
term in \eqref{re:unsignedbound} by the smaller full rectangle bound
\[
 \max_{c=c_-,c_+;,k=k_l,k_u;,\psi=\varphi_-,.4;,
             v=-1/2,\,2M/3-1/2}|ck+\psi v|+s_F.
\]
All sixteen corners are retained. If $b_0>0$, the opposite-sign event
also has probability at least $(1+b_0^2)^2/(2u^2)$, from
$m\ge(1+\beta^2)/(2\rho)$ and $m^2\le pq$. It may improve the complete
HIGH value. These are intersections with necessary conditions; the cover
still examines every $\rho$-cell.

\paragraph{Moving from a small defect to probability coordinates.}
For $0<d\le.7$, start with
\[
 A_l=\frac{1.47\sqrt{2-d}-\sqrt{3.7d}}{3.52},\qquad
 B_l=\frac{2.04\sqrt{2-d}-\sqrt{3.7d}}{3.52}.
\]
They follow from $A+B\ge\sqrt{2-d}$ and
$|\sqrt{\sigma+1}A-\sqrt{\sigma-1}B|\le\sqrt{3.7d}$;
their positivity is checked, rather than assumed at every $d$.
The pointwise inequality
$x(x+b)\ge(29/5)b(x-b)$ for $x>b>0$, applied to the two residual
third moments, gives
$D\ge(29/5)\min((\sigma+1)A,(\sigma-1)B)\theta$.
Start $\Theta=1$ and perform the declared finite number of updates:
\begin{align*}
 \Theta&\leftarrow\min\{\Theta,d/[(29/5)
                 \min((104/25)A_l,(54/25)B_l)]\},\\
 \epsilon&=\sqrt{d/[(616/25)(1-\Theta)^2]},\qquad
 [l,u]=[(4-\sigma_u)/2-\epsilon,(4-\sigma_l)/2+\epsilon],\\
 A_l&\leftarrow\max\{A_l,\sqrt{(1-\Theta)l/(1-l)}
                           -d/[2(1-u)(179/25)A_l^2]\},\\
 B_l&\leftarrow\max\{B_l,\sqrt{(1-\Theta)(1-u)/u}
                           -d/[2l(129/25)B_l^2]\}.
\end{align*}
The last two bounds use
$D\ge2(\sigma+2)B^2M_++2(\sigma+4)A^2M_-$ and
$\Delta_A=M_-/q$, $\Delta_B=M_+/p$.
Every division, interval endpoint and variance domain is gated.
Subdivide this entire probability interval, then use the common
cell update and coefficient atlas. Thus the unsigned first barrier,
the signed initialization and the final probability refinements are
three data choices for the same finite-enclosure theorem.

\section{One confinement lemma with two concentration choices}\label{conf:re:confinement}
A proved amplitude is used as a support update, never as a premise of
its own proof. The following derivation applies to every acquired row,
with the same formulas for the central and left probability ranges.
\subsection{Statement and hypotheses}
Let $n\ge N=47$, $Y=B+U$, $\Pp(B=1)=p<1/2$ and $\E(U\mid B)=0$.
Put $q=1-p$, $s=\E U^2$, $\lambda=ns$ and $X=\ell(Y-p)$, where
$\E X^2=1$. Thus $\ell^2=(1-\theta)/(pq)$ and $\theta=\ell^2s$.
The conditional label-zero support is at most $a\ge0$ and the label-one
support is at least $b\le1$, where $a,b$ are support points and
$b-a=1-d$. In particular $0\le a\le d$, $1-d\le b\le1$.
Assume the initial strip $\operatorname{supp}Y\subset(-.55,1.55)$.
Write $\rho=\E|X|^3$, $\beta=\E[X|X|]$ and $\epsilon=\E|U|$.

Select an acquired closed probability panel $[l,h]$.
Its fields give upper bounds $\Theta,d_0,\rho_0,\beta_+,T,\Lambda$,
a lower bound $\beta_-$, and separate bounds $H_+,H_-$ for the two signs of
$H$. Assume
\[
 \theta\le\Theta,\quad d\le d_0,\quad \rho\le\rho_0,
 \quad \beta_-\le\beta\le\beta_+,\quad 0\le\lambda\le\Lambda,
 \quad .4<C\le C_0=.4253.
\]
For the positive sign $0\le H\le H_+$; for the negative sign
$-H_-\le H\le0$. Let $x=np+H/\ell$ and
$\nu=\mathcal L(x-\sum_{i<n}Y_i)$. Assume a contact majorant $Q$ for
$g(y)=\nu([y,\infty))$, with equality at each support point of $Y$,
and derivative
\begin{equation}\label{conf:derivative}
 \sqrt n(-Q'(y))=\ell\phi(H/\sqrt n)
 +\frac\ell n\{3C\beta+(H\phi(H/\sqrt n)+3C\rho)t-3Ct|t|\},
 \quad t=\ell(y-p).
\end{equation}
Both actual deleted convolutions satisfy, separately,
\begin{equation}\label{conf:concentration}
 Q_j(w)\le\frac{\ell\phi(0)w+2C_0\rho}{\sqrt j},
 \qquad j=n-1,n-2,
\end{equation}
for every interval of width $w\ge0$.

Here is the precise common resource assumption. Put
$\kappa=997/15000<\phi(0)/6$, $J=\int_{[a,b)}(y-a)\nu(dy)$,
and note that the central gap contributes $m\ell J=(1-\theta)J$ to $L$.
Thus $L\ge(1-\theta)J$; retain $e\ge0$ and
\begin{equation}\label{conf:budget}
 n e\rho+n\phi(0)D/6+n^{3/2}L<\mathcal B(H),\qquad
 \mathcal B(H)=\min\{T,\max_j(A_j+F_{j,\pm}|H|-Q_jH^2)\}.
\end{equation}
The four quadratics are the table's \texttt{quadratics} rows, using
\texttt{Qbudget}, not \texttt{Qraw}. The Gaussian phase contribution
has already been included in \texttt{Qbudget}.

For each panel let $E,F_\ell,m_0$ be respectively strict lower, upper,
and lower rational bounds for
\[
 \sqrt{(1-\Theta)/(h(1-h))},\qquad
 1/\sqrt{l(1-l)},\qquad \sqrt{(1-\Theta)l(1-l)}.
\]
The bounds use the grid $10^{-6}$, with an extra unit outward. Define
\begin{align*}
 c_+&=2(5.16)E^2(1-h-d_0)^2,&
 c_-&=2(7.16)E^2(l-d_0)^2,\\
 b_*&=\min\{3(4.16)E(1-h-d_0),3(6.16)E(l-d_0)\},\\
 P&=m_0E^2\pos{3.5\operatorname{dist}([l,h],[.41886,.41887])-2.041d_0}^2.
\end{align*}
The simultaneous moment inequality has the following usable form:
\begin{equation}\label{conf:moment}
 D\ge c_+M_++c_-M_-+b_*\theta+P,\quad
 d=(M_+/p+M_-/q)/\ell,\quad
 \epsilon\le2(M_++M_-)/\ell,\quad M_\pm\ge0.
\end{equation}
To derive it, the endpoint expansion of \eqref{re:gapexact} gives
$D\ge m(\sqrt{\sigma+1}A-\sqrt{\sigma-1}B)^2
+2(\sigma+2)B^2M_++2(\sigma+4)A^2M_-
+3(\sigma+1)B\theta_++3(\sigma+3)A\theta_-$.
Insert $A\ge E(l-d_0)$, $B\ge E(1-h-d_0)$ and the rational lower
coefficients displayed above. For the square, write
$A=\ell(p-a)$, $B=\ell(1-p-(1-b))$, with
$a+(1-b)=d\le d_0$. Its unshifted coefficient is
$(\sqrt{\sigma+1}+\sqrt{\sigma-1})(p-p_*)$;
$\sqrt{\sigma+1}+\sqrt{\sigma-1}>3.5$ and
$\sqrt{\sigma+1}<2.041$ prove the stated lower square $P$.
For completeness, on the positive class put $Z=X-B\ge0$.
Then $T_+=2B^2M_++3B\E[Z^2;X>0]+\E[Z^3;X>0]$ and
$R_+=BM_++\E[Z^2;X>0]$. Consequently
$(\sigma-1)T_++6BR_+\ge2(\sigma+2)B^2M_+
+3(\sigma+1)B\theta_+$, since $\E[Z^2;X>0]\ge\theta_+$.
The reflected expansion gives the negative-class bound.
The centered absolute deviation on each class is at most twice its
mean distance from the inner endpoint, proving the bound on $\epsilon$.
All quantities refer to the same marginal.

\subsection{Concentration from an antichain}
\begin{lemma}[Strict-gap concentration]
Suppose $Y=U_0+B(1+U_1-U_0)$, where $B$ is Bernoulli$(p)$ independent
of the pair $(U_0,U_1)$, and $1+U_1-U_0\ge c>0$ almost surely.
For independent copies $Y_1,\ldots,Y_j$ and any interval $I$ of diameter
strictly less than $c$,
\[
 \Pp\left(\sum_{i=1}^jY_i\in I\right)
 \le \max_{0\le k\le j}\binom jk p^k(1-p)^{j-k}.
\]
The same bound holds after adding any independent random variable.
\end{lemma}
\begin{proof}
Condition on all noise pairs. The sum is a fixed constant plus
$\sum B_iw_i$, with every $w_i\ge c$. The family $\mathcal F$ of subsets
whose weighted sums lie in $I$ is an antichain: distinct comparable
subsets have sums differing by at least $c$, larger than the diameter.
A uniformly random permutation of $\{1,\ldots,j\}$ defines the maximal
chain of its initial subsets. The chain contains a fixed $k$-subset with
probability $1/\binom jk$, and meets an antichain at most once. Hence
\[
 \sum_{S\in\mathcal F}\binom j{|S|}^{-1}\le1.
\]
Multiplying each summand by its binomial level mass proves the bound
conditionally. Integration removes the conditioning. Conditioning on an
additional independent variable proves the last assertion. This argument
includes all atom collisions and any open or closed interval endpoints.
\end{proof}
For the actual marginal we may couple the two conditional noises
independently and then sample an independent label. Since $U_0\le a$ and
$U_1\ge b-1$, the step is at least $b-a=1-d$.
Strict diameter is necessary: a Bernoulli law on $\{0,c\}$ assigns mass
one to $[0,c]$, exceeding its largest atom when $0<p<1$.

\paragraph{A rational cap for every remaining index.}
Put $J_0=N-2=45$ and define
\[
 r_L=l-(1-l)/J_0,\qquad r_R=h+h/J_0,\qquad
 v_* =\min\{r_L(1-r_L),r_R(1-r_R)\}.
\]
The checker verifies $0<r_L\le r_R<1$. Adjacent binomial-mass ratios
show that any mode $k$ at index $j\ge J_0$ satisfies
$(j+1)p-1\le k\le(j+1)p$, hence $r=k/j\in[r_L,r_R]$.
In particular $1\le k\le j-1$. Finite Stirling bounds give
\[
 \sqrt j\binom jk p^k(1-p)^{j-k}
 \le\frac{\exp(1/(12j)-jD(r\Vert p))}{\sqrt{2\pi r(1-r)}}
 <\frac{.4}{\sqrt{v_*}(1-1/(12J_0))}.
\]
Here $D(r\Vert p)\ge0$ follows from Jensen's inequality for the logarithm;
$\exp x\le(1-x)^{-1}$ for $0\le x<1$ follows term by term from the
power series. The finite Stirling inequalities used are
$\sqrt{2\pi m}(m/e)^m<m!<\sqrt{2\pi m}(m/e)^m\exp(1/(12m))$.
Choose the strict $10^{-6}$ grid upper square root $M$ of
$.16/[v_*(1-1/(12J_0))^2]$. Thus, whenever $w<1-d$,
\begin{equation}\label{conf:antichainconcentration}
 Q_j(w)<M/\sqrt j,\qquad j\ge n-2\ge J_0.
\end{equation}
This covers the $n-1$ full summands of $\nu$ and the $n-2$ full summands
remaining after one mark or deletion. An independent remaining conditional
noise preserves the cap. We do not replace the latter index by $n-1$.


Either the score-derived cap or the strict-gap cap may be selected on
an interval where its hypotheses hold. The first confinement uses the
score cap; later central and left schedules use the strict-gap cap.
For the left panels replace the conditional cubic coefficient 24 below
by 28; in general any $c_\rho$ satisfying
$c_\rho^2[l(1-l)]^3>[3(1-l)+.55]^2$ is valid.
\subsection{Signed deletion, including atoms}
Let $m=n-1$, $Z=x-\sum_{i=1}^mY_i$, $\mu=x-mp$ and let $K=\sum B_i$.
For a fixed removed label put
$T_b(I)=\Pp(x-\sum_{i<m}Y_i-b-U_b\in I)$.
Exchangeability and $mp-K=Z-\mu+\sum U_i$ give the exact identities
\[
 \nu(I)-T_1(I)=\frac{\E[(Z-\mu+\sum U_i)1_I(Z)]}{mp},\qquad
 T_0(I)-\nu(I)=\frac{\E[(Z-\mu+\sum U_i)1_I(Z)]}{mq}.
\]
There is no absolute value around the sum of noises.
For a full removed pair $(Y_i,U_i)$, condition on the pair. The remaining
sum gives a function $f(Y_i)\in[0,Q_{m-1}(w)]$. Since $\E U_i=0$,
\[
 \E[U_i f(Y_i)]\le Q_{m-1}(w)\E(U_i)_+
 =\tfrac12\epsilon Q_{m-1}(w).
\]
This also works when one conditional law is deterministic.

Let independent conditional noises satisfy $\Delta=U_1-U_0\ge-d$.
For $\tau>0$ put $I^+=I+1+[-d,\tau]$. On $\Delta\le\tau$, the
point for $T_1$ in $I$ maps to the point for $T_0$ in $I^+$. On the
complement, conditioning on both noises gives error at most
$Q_{m-1}(w)\Pp(\Delta>\tau)$. Moreover
\begin{equation}\label{conf:bad}
 \Pp(\Delta>\tau)\le
 \min\left\{\frac{s}{s+p\tau^2},\frac d{d+\tau}\right\}.
\end{equation}
Indeed $\E\Delta=0$, $\operatorname{Var}\Delta\le s/p$, and
$\Delta+d\ge0$ has mean $d$. Cantelli and Markov prove the two bounds.
The exceptional event is strict, so $\Delta=\tau$ belongs to the good event.

If $I,I^+$ are disjoint, combine the two signed noise terms before
bounding them. For the same remaining-sum law their conditional
probabilities $f_I,f_J$ satisfy
\[
 0\le f_I\le Q_I,\quad0\le f_J\le Q_J,\quad f_I+f_J\le Q_{\rm hull}.
\]
Since $p<q$,
\[
 \frac{f_I}{p}+\frac{f_J}{q}
 \le\min\left\{\frac{Q_I}p+\frac{Q_J}q,
 \frac{Q_{\rm hull}}q+\left(\frac1p-\frac1q\right)Q_I\right\}.
\]
Multiplication by $\E U_+=\epsilon/2$ proves the combined bound.
Closed, open and half-open intervals are all covered by containment in
the corresponding closed interval. No density or atom separation is used.

\subsection{Slope bounds and the two exterior intervals}
Split $|H|$ into 16 closed equal cells $[h_l,h_h]$ and $\lambda$ into
24 closed equal cells $[s_l,s_h]$. On each box put
\[
 \begin{gathered}
 \rho_* =\min\{\rho_0,\rho_B(l)+24s_h/N\},\\
 b_- =\max\{\beta_-,1-2h-2(1-h)\Theta\},\qquad
 b_+ =\min\{\beta_+,1-2l+\Theta\}.
 \end{gathered}
\]
Here $\rho_B(p)=(1-2p(1-p))/\sqrt{p(1-p)}$. Write $s_i=\E U_i^2$. Conditional centering and the two sign supports give
\[
 \rho=\ell^3\{pq(1-2pq)+3pq(s_0+s_1)
                  +p\E U_1^3-q\E U_0^3\}.
\]
The initial strip and $d_0<.55$ imply $|U_i|<.55$.
Thus $\rho\le\rho_B+\ell^3(3q+.55)s\le\rho_B+24s$; the last
coefficient inequality is checked as
$24^2[l(1-l)]^3>[3(1-l)+.55]^2$.
Similarly $\beta=(1-\theta)(1-2p)+\ell^2(ps_1-qs_0)$, and
$|\ell^2(ps_1-qs_0)|\le\theta$. These identities prove the stated
outward bounds on $\rho$ and $\beta$.
Set $f=.4F_\ell$, $\varphi_-=.39886(1-h_h^2/(2N))>0$.
The exact upper bound for $H\phi(H/\sqrt n)$ used below is
\[
 v_+=\begin{cases}.4h_h&H\ge0,\\-\varphi_-h_l&H\le0.\end{cases}
\]
For all $t\le0$, completing the square in \eqref{conf:derivative} gives
\[
 \mathcal M_-=F_\ell\pos{(v_++3C_0\rho_*)^2/(12(.4))-3(.4)b_-},
 \qquad L_-=E\varphi_- -\mathcal M_-/N.
\]
A positive bound $B_*>\sqrt{(1-l)/l}$ gives valid central-gap lower
and positive-half-line upper slope bounds
\begin{align*}
 \mathcal M_+&=F_\ell\pos{.4h_hB_*+3C_0B_*(B_*-1)-1.2b_-},\\
 L_g&=E\varphi_--\max(\mathcal M_-,\mathcal M_+)/N,\\
 U_+&=f+\frac{F_\ell}{N}
 \left\{3C_0b_++\frac{(.4h_h+3C_0\rho_*)^2}{12(.4)}\right\}.
\end{align*}
They follow by completing a square on the positive half-line and using
$t(t-1)\le B_*(B_*-1)$ for $1\le t\le B_*$; on $0\le t\le1$ that
term has the favorable sign. Negative values are covered by $\mathcal M_-$.
All divisions and positivity requirements are explicit gates.

\paragraph{The sharper lower-ray slope.}
Take $\tau_-=5r_0/12$, $\eta=10^{-5}$ and
$I=[a-r,a-d-\tau_--\eta]$, where $r\ge r_0$.
Since $a\le d$, the entire interval has
$t\le t_*=-E(l+\tau_-+\eta)<0$.
Write $B(t)=b_-+\rho_*t+t^2$ and
$B_*^{\min}=B(\min(t_*,-\rho_*/2))$. For the actual $C$ and $t\le t_*$,
\[
 3C(\beta+\rho t+t^2)+H\phi t
 \ge1.2B(t)+v_+t+3(C_0-.4)\min(0,B_*^{\min}).
\]
The inequality uses $t<0$ and separates the possible negative part of
$B(t)$ before varying $C$. Minimize the right side at the clipped vertex
$\min(t_*,-(v_++1.2\rho_*)/2.4)$ and call its value $q_*$. Then
\[
 L=E\varphi_--F_\ell\max(0,-q_*)/N
\]
is a lower slope on the whole interval, for every $n\ge N$.
Positive corrections are discarded. Therefore contact gives
\[
 \sqrt n\nu(I)\ge L(r-d-\tau_--\eta).
\]
This one slope replaces the previous separate bounds on the long and
short parts of the interval.

\paragraph{Lower interval geometry.}
Its translate is $I^+=[a-r+1-d,b-\eta]$, of width $r-\eta$.
The initial strip implies $I^+\subset[a+g,b)$,
$g=1-.55-2d_0>0$, hence $\nu(I^+)\le J/g$.
The intervals are disjoint throughout $r\le.55+d_0$, since
$\inf I^+-\sup I=1-r+\tau_-+\eta>0$.
At $r=r_0$, concentration and directed distance bounds are
\begin{align*}
 Q_I&=f(r_0-\tau_--\eta)+2C_0\rho_*,&
 Q_J&=f(r_0-\eta)+2C_0\rho_*,\\
 Q_{\rm hull}&=f(1+r_0-\eta)+2C_0\rho_*,\\
 D_I&=\pos{v_H-l-\tau_--\eta},&
 D_J&=\pos{v_H+1-l-\eta},
\end{align*}
where $v_H=-h_l/F_\ell$ for $H\ge0$ and $v_H=h_h/E$ for $H\le0$.
These signs follow from $\mu=p+H/\ell$.

\paragraph{Upper interval geometry.}
For a support point $y\ge1+r_0$, take $\tau_+=7r_0/15$,
$I=[a,a+\tau_+]$ and $I^+=[b,a+1+2\tau_+]$.
The strict gates $d_0+2\tau_+<r_0$ and $1-d_0-\tau_+>0$ give
$I^+\subset[b,y)$ and disjointness. Contact, including the atom at $b$,
and the leakage definition give
\[
 \sqrt n\nu(I)\ge L_g(1-d)-\sqrt nJ/\tau_+,\qquad
 \sqrt n\nu(I^+)\le U_+(.55+d).
\]
Use
\begin{align*}
 Q_I&=f\tau_++2C_0\rho_*,&Q_J&=f(2\tau_++d_0)+2C_0\rho_*,\\
 Q_{\rm hull}&=f(1+2\tau_+)+2C_0\rho_*,\\
 D_I&=\pos{d_0+\tau_++h_h/E-l},&
 D_J&=\pos{1+d_0+2\tau_++h_h/E-l}.
\end{align*}
The upper side deliberately keeps the absolute phase bound.


\paragraph{The strict-gap alternative.}
If $1-.55-2d_0>0$ and the displayed lower intervals have width below
$1-d$ throughout $r\in[r_0,.55+d_0]$, take $Q_I=Q_J=M$.
For the upper intervals the corresponding gate is
$1-2d_0-2\tau_+>0$; then again $Q_I=Q_J=M$.
Use the same directed distances. Their concentration caps are now
independent of $r$, and use
$a_e=F_dM\{1/l+1/(1-h)\}/2$ in the next formulas.
The lower-ray gate reduces to $L>0$: its contact side increases with
$r$, while the entire error is constant on the acquired finite strip.
For the score-derived cap retain the explicit affine slope gate below.
No assertion outside that strip is needed.

\subsection{One budget, two error bounds}
Let $F_d$ be the strict grid upper bound for
$\sqrt{\max\{(47/46)^3,47/45\}}$. It simultaneously covers the two
concentration factors in \eqref{conf:concentration} after multiplication by
$\sqrt n$, uniformly for $n\ge47$. For either pair of intervals set
\begin{align*}
 a_e&=\frac{F_d}{2}\min\left\{\frac{Q_I}l+\frac{Q_J}{1-h},
 \frac{Q_{\rm hull}}{1-h}+\left(\frac1l-\frac1{1-l}\right)Q_I\right\},\\
 A&=F_d\left\{Q_ID_I/l+Q_JD_J/(1-h)\right\},\qquad B=F_dQ_I.
\end{align*}
The probability denominators are outward because $1/p-1/(1-p)$ is
decreasing on $(0,1/2)$.
Use $a_d=L$, $a_J=1/g$ on the lower side, and
$a_d=L_g+U_+$, $a_J=1/\tau_+$ on the upper side.
For the phase cell, maximize each concave budget quadratic at its
vertex clipped to $[h_l,h_h]$, take their maximum, then the minimum
with $T$; call this upper bound $T_*$. Put
\[
 R=T_*/\kappa-b_*E^2s_l,
 \quad\alpha=\min(c_+lE,c_-(1-h)E),\quad G=(R-NP)/\alpha.
\]
If $G\le0$, the strict common budget excludes the entire closed box.
Otherwise $nd<G$ and
\[
 B_\tau=\min\{s_h/(s_h+Nl\tau^2),G/(N\tau+G)\}
\]
bounds \eqref{conf:bad}. With
\[
 z=\max\left\{\frac{a_d/l+2a_e}{Ec_+},
 \frac{a_d/(1-h)+2a_e}{Ec_-},\frac{\kappa a_J}{1-\Theta}\right\},
\]
one valid total error bound is
\begin{equation}\label{conf:oldbound}
 \mathcal E_0=(zR+A)/N-zP+B B_\tau.
\end{equation}
This follows by multiplying the single budget by $z/\kappa$.
The inverse-index coefficient $zR+A$ is checked positive.

A stronger bound spends the same resource on the pair event. Let
$R_0=R-NP>0$ and define the three pairs
\[
 (v_+,w_+)=\left(\frac{a_d/l+2a_e}{Ec_+}\frac{R_0}{N},
                    \frac{R_0}{Ec_+l}\right),
\]
\[
 (v_-,w_-)=\left(\frac{a_d/(1-h)+2a_e}{Ec_-}\frac{R_0}{N},
                    \frac{R_0}{Ec_-(1-h)}\right),\qquad
 (v_J,w_J)=\left(\frac{\kappa a_J}{1-\Theta}\frac{R_0}{N},0\right).
\]
They are upper cost/displacement coordinates at the vertices of the
simplex with resources $nc_+M_+$, $nc_-M_-$ and $n^{3/2}L/\kappa$.
For any $t\ge0$,
\[
 \frac{D}{N\tau+D}\le\frac{N\tau}{(N\tau+t)^2}D+
 \frac{t^2}{(N\tau+t)^2}\qquad(D\ge0).
\]
The difference is $N\tau(D-t)^2/[(N\tau+D)(N\tau+t)^2]$.
Consequently a further valid error bound is
\[
 \mathcal E_t=\frac AN+
 \max_{i\in\{+,-,J\}}\left(v_i+B\frac{N\tau}{(N\tau+t)^2}w_i\right)
 +B\frac{t^2}{(N\tau+t)^2}.
\]
The checker takes the minimum of \eqref{conf:oldbound} and these eleven
bounds with $t=jG/10$, $j=0,\ldots,10$. Each is individually valid;
no numerical minimization or derivative of their minimum is assumed.

The lower contradiction is the strict gate
$Lr_0-L(\tau_-+\eta)>\mathcal E$.
To extend it to every $r\ge r_0$, bound increments directly in the original
affine-in-width signed deletion bound for the same law. The budget gives
$n\epsilon<E_*=2T/(\kappa E\min(c_+,c_-))$. The sufficient slope gate is
\[
 L>\frac{F_df}{N}\left\{D_I/l+D_J/(1-h)
 +\frac{E_*}{2}\left(1/l+1/(1-h)\right)\right\}+F_dfB_\tau.
\]
Using the uncoupled increment is safe even when the hull bound was used
at $r_0$. The upper contradiction is $L_g-.55U_+>\mathcal E$.
Thus no support point lies below $a-r_0$ or above $1+r_0$.
Since $a\ge0$, $b\ge1-d_0$ and $d_0<r_0$, the label-zero and label-one
supports both satisfy $|U|<r_0$.


\paragraph{The selected schedules.}
The ordinary, central-antichain, and left-antichain certificates substitute
their recorded $r_0$ into these same gates. Noise-dependent amplitudes
are the same theorem with a positive lower noise endpoint in
$R=T_*/\kappa-b_*E^2s_l$. In the central rows the additional amplitudes
are $.35,.30,.275$ on $[.05,.08],[.08,.10],[.10,\Lambda]$;
on the left they are $.375,.325,.30$. Intersect each interval with its
own row's $\Lambda$ and each amplitude with that row's global amplitude.
The fixed acceptance checks both signs of $H$, all closed phase/noise
cells, clipped endpoints, and strict budget exclusions. No proof of
LOW, MID, or a binomial score comparison is used by this confinement lemma.


% Source: local_comparison.tex
\section{One local comparison theorem}\label{shortloc:section}

A law concentrated near two points has a natural integer count and a small
centered error. We compare it with its own Bernoulli count. The same exact
identity handles every noise regime. There are two useful estimates of its
right side: discard all but the nearest outgoing count, or retain all counts
through a random cutoff. The estimates share their normalization, moments,
and tail bounds; only the final scalar witness changes.

\subsection{The identity and the common numerical test}

Let $0<p\le1/2$, $q=1-p$, $v=pq$, and let
\[
 Y=B+U,\qquad \Pr(B=1)=p,\qquad \mathbb E(U\mid B)=0.
\]
The two actual conditional laws $U_0,U_1$ obey
\begin{equation}\label{shortloc:support}
 |U_j|\le A,\quad U_0\le a_0,\quad U_1\ge-a_1,\quad
 a_0,a_1\ge0,\quad d=a_0+a_1\le d_0<\min(p,q),\quad d_0\le A.
\end{equation}
Write $s_j=\mathbb E U_j^2$, $s=qs_0+ps_1$, $\lambda=ns$,
$\ell=(v+s)^{-1/2}$, $\theta=s/(v+s)$ and
$\rho=\mathbb E|\ell(Y-p)|^3$. At the actual threshold choose
\[
 x=np+H/\ell=k+u,\qquad k\in\{1,\ldots,n-1\},\quad |u|\le\tfrac12,
 \qquad \delta=k-np=H/\ell-u.
\]
Either nearest integer is allowed at a tie. Put
\[
 \pi_j=\binom njp^jq^{n-j},\quad F_B(j)=\Pr(\operatorname{Bin}(n,p)\le j),
 \qquad \rho_B=(1-2v)/\sqrt v,
\]
with $\pi_j=0$ outside $[0,n]$. Let $W_j$ be an independent sum of $j$
copies of $U_1$ and $n-j$ copies of $U_0$. The signed scores are
\[
 C=\frac{\sqrt n}{\rho}\{\Pr(\textstyle\sum Y_i\le x)-\Phi(H/\sqrt n)\},
 \qquad C_B=\frac{\sqrt n}{\rho_B}\{F_B(k)-\Phi(\delta/\sqrt{nv})\}.
\]
Throughout this section $C\ge c_*>0$; the selected application uses $c_*=.4$.
Define
\begin{align}
 f(u)&=\frac{\Phi((\delta+u)/\sqrt{n(v+s)})-
                    \Phi(\delta/\sqrt{n(v+s)})}{\pi_k},\label{shortloc:f}\\
 Q&=\Pr(\textstyle\sum Y_i\le k+u)-F_B(k),\nonumber\\
 \mathcal R&=\sqrt n\{\Phi(\delta/\sqrt{nv})-
                 \Phi(\delta/\sqrt{n(v+s)})\}+C(\rho_B-\rho).\nonumber
\end{align}
Direct subtraction gives the exact identity
\begin{equation}\label{shortloc:master}
 \rho_B(C-C_B)=\sqrt n\pi_k\{Q/\pi_k-f(u)\}+\mathcal R.
\end{equation}
In particular, the phase $u$, the sign of $H$, and both conditional laws
refer to the same marginal throughout.

\begin{samepage}
\begin{theorem}[Parameterized local comparison]\label{shortloc:theorem}
Fix $n\ge N$, $\lambda>0$ and a closed parameter cell on which
\eqref{shortloc:support} holds. Suppose $\sqrt n\pi_k\ge m>0$.
Either of the following witnesses proves $\rho_B(C-C_B)<-\varepsilon\lambda$.
\begin{enumerate}
\item \emph{Nearest-count witness.} Valid bounds $O,T,R_0,\kappa$ satisfy
\[
 \Pr(W_k>u)+f(u)\ge O,\quad
 \sum_{i\ge1}\Pr(W_{k+i}\le u-i)\le T,\quad
 \pi_{k+i}/\pi_k\le R_0,\quad \mathcal R\le\kappa\lambda/n.
\]
Only legal counts occur in the sum; adding positive upper bounds for
other counts is allowed. If $D=O-R_0T>0$ and
\begin{equation}\label{shortloc:directtest}
             mD-\kappa\lambda/n>\varepsilon\lambda,
\end{equation}
the conclusion follows.
\item \emph{Whole-count witness.} Bounds $B(\lambda)<0$ and $I(\lambda)\ge0$
satisfy
\[
 \rho_B(C-C_B)\le\sqrt n\pi_k\lambda B(\lambda)
                                      +\lambda I(\lambda)/n.
\]
It is enough that
\begin{equation}\label{shortloc:whole-test}
                  mB(\lambda)+I(\lambda)/N<-\varepsilon.
\end{equation}
\end{enumerate}
The witnesses may differ from cell to cell. A finite closed cover and a
separate treatment of $\lambda=0$ therefore suffice for a complete local
comparison.
\end{theorem}
\end{samepage}
\begin{proof}
Conditioning on the binomial count gives, with all atom conventions intact,
\begin{equation}\label{shortloc:count}
 Q=-\sum_{i=0}^{k}\pi_{k-i}\Pr(W_{k-i}>u+i)
       +\sum_{i=1}^{n-k}\pi_{k+i}\Pr(W_{k+i}\le u-i).
\end{equation}
Dropping the outgoing terms with $i>0$ gives $Q/\pi_k-f(u)\le-D$.
Use \eqref{shortloc:master} and $D>0$. For the second witness, use
$B<0$ before replacing $\sqrt n\pi_k$ by $m$, and $I\ge0$ before replacing
$n$ by $N$. These sign requirements are part of the test, not conventions.
\end{proof}

\subsection{Normalization and the repair, once for every witness}

Here is an explicit finite recipe for the constants required above.
Let $p\in[p_-,p_+]\subset(0,1/2)$, $n\ge N$, $\lambda\le\Lambda$,
$\theta\le\Theta$, and $|\delta|\le\Delta$, where $\Delta\ge1$ also includes
every binomial mode. Choose $0<r_-<p_- -\Delta/N$ and
$p_++\Delta/N<r_+<1/2$. Let $v_*>0$ be a lower bound for $p(1-p)$.
Choose positive $E,L,R$ bounding, on $[r_-,r_+]$,
\[
 |D'''(r\Vert p)|=\frac{1-2r}{r^2(1-r)^2},\qquad
 |(\log(r(1-r)))'|=\frac{1-2r}{r(1-r)},\qquad
 \frac1{12r(1-r)},
\]
respectively, where $D(r\Vert p)=r\log(r/p)+(1-r)\log((1-r)/q)$.
Then, with $Z(t)=Lt/2+R+Et^3/(6N)$,
\begin{equation}\label{shortloc:mass}
 \sqrt n\pi_k=\frac{\phi(0)}{\sqrt v}
       e^{-\delta^2/(2nv)+\zeta},\qquad n|\zeta|<Z(|\delta|).
\end{equation}
For completeness, the finite Stirling formula used here is
\[
 \log(j!)=(j+\tfrac12)\log j-j+\tfrac12\log(2\pi)+r_j,
 \qquad 0<r_j<1/(12j).
\]
If $a_j=\log(j!)-(j+1/2)\log j+j$ and $z=(2j+1)^{-1}$, the integrated
geometric series gives
$0<a_j-a_{j+1}=\sum_{i\ge1}z^{2i}/(2i+1)<1/[12j(j+1)]$.
Summing gives the remainder relative to $\lim a_j$. The integration-by-parts
recursion for $\int_0^{\pi/2}\sin^j t\,dt$ and the inequalities between three
successive integrals give Wallis's product; its finite product equals
$e^{4a_j-2a_{2j}}j/[2(2j+1)]$ and tends to $\pi/2$. Thus
$\lim a_j=\tfrac12\log(2\pi)$. Applying this formula to $n!,k!,(n-k)!$
and Taylor-expanding $D$ and the logarithmic prefactor proves
\eqref{shortloc:mass}; the factorial error is at most $1/[12nr(1-r)]$.

Consequently a valid mass lower bound is any positive
\begin{equation}\label{shortloc:m}
 m<\frac{.39886}{\sqrt{p_+(1-p_+)}}
          \left[1-\frac{\Delta^2/(2v_*)+Z(\Delta)}N\right].
\end{equation}
The same formula at a mode proves $\sqrt n\max_j\pi_j<1$ whenever
$v_*[1-Z(\Delta)/N]^2>.4^2$ and $Z(\Delta)<N$.
The rational bounds $.39886<\phi(0)<.4$ follow, for example, by the
alternating arctangent series in Machin's identity.

Differentiating \eqref{shortloc:f} gives
\begin{equation}\label{shortloc:logder}
 \log f'=-\tfrac12\log(1+s/v)+\delta^2/(2nv)
                   -(\delta+u)^2/(2n(v+s))-\zeta.
\end{equation}
Thus $n|\log f'|\le B_0\lambda+P(D,a)$ on $|\delta|\le D\le\Delta$,
$|u|\le a\le1/2$, where
\[
 B_0>\frac1{2v_*}+\frac{\Delta^2}{2Nv_*^2},\qquad
 P(D,a)=\frac{2Da+a^2}{2v_*}+Z(D).
\]
Choose $0<b<1$ above $[B_0\Lambda+P(\Delta,1/2)]/N$ and put
$F=(1-b)^{-1}$. The elementary estimates $e^z\le(1-z)^{-1}$ and
$1-e^{-z}\le z$ imply
\begin{equation}\label{shortloc:derivative}
 |f'-1|\le(B_1\lambda+P_1)/n,\qquad B_1=FB_0,\quad P_1=FP(D,a).
\end{equation}
This supplies both a derivative interval $l_f\le f'\le h_f$ and the
finer phase-error bound $|u-f(u)|\le a(B_1\lambda+P_1)/n$.

The exact count variance is
\begin{equation}\label{shortloc:variance}
 t_j=\mathbb E W_j^2=\lambda+(j-np)(s_1-s_0),\qquad
 |s_1-s_0|\le s/p\le c_v s,\quad c_v\ge1/p_-.
\end{equation}
It gives $\omega_-\lambda\le t_k\le\omega_+\lambda$ with
$\omega_\pm=1\pm c_vD/N$. All uses below require the checked condition
$\omega_->0$. Also
$t_{k+i}\le iV_*\lambda$ for $V_*=1+c_v(\Delta+1)/N$.
The binomial recursion gives, for every legal $i\ge1$,
\begin{equation}\label{shortloc:ratio}
 \frac{\pi_{k+i}}{\pi_k}\le R_0,
 \qquad R_0>1+\frac{p_++\Delta-1}
 {Nv_*[1-(\Delta-1)/(Np_-)]},
\end{equation}
provided the denominator is positive and $p_++\Delta<2$.
Indeed the first ratio has this bound and all subsequent rising factors
are at most one. No count outside $[0,n]$ is used in this recursion.

Finally \eqref{shortloc:support} fixes the signs inside the absolute cubic.
Expansion and centering give
\[
 \mathbb E|Y-p|^3=v(1-2v)+3pq(s_0+s_1)
                   +p\mathbb E U_1^3-q\mathbb E U_0^3
 \ge v(1-2v)+(3p-d)s.
\]
The inequality $(1+z)^{3/2}\le1+3z/2+3z^2/8$ for $z\ge0$ yields
\begin{equation}\label{shortloc:cubic}
 (v+s)^{3/2}(\rho-\rho_B)\ge s
     \left(6p-3p^2-\tfrac32-d-\frac{3(1-2v)s}{8v}\right).
\end{equation}
Let $Q_*>0$ be a lower bound for the bracket and
$m_c\le c_*/[p_+(1-p_+)+\Lambda/N]^{3/2}$.
The normal variance correction is nonpositive for $\delta\le0$;
for $\delta>0$ its derivative in $s$ is at most $.4\delta/(2v^{3/2})$.
Therefore
\begin{equation}\label{shortloc:repair}
 \mathcal R\le(P_*D_+-m_cQ_*)\lambda/n,
 \qquad P_*\ge\frac{.4}{2v_*^{3/2}},\quad D_+\ge\delta_+.
\end{equation}
One may take $D_+=D$, but retaining the positive part is often decisive.
All bounds concern $v+s$, which can exceed $1/4$.

\subsection{Two elementary tools for probabilities}

\begin{lemma}[A shifted quartic]\label{shortloc:quartic}
Suppose $W$ is a sum of independent centered variables bounded in absolute
value by $A$, and $t=\mathbb E W^2$. Let $K_0\ge(9+6\sqrt3)/4$ and $a>0$.
For $r\ge0$ and $\sigma\in\{-1,1\}$,
\begin{equation}\label{shortloc:quarticformula}
 \Pr(W+\sigma r>0)\ge
 \frac{(3a^2-A^2)t-3t^2-4Art-6r^2t-r^4
                      +3a^2r^2+2\sigma a^3r}{K_0a^4}.
\end{equation}
\end{lemma}
\begin{proof}
The polynomial $h_a(z)=(-z^4+3a^2z^2+2a^3z)/(K_0a^4)$ is below
$\mathbf1_{\{z>0\}}$: it is nonpositive for $z\le0$, and its maximum for
$z>0$ is at $z/a=(1+\sqrt3)/2$ and is at most one. Independence gives
$|\mathbb E W^3|\le At$ and $\mathbb E W^4\le A^2t+3t^2$.
Expand $\mathbb Eh_a(W+\sigma r)$ and apply these two inequalities.
The value at the threshold is zero, so atoms require no limiting argument.
\end{proof}

\begin{lemma}[One common tail bound]\label{shortloc:tails}
If independent centered $X_i\le A$ have total variance at most $V$, then,
for $b>0$, their sum has upper tail bounded by the minimum of one and
\begin{equation}\label{shortloc:tailmin}
 \begin{gathered}
 \frac{V}{V+b^2},\qquad
 \mathcal B(A,V,b)=e^{b/A}\left(\frac{V}{V+Ab}\right)^{b/A+V/A^2},\qquad
 \frac{\mathbb E(P_\mu-j)_+^3}{(b/A+\mu-j)^3},\\
 \mu=V/A^2,\qquad 0\le j<b/A+\mu.
 \end{gathered}
\end{equation}
where $P_\mu$ is Poisson. The same holds for a lower tail when $|X_i|\le A$.
All bounds apply to closed events, hence also to strict events.
\end{lemma}
\begin{proof}
Cantelli's bound follows by applying Markov to $(\sum X_i+V/b)^2$.
Taylor's integral remainder gives
$\mathbb E e^{zX_i}\le\exp\{\mathbb E X_i^2(e^{zA}-1-zA)/A^2\}$
for $z\ge0$. Multiply and use $z=A^{-1}\log(1+Ab/V)$ to get $\mathcal B$.
The zero-variance case is immediate.

For the third bound, smooth $f(x)=(x-t)_+^3$ by a nonnegative smooth
compactly supported density. Then $f'',f''',f''''\ge0$. For
$\mathcal Af(x)=f(x+A)-f(x)-Af'(x)$, the Taylor quotient
\[
 \frac{f(x+y)-f(x)-yf'(x)}{y^2}
       =\int_0^1(1-r)f''(x+ry)\,dr
\]
is nondecreasing in $y$. A centered $X\le A$ of variance $v$ consequently
satisfies $\mathbb E f(x+X)\le f(x)+(v/A^2)\mathcal Af(x)$.
For $\mathcal P_zf(x)=\mathbb E f(x+A(P_z-z))$,
$\partial_z\mathcal P_zf=\mathcal P_z\mathcal Af$ and
\[
 \partial_z^2\mathcal P_zf=\mathcal P_z\mathcal A^2f\ge0,\qquad
 \mathcal A^2f(x)=\int_0^A\!\int_0^A(A-r)(A-w)f''''(x+r+w)\,dr\,dw.
\]
Thus $\mathcal P_{v/A^2}f\ge f+(v/A^2)\mathcal Af$. Replace the independent
variables successively; translations and expectations preserve the derivative
signs and Poisson rates add. Increasing the total variance to $V$ is permitted
by Jensen after adding independent centered Poisson noise. Removing the
smoothing is justified by bounded actual upper supports and finite Poisson
third moments. Markov at $t=A(j-\mu)$ proves the displayed bound.
\end{proof}

These bounds have finite rational evaluations. For example,
\begin{align}\label{shortloc:poissonfinite}
 \mathbb E(P_\mu-j)_+^3={}&\mu^3+3(1-j)\mu^2+(1-3j+3j^2)\mu-j^3\\
 &+e^{-\mu}\sum_{h=0}^{\lfloor j\rfloor}\frac{\mu^h(j-h)^3}{h!}.\nonumber
\end{align}
Replacing $e^{-\mu}$ by $(\sum_{h=0}^{24}\mu^h/h!)^{-1}$ enlarges the
nonnegative correction. Alternatively sum the positive Poisson series
through 32; the remaining tail is at most its first term divided by
$1-\mu(34-j)^3/[34(33-j)^3]$, provided $j<33$ and that ratio is below one.
For Bennett, put $x=Ab/V$ and
$L=[(b/A+\mu)\log(1+x)-b/A]$. The positive series
$\log(1+x)=2\sum_{h\ge0}z^{2h+1}/(2h+1)$, $z=x/(2+x)$, gives lower
bounds for $L$; any positive lower exponential series gives an upper bound
for $e^{-L}$. Rounding a positive lower $L$ down before exponentiation is
also valid. These are choices of rational witnesses, not optimization claims.

\subsection{The nearest-count witness, including vanishing noise}

For $u=-r\le0$, add $-h_fr$ to \eqref{shortloc:quarticformula} with
$\sigma=1$; for $u=r\ge0$, add $l_fr$ with $\sigma=-1$.
Both bound $\Pr(W_k>u)+f(u)$ from below. For a cell
$t\in[t_-,t_+]$, $r\in[r_-,r_+]$, this expression is concave in $t$
(second derivative $-6/(K_0a^4)$), so only $t_-,t_+$ are needed.
At each endpoint, bound a monomial with positive coefficient at $r_-$
and one with negative coefficient at $r_+$. Any rational $a>0$ is a valid
witness. Other valid choices are
\begin{equation}\label{shortloc:ownother}
 \frac{r_-^2}{t_++r_-^2}-h_fr_+\quad(u\le0),\qquad l_fr_-\quad(u\ge0).
\end{equation}
They follow from Cantelli and nonnegativity, respectively. Their maximum
or one selected witness is valid without an optimality proof.

For incoming counts use $t_{k+i}\le iV_*\lambda$. On $u\in[-r_+,-r_-]$
the first distance is $1+r_-$ and every later distance is at least $i$.
On $u\in[r_-,r_+]$ they are $1-r_+$ and $i(1-r_+/2)$ for $i\ge2$.
Let $b,b_o$ be these first and later distances. On $\lambda\in[L,U]$,
let $T_1$ be any bound from Lemma~\ref{shortloc:tails} at $(A,V_*U,b)$.
Since $\mathcal B(A,iV,ib)=\mathcal B(A,V,b)^i$, all later counts total at most
$B^2/(1-B)$ for $B=\mathcal B(A,V_*U,b_o)<1$. Thus
\begin{equation}\label{shortloc:midtest}
 D=O-R_0\{T_1+B^2/(1-B)\}>0,\qquad
 mD-\kappa U/N>\varepsilon U
\end{equation}
proves the entire cell when $\kappa\ge0$. The selected MID normalizers
explicitly verify this sign. This is the common central and left MID proof.

A cell touching $\lambda=0$ requires estimates proportional to $\lambda$;
a positive lower endpoint may not be substituted. The following rescaling
provides them. Put $r=zt$, $t_*=\omega_+U$ and $d_a=K_0a^4$.
The shifted quartic plus $br$, divided by $t$, becomes
\begin{equation}\label{shortloc:scaledown}
 F_{\sigma,b}(z,t)=\frac{3a^2-A^2-3t}{d_a}
 +\left(\frac{2\sigma}{K_0a}+b-\frac{4At}{d_a}\right)z
 +\frac{3a^2t-6t^2}{d_a}z^2-\frac{t^3z^4}{d_a}.
\end{equation}
It extends to $t=0$ and is concave in $t$. On $0\le z\le z_0$, replacing
$z^4$ by $z_0^2z^2$ gives a lower quadratic at $t=t_*$; at $t=0$ it is
linear. Endpoints and a clipped quadratic vertex therefore give an exact
lower bound $J(\sigma,b,z_0)$ for the whole rectangle.

For negative phase choose $z_-,c_->0$ and verify
\begin{gather}\label{shortloc:tinynegative}
 J(1,-h_f,z_-)>c_-,\quad
 \frac{z_-^2}{1+z_-^2t_*}-h_fz_->c_-,\quad
 \frac{2z_-}{(1+z_-^2t_*)^2}>h_f,\\
 \frac1{1+4t_*}-\frac{h_f}{2}>c_-t_*,\qquad z_-t_*<1/2.\nonumber
\end{gather}
The quartic handles $r\le z_-t$. Beyond that point use
$r^2/(t+r^2)-h_fr$. Its derivative at $r=z_-t$ is positive by the third
gate; its second derivative has just one sign change, so any later critical
point is a maximum. Its minimum is at $z_-t$ or $1/2$, covered by the other
gates. Hence the own count exceeds $c_-t$.
For positive phase choose $z_+,c_+>0$ with
\begin{equation}\label{shortloc:tinypositive}
 J(-1,l_f,z_+)>c_+,\qquad l_fz_+>c_+,\qquad z_+t_*<1/2.
\end{equation}
The quartic handles $r\le z_+t$ and the normal term handles the rest.
The own count exceeds $c_+t$ and also exceeds $l_fr$.

For incoming tails proportional to $\lambda$, fix a Poisson anchor $j$.
A positive series term of index $h$ divided by
$\lambda(b/A+\mu-j)^3$ has logarithmic derivative, multiplied by $\lambda$,
\[
 h-1-\mu-\frac{3\mu}{b/A+\mu-j},\qquad \mu=V_*\lambda/A^2.
\]
If $b/A-j>0$ and this is nonnegative at $\lambda=U$ for the first positive
index, it is nonnegative for every term and all $0<\lambda\le U$.
Therefore the Poisson bound divided by $\lambda$ is largest at $U$.
Likewise $B^2/[\lambda(1-B)]$ is increasing whenever
$\lambda B'/B\ge1/2$. Differentiating the Bennett formula gives
\[
 \lambda B'/B=b/A-\mu\log(1+b/(A\mu));
\]
The function $g(\mu)=b/A-\mu\log(1+b/(A\mu))$ decreases because
$\log(1+x)>x/(1+x)$ for $x>0$. Moreover the logarithmic derivative of
$B^2/[\lambda(1-B)]$, multiplied by $\lambda$, is
$2g-1+gB/(1-B)\ge0$ when $g\ge1/2$. Thus checking the gate at $U$
proves the whole interval. A rational upper bound for the logarithm
verifies that gate. On cells using the fixed positive bound $l_fr_-$ instead,
unscaled endpoint tails suffice: after a positive bracket is checked, divide
it by $U$ to obtain a lower bound per $\lambda$. These arguments, followed by
\eqref{shortloc:directtest}, prove TINY uniformly down to arbitrary positive
noise, including a deterministic conditional component.

\subsection{The whole-count witness}

Assume $c=1-d_0>1/2$. Couple all $2n$ conditional variables independently and
set $W_j=\sum_{i\le j}U_{1,i}+\sum_{i>j}U_{0,i}$, extending it constantly
outside $[0,n]$. Every increment of $j+W_j$ is at least $c$. Define
\[
 K=\max\{j\in\mathbb Z:j+W_j\le k+u\},\qquad R=K-k.
\]
An independent binomial selector has the correct actual count laws, so
\begin{equation}\label{shortloc:cutoff}
                  \Pr(\textstyle\sum Y_i\le k+u)=\mathbb EF_B(k+R).
\end{equation}
Let $\chi=(n+1)p-k<2$ and $m_B=\max_j\pi_j$. For every integer $r$,
\begin{equation}\label{shortloc:curvature}
 F_B(k+r)-F_B(k)-r\pi_k\le\frac{m_B}{(n+1)v}
 \left[(\chi-1)_+r_++\chi_+\binom{r_-}{2}+\binom{r_-}{3}\right].
\end{equation}
Indeed, the positive part of $\pi_{k+i}-\pi_{k+i-1}$ is at most
$m_B(\chi-i)_+/((n+1)v)$, by
$\pi_j-\pi_{j-1}=\pi_j((n+1)p-j)/(p(n-j+1))$.
Double summation gives \eqref{shortloc:curvature}: for $r>0$ only $i=1$
can contribute; for $r=-h<0$ use $(\chi+i)_+\le\chi_++i$ and sum the
multiplicities $h-1-i$. Zero extension makes the same calculation valid
outside the legal count range.

Write $t=\mathbb EW_k^2$. Crossing monotonicity and Cauchy--Schwarz give
\begin{equation}\label{shortloc:crossingmoments}
 \mathbb ER_+\le\frac{t}{(c-u)^2},\quad
 \mathbb E\binom{R_-}{2}\le\frac{t}{[c+\min(u,c/3)]^2},\quad
 \mathbb E\binom{R_-}{3}\le\frac{t\sqrt{A^2+3t}}{(2c-1/2)^3}.
\end{equation}
For the first, $R=i\ge1$ implies $-W_k\ge ci-u$ and
$(ci-u)^2\ge i(c-u)^2$. For $R=-h\le-2$, the next count is outside the
closed event, so $W_k>c(h-1)+u$. Put $j=h-1$.
The function $j(j+1)/(2(cj+u)^2)$ decreases for $j\ge1$ if $u\le c/3$;
for larger $u$ reduce to that endpoint. Also
$(j^3-j)/(6(cj-1/2)^3)$ decreases for $j\ge2$, since its derivative has
the sign of $2cj-(3j^2-1)/2<0$. These give the last two bounds using
$\mathbb E(W_k)_+^3\le\sqrt{\mathbb EW_k^2\mathbb EW_k^4}$.

\paragraph{The marked crossing.}
Let $D_i=U_{1,i}-U_{0,i}$, $\vartheta=\mathbb E(D_i)_+=\mathbb E(D_i)_-$,
$0<\eta\le1$, and $R_f=\mathbb E\{u+\eta-W_k\}$. The two deleted laws are
\[
 A_0\overset{d}=kU_1+(n-k-1)U_0,\qquad
 A_1\overset{d}=(k-1)U_1+(n-k)U_0.
\]
Then
\begin{align}
 \mathbb ER&\le u+\eta-R_f+J,\label{shortloc:mark}\\
 J&=\frac{\vartheta}{\eta}\sum_{i\ge1}
 \{\Pr(A_0\le u+d_0-ci)+\Pr(A_1>u+\eta+i-A)\}.\nonumber
\end{align}
Here are the pathwise details. Put $z=u-W_k$ and mark
\[
 Z_f=\sum_{i=1}^{n-k}(D_{k+i})_-\mathbf1_{\{W_k\le u-ci\}},\quad
 Z_b=\sum_{i=1}^{k}(D_{k-i+1})_+\mathbf1_{\{W_k>u+\eta+i\}}.
\]
If $R>z+\eta$, the crossing activates the marks through $R$: for $R>0$
use $z\ge cR$, while for $R<0$ use $z<R-\eta$. In either case
$R\le z+Z_f+Z_b$. Hence if $h=R-\lfloor z+\eta\rfloor>0$,
$Z_f+Z_b>h-1+\eta\ge\eta h$. Thus
$R\le\lfloor z+\eta\rfloor+(Z_f+Z_b)/\eta$ on every path.
On an active forward mark $U_0>U_1\ge-d_0$; deleting its pair implies
$A_0\le u+d_0-ci$. On a backward mark delete $U_1\le A$ to obtain the
strict event for $A_1$. The remaining sum is independent of the entire
pair, so the expectations factor. Extending to all $i\ge1$ adds only
nonnegative bounds. Finally $|R|\le1+(|W_k|+|u|)/c$ proves integrability
and permits expectations. No coupling has replaced either deleted marginal.

\paragraph{A finite bound for every mark.}
Fix $\lambda\in[L,U]$, $L>0$, and $t\le\omega_+\lambda$; put
$T_L=\omega_+L$, $T_U=\omega_+U$. On a phase subcell $u\in[v_l,v_h]$,
for $i=1,2$ let $b_1=ci-d_0-v_h>0$ and $b_2=i-A+v_l>0$.
Lemma~\ref{shortloc:tails} gives a bound $\omega_+\lambda K_i^{\rm ind}$,
where
\[
 K_i^{\rm ind}=\sum_{j=1}^2\min\left\{\frac1{b_j^2+T_L},
                          \frac{\overline P(A,T_U,b_j)}{T_L}\right\}.
\]
The denominator $T_L$ is a lower bound for the chosen upper envelope
$\omega_+\lambda$, not for any actual variance.
Alternatively couple $A_0=Z+U_1$, $A_1=Z+U_0$ with
$Z\overset d=(k-1)U_1+(n-k-1)U_0$. The two events project into
$Z\le-\alpha$ and $Z>\beta$, where
$\alpha=ci-2d_0-v_h>0$, $\beta=i-A-d_0+v_l>0$.
For any $z\in(-\alpha,\beta)$ their sum is at most
$(\omega_+\lambda+z^2)/\min(z+\alpha,\beta-z)^2$.
Consequently a second kernel is
\[
 K_i^{\rm com}=\frac{1+z^2/T_L}{\min(z+\alpha,\beta-z)^2},\qquad
 z=\max\{-T_L/\beta,\min((\beta-\alpha)/2,T_L/\alpha)\}.
\]
This center is inside the interval; differentiating the two quadratic
ratios on either side of the midpoint proves its optimality, although
any interior rational center suffices. Take the smaller kernel on each
phase subcell and the maximum over their closed cover. Denote it $K_i$.
For $i\ge3$ use fourth moments and the convex midpoint integral for
$(cx-b)^{-4}$. On $u\in[u_l,u_h]$ this gives
\begin{equation}\label{shortloc:K}
 K=K_1+K_2+\frac{A^2+3T_U}{3c}
  \left\{\frac1{((5/2)c-A-u_h)^3}
                 +\frac1{((5/2)c-A+u_l)^3}\right\},\qquad
 J\le\frac{\vartheta}{\eta}\omega_+\lambda K.
\end{equation}
All denominators must be positive. Indeed, the tail fourth moment is at
most $A^2t+3t^2$, and integration from $5/2$ to infinity covers every
remaining integer, including the harmless added illegal counts.

If the same-law resource supplies $nd<G$ and
$G^2<c_vNL$, then
\begin{equation}\label{shortloc:theta}
 \vartheta\le\frac{c_vG\lambda}{c_vn\lambda+G^2}.
\end{equation}
To prove this, for centered $X\ge-b$, $\mathbb EX^2\le V$ and $b^2\le V$,
the quadratic $b(X-V/b)^2/(b+V/b)^2$ majorizes $X_-$ on $[-b,\infty)$.
Its expectation is at most $bV/(V+b^2)$. Apply this to $D_i\ge-d$,
using $\operatorname{Var}(D_i)=s_0+s_1\le c_v\lambda/n$, and monotonicity
in $b\le\sqrt V$. The case $d=0$ is immediate. With $\eta=z\lambda$,
$z>0$, $zU\le1$, \eqref{shortloc:K} becomes
\begin{equation}\label{shortloc:Jfinal}
 \frac J\lambda\le\frac{\omega_+Kc_vG}{z(c_vN\lambda+G^2)}.
\end{equation}

\paragraph{The fractional-part gain.}
If each centered summand has support span at most $b<1/2$ and the sum has
variance $t$, then, for any real $w$ and
$0<r_*\le2\sin^2(\pi b)/b^2$,
\begin{equation}\label{shortloc:harmonic}
             \mathbb E\{w-W\}\ge\frac5{23}(1-e^{-r_*t}).
\end{equation}
Indeed $1-\cos(2\pi x)\le(23/5)x$ for $0\le x\le1$.
A finite proof of this sole scalar inequality is recorded below.
For an independent copy $Z'$ of one summand, concavity of sine on
$[0,\pi b]$ gives
\[
 |\mathbb Ee^{2\pi iZ}|^2
 =1-2\mathbb E\sin^2\pi(Z-Z')
 \le e^{-4\sin^2(\pi b)\operatorname{Var}(Z)/b^2}.
\]
Multiply square roots over the independent summands and use periodicity.
This argument includes zero variances and integer collisions.
Here the conditional spans are at most $A+d_0$, so this is the $b$ to use.

For a reproducible scalar proof, $x\ge10/23$ follows from $1-\cos\le2$.
On $[0,10/23]$ use the degree-ten alternating upper polynomial for
$1-\cos(2\pi x)$, with directed coefficients:
\[
 Q(x)=\sum_{j=1}^5\frac{(-1)^{j+1}(2\pi_jx)^{2j}}{(2j)!},\qquad
 \pi_j=\begin{cases}3.141592653589794,&j\text{ odd},\\
                       3.141592653589793,&j\text{ even}.\end{cases}
\]
Put
$P(x)=23/5-Q(x)/x$. On the eight intervals with left endpoint $a_j=5j/92$
and width $w=5/92$, its degree-nine Bernstein coefficients are
\[
 b_{jk}=\sum_{i=0}^k\frac{\binom ki}{\binom9i}
           \sum_{h=i}^9[x^h]P(x)\binom hi a_j^{h-i}w^i,
 \qquad 0\le j<8,\quad0\le k\le9.
\]
The 80 rational coefficients are positive; the retained scalar acceptor
reconstructs them. Expansion in powers of the unit-interval coordinate
proves the formula; Bernstein weights are nonnegative and sum to one.
The alternating remainder and the stated rational bounds on $\pi$ follow
from the same elementary series above. This completely specifies the
finite scalar obligation. A rational rate follows by replacing
$\sin(\pi b)$ with $x-x^3/6+x^5/120-x^7/5040$, $x=3.14159b$, on the
checked alternating-series domain.

Set $r=r_*\omega_-$, require $rU<3$, and define
\begin{equation}\label{shortloc:gain}
 g=\frac5{23}\left(\frac{r^3L^2}{6}-\frac{r^4L^3}{24}
                     +\frac{r^5L^4}{120}-\frac{r^6L^5}{720}\right).
\end{equation}
The sixth-order Taylor lower bound for $1-e^{-r\lambda}$ gives
$R_f/\lambda\ge5r/23-5r^2\lambda/46+g$. The discarded higher-order
part is increasing as a function of $\lambda$ for $r\lambda<3$:
its derivative has the sign of $48-18x+(24/5)x^2-x^3>0$ on $[0,3]$.
For instance this polynomial has positive degree-three Bernstein
coefficients $48,30,132/5,51/5$ there. Thus freezing it at $L$ is valid.

\paragraph{One final expression.}
On a signed phase cell find an interval $[x_l,x_h]$ containing $H/\ell$
by its four endpoint quotients, using positive bounds for $\ell$.
Let
\[
 D=\max(|x_l-u_h|,|x_h-u_l|),\quad
 D_+=\max(0,x_h-u_l),\quad a=\max(|u_l|,|u_h|),\quad b_\chi=p_+-x_l.
\]
The same formula allows coarser upper bounds $D,D_+,b_\chi$, as used by
the central certificate. Check $D\le\Delta$ and $b_\chi+u_h<2$.
With $\omega_\pm=1\pm c_vD/N$, the curvature coefficient is
\begin{equation}\label{shortloc:V}
 V_0=\frac{\omega_+}{v_*}\left[
 \sup_{u_l\le u\le u_h}\left\{
 \frac{(b_\chi+u-1)_+}{(c-u)^2}
 +\frac{(b_\chi+u)_+}{[c+\min(u,c/3)]^2}\right\}
 +\frac{\sqrt{A^2+3T_U}}{(2c-1/2)^3}\right].
\end{equation}
On a rational phase subcell, both numerators are bounded at the right
endpoint; the first denominator is smallest there, and the second is
smallest at the left endpoint. This proves the supremum bound, including
both kinks. Equations \eqref{shortloc:cutoff}--\eqref{shortloc:crossingmoments},
$\sqrt n m_B\le1$, and \eqref{shortloc:master} now give the whole-count
witness with
\begin{align}\label{shortloc:BI}
 B(\lambda)&=z+\frac{\omega_+Kc_vG}{z(c_vN\lambda+G^2)}
                      -\frac{5r}{23}+\frac{5r^2\lambda}{46}-g,\\
 I(\lambda)&=V_0+aB_1+\frac{aP_1}{\lambda}+P_*D_+-m_cQ_*.
 \nonumber
\end{align}
We used $|u-f(u)|\le a(B_1\lambda+P_1)/n$ and $\sqrt n\pi_k\le1$.
Every constant is fixed before $\lambda$ varies. Both $B$ and
$mB+I/N$ are convex on $\lambda>0$; $I$ is decreasing. Thus their two
endpoint gates $B<0$, $I>0$, $mB+I/N<-.001$ prove the entire closed noise
cell. This proves all three LOW variants by exactly the same argument.

\subsection{The selected finite cover and the zero-noise endpoint}

The retained tables contain 45 central probability intervals covering
$[.38,.47]$ and 38 left intervals covering $[.304,.38]$. Each row supplies
its own $\Theta$, signed $H$ bounds, $d_0$, $\Lambda$, and global amplitude;
the confinement profiles may improve the amplitude with $\lambda$.
All these are hypotheses of the present section, acquired earlier without
using its conclusion. The shared resource has the form
\[
 nD_{\rm mom}<\frac{15000}{997}T_H,\qquad
 D_{\rm mom}\ge\alpha d+\beta\theta+P.
\]
Consequently a noise cell $[L,U]$ has
\[
 nd<G=\{(15000/997)T_H-NP-\beta\ell_-^2L\}/\alpha.
\]
Here $T_H$ is the maximum of the acquired concave quadratic budget on the
closed signed $H$ cell, further bounded by its scalar ceiling. It already
contains the phase charge. If $G\le0$, the entire cell is impossible and
receives no score test; the exclusion is strict even at $G=0$.

For positive noise the complete cover is the following. Each interval is
intersected with the row ceiling $\Lambda$; shared endpoints are included.
\begin{center}\small
\begin{tabular}{@{}llll@{}}\toprule
Noise & Central phase & Left phase & Witness\\\midrule
$(0,.05]$ & all & all & nearest count (scaled)\\
$[.05,.06]$ & $|u|\ge.09$ & $|u|\ge.09$ & nearest count\\
$[.06,.08]$ & $|u|\ge.13$ & $|u|\ge.14$ & nearest count\\
$[.08,.10]$ & $|u|\ge.18$ & $|u|\ge.19$ & nearest count\\
$[.10,.12]$ & $|u|\ge.29$ & $u\le-.29$ & nearest count\\
$[.10,.12]$ & --- & $u\in[.29,.5]$ & whole count\\
$[.05,\Lambda]$ & remaining phase & remaining phase & whole count\\\bottomrule
\end{tabular}
\end{center}
Beyond $.12$ the last row covers the full phase interval. The three LOW
objects are one central cover, one left cover with phase slices, and one
extra positive left cover. Their distinction is only the closed domain
and the chosen outward constants. Signed $H$ cells meet at zero; nearest
count ties meet at $u=\pm1/2$. The fixed coverage checker also checks every
row-ceiling singleton and every entirely excluded LOW row.

The common TINY parameters are
\[
 N=47,\quad p\in[.304,.47],\quad |H|\le1.393,\quad
 \Theta=.0154,\quad d_0=.0181,\quad A=.425,\quad\Lambda=.05.
\]
For its scaled own witness use $a=.4784$, $K_0=9/(4\cdot.4641)$,
$(z_-,c_-)=(1.66,.514)$ and $(z_+,c_+)=(1.2,1.036)$.
Its six phase cells are $[-.5,0],[0,.1],\ldots,[.4,.5]$.
The first two use scaled Poisson anchors $1.25,1$, respectively; the other
four use the fixed positive normal bound. The exact normalizers and gates
specified in the accompanying certificate map verify these constants.
The central MID amplitudes are $.35,.35,.30,.275$ in its four rows.
The left ones are $.375,.375,.325,.30$, each further reduced by the row's
global amplitude. No interpolation of the tables is used.

All six accepted comparisons give the same strict conclusion
\[
                        \rho_B(C-C_B)<-.001\lambda.
\]
Their proofs have now reduced to Theorem~\ref{shortloc:theorem}, the shared
normalization, and the two explicitly derived witnesses. The arithmetic
objects still have their original domains and targets; no certificate
branch has been silently deleted or relabeled.

If $\lambda=0$, conditional centering and $p,q>0$ imply $U_0=U_1=0$
almost surely. The law is exactly Bernoulli. For an attained positive closed
maximum its threshold is a binomial atom: lowering a noninteger threshold
to the preceding atom leaves the CDF unchanged and strictly improves the
normal term, while thresholds below the first atom have nonpositive score.
Thus $u=0$ at such a maximum and $C=C_B$. Schulz's all-probability
Bernoulli theorem~\cite[Theorem~1]{Schulz2016} gives $C_B\le C_E$ for every
$0<p<1$ and in particular for all the indices $n\ge47$ used here. The same
external theorem closes the positive-noise conclusion $C<C_B\le C_E$.
At zero noise $C=C_B\le C_E$ already contradicts the strict contrary
hypothesis $C> C_E$ used by the main first-failure argument; no
finite Bernoulli reserve is needed. No estimate
in this section divides by zero or passes from fixed-law convergence to
uniformity in varying laws.


% Source: joint_resource.tex
\section{Noise and inward displacement share one budget}



The local comparison near a two-point law uses both a variance bound and
a bound on the displacement of its inner support endpoints. These two
bounds come from the same inequality. Keeping their joint constraint
shrinks the allowed displacement as the variance increases. We give an
elementary, self-contained proof, including the gap identity on which it
rests. The final section explains its use in the exact local comparisons;
the numerical acquisition of the input bounds is a separate obligation.


\subsection{The idea}
Suppose a proof has established, for nonnegative quantities $d$ and
$\lambda$, that
\[
             \frac{d}{d_*}+\frac{\lambda}{\Lambda}\le1.
\]
Replacing this by $d\le d_*$ and $\lambda\le\Lambda$ loses information.
For example, if $\lambda\ge\Lambda/2$, the original inequality gives
$d\le d_*/2$. The local characteristic-function and tail bounds improve
when $d$ decreases, so this elementary observation is useful precisely
where the noise is largest.

Here $d$ measures how far the actual inner endpoints lie inside the
conditional means, and $\lambda$ is the total conditional noise variance
in a sum of $n$ independent copies. The point is to establish the joint
inequality for one and the same marginal distribution.

\subsection{Definitions and the gap identity}
Let $X$ be bounded, $\E X=0$, and $\E X^2=1$. Assume that its closed
support has a gap $(-A,B)$ containing zero, with $A,B>0$ and endpoints in
the support. Write
\[
 \sigma=\sqrt{10},\quad \rho=\E|X|^3,\quad \gamma=\E X^3,
 \quad p=\Pp(X>0),\quad q=1-p,\quad m=\E X_+=\E X_-.
\]
Thus $0<p,q<1$ and $m>0$. For every bounded open gap $g=(a,b)$ of the
closed support, let $\tau_g$ be the constant value on that gap of
$\tau(t)=\E[X\mathbf1_{\{X\ge t\}}]$. Set
\[
 J=\sum_g\tau_g(b-a)^2,
 \qquad D=(3+\sigma)\rho-\gamma-3J.
\]
All summands in $J$ are nonnegative. The bounds below also establish
convergence of its outer-gap part. Endpoint atoms cause no ambiguity:
$\tau_g$ denotes the value in the open gap.

Define
\[
 T_+=\E X_+^3-mB^2,\quad T_-=\E X_-^3-mA^2,
 \qquad R_+=\E X_+^2-Bm,\quad R_-=\E X_-^2-Am.
\]
Each is nonnegative because the nonzero positive and negative magnitudes
are at least $B$ and $A$, respectively. Put $J_o=J-m(A+B)^2$.

\begin{lemma}[Gap remainder]\label{lem:gap}
The outer gaps satisfy
\[
 J_o\le T_++T_--2BR_+-2AR_-.
\]
Consequently
\begin{align}
 D={}&m\bigl(\sqrt{\sigma+1}A-\sqrt{\sigma-1}B\bigr)^2
 +(\sigma-1)T_++(\sigma+1)T_-\notag\\
 &\hspace{28mm}+3(T_++T_--J_o).\label{eq:exact}
\end{align}
In particular $D\ge0$.
\end{lemma}
\begin{proof}
On a positive outer gap $(u,v)$, where $u\ge B$,
\[
 (v-u)^2\le 2\int_u^v(t-B)\,dt.
\]
Multiply by $\tau_g$ and sum. Tonelli's theorem, applied to nonnegative
integrands, gives
\[
 J_o^+\le2\int_B^\infty\tau(t)(t-B)\,dt
 =\E[X_+(X_+-B)^2]=T_+-2BR_+.
\]
Reflection proves the corresponding negative bound. Their sum is the
first assertion. To check the identity, substitute
$\rho=m(A^2+B^2)+T_++T_-$,
$\gamma=m(B^2-A^2)+T_+-T_-$, and
$J=m(A+B)^2+J_o$. The cross coefficient in the square is $-6AB$
because $(\sigma+1)(\sigma-1)=9$. Nonnegativity follows from the first
assertion and $R_\pm,T_\pm\ge0$.
\end{proof}

\subsection{The simultaneous budget}
Let $V=\E[X\mid\mathbf1_{\{X>0\}}]$, and let $W=X-V$. Define
\[
 \theta_+=\E[W^2;X>0],\qquad
 \theta_-=\E[W^2;X<0],\qquad \theta=\theta_++\theta_-.
\]
The two conditional means are $\bar B=m/p$ and $-\bar A=-m/q$.
Therefore
\[
 m^2=(1-\theta)pq,\qquad
 \Delta_A=\bar A-A\ge0,\qquad\Delta_B=\bar B-B\ge0.
\]
In particular $\theta<1$. The following inequality retains three distinct
nonnegative costs: departure of $p$ from its distinguished value,
inward displacement, and conditional variance.

\begin{theorem}[One-law gap budget]\label{thm:budget}
With $p_*=(4-\sqrt{10})/2$,
\begin{align}
 D\ge{}&\frac{(6+2\sigma)(1-\theta)^{3/2}}{\sqrt{pq}}(p-p_*)^2
 +6m\bigl\{(A+B)(\Delta_A+\Delta_B)+\Delta_A\Delta_B\bigr\}\notag\\
 &+3(\sigma+1)B\theta_++3(\sigma+3)A\theta_-.
 \label{eq:budget}
\end{align}
\end{theorem}
\begin{proof}
Conditional expansion of the cubic gives
\[
 \E X_+^3=m\bar B^2+
 \E[(X+2\bar B)(X-\bar B)^2;X>0]
 \ge m\bar B^2+3B\theta_+.
\]
The reflected inequality has $\bar A,A,\theta_-$ in place of
$\bar B,B,\theta_+$. Also
$R_+=\theta_++m\Delta_B$ and
$R_-=\theta_-+m\Delta_A$.
Substitute these statements in Lemma~\ref{lem:gap}. The variance
coefficients become $3(\sigma+1)B$ and $3(\sigma+3)A$.
The remaining deterministic expression is
\begin{align*}
 &m\{(\sigma+1)\bar A^2+(\sigma-1)\bar B^2-6\bar A\bar B\}\\
 &\qquad+6m\{\bar A\bar B-AB+B\Delta_B+A\Delta_A\}.
\end{align*}
The second brace equals
$(A+B)(\Delta_A+\Delta_B)+\Delta_A\Delta_B$.
For the first, substitute $\bar A=m/q$, $\bar B=m/p$ and use
\[
 (\sigma+1)p^2+(\sigma-1)q^2-6pq
       =(6+2\sigma)(p-p_*)^2,
 \qquad \frac{m^3}{p^2q^2}=\frac{(1-\theta)^{3/2}}{\sqrt{pq}}.
\]
This proves every term in \eqref{eq:budget}.
\end{proof}

\subsection{A triangle of allowed noise and displacement}
Introduce Bernoulli coordinates
\[
 X=\ell(B_0+U-p),\quad B_0=\mathbf1_{\{X>0\}},\quad
 \ell=\bar A+\bar B=\frac m{pq},\quad \E[U\mid B_0]=0.
\]
Then
\[
 \ell^2=\frac{1-\theta}{pq},\quad m\ell=1-\theta,
 \qquad d=\frac{\Delta_A+\Delta_B}{\ell},\quad
 \lambda=n\E U^2=\frac{npq\theta}{1-\theta}.
\]

\begin{theorem}[Joint inward/noise constraint]\label{thm:triangle}
Suppose $n\ge N>0$, and bounds for this same law give
\[
 nD\le N d_D,\quad \theta\le\Theta<1,\quad pq\le v_*,
 \quad A+B\ge S_*>0,
\]
and
\[
 D\ge b+6mS_*(\Delta_A+\Delta_B)+K\theta,
 \qquad b\ge0,\quad K>0.
\]
Set $r=d_D-b$. If $r>0$, define
\[
 d_* =\frac{r}{6(1-\Theta)S_*},\qquad
 \Lambda=\frac{Nv_*r}{K(1-\Theta)}.
\]
Then
\begin{equation}\label{eq:triangle}
                 \boxed{\displaystyle
                   \frac d{d_*}+\frac\lambda\Lambda\le1.}
\end{equation}
If $r<0$, the premises are impossible. If $r=0$, they force
$d=\lambda=0$.
\end{theorem}
\begin{proof}
Multiply the assumed joint lower bound for $D$ by $n$. Since all
coefficients being weakened are nonnegative,
\begin{align*}
 Nd_D\ge nD
 &\ge nb+6n(1-\theta)S_*d+K\frac{(1-\theta)\lambda}{pq}\\
 &\ge Nb+6N(1-\Theta)S_*d+\frac{K(1-\Theta)}{v_*}\lambda.
\end{align*}
Subtract $Nb$. For $r>0$, division by $Nr$ proves
\eqref{eq:triangle}. For $r\le0$, positivity of both coefficients gives
the stated alternatives without dividing by zero.
\end{proof}

The hypotheses are direct consequences of Theorem~\ref{thm:budget}.
For example, on $p\in I=[l,u]$, take positive lower bounds $A_*,B_*$,
put $S_*=A_*+B_*$ and
\[
 v_* =\max_{p\in I}p(1-p),\qquad
 K=\min\{3(\sigma+1)B_*,3(\sigma+3)A_*\},
\]
and use
\[
 b=\frac{(6+2\sigma)(1-\Theta)^{3/2}}{\sqrt{v_*}}
           \operatorname{dist}(p_*,I)^2.
\]
Drop the nonnegative product $6m\Delta_A\Delta_B$.
The remaining bound is the hypothesis of Theorem~\ref{thm:triangle}.

\begin{corollary}[Closed noise bands and outward rounding]
If $0\le L\le\lambda\le U\le\Lambda$, then
\[
                       d\le d_*(1-L/\Lambda).
\]
Both intercepts in \eqref{eq:triangle} may first be replaced by larger
positive numbers. The same conclusion then holds using those enlarged
intercepts and the enlarged noise ceiling.
\end{corollary}
\begin{proof}
Rearrange \eqref{eq:triangle} and use $\lambda\ge L$.
Increasing either denominator weakens its nonnegative summand, so it
preserves the joint inequality. This also proves the rounding assertion.
\end{proof}

\subsection{Use in the computer-assisted proof}
The acquired tables store a weighted defect bound $nD\le Nd_D$ and the
two intercepts $d_*,\Lambda$. Those are simultaneous statements about
the actual extremizing law. An unweighted bound $D\le d_D$ alone would
not justify the $n$-uniform noise ceiling.

On each closed noise interval, the local programs recompute the inward
gap ceiling from the corollary. This improves the available support
separation, the conditional cubic correction, the Poisson and common
background tail distances, and the characteristic-function decay rate.
The proof uses the resulting smaller allowed set. It does not add two
different estimates that each spend the entire budget.

In the current exact central comparison the amplitude ceilings on the
noise bands are
\[
\begin{array}{c|cccc}
\lambda &[.05,.08]&[.08,.10]&[.10,.12]&[.12,\Lambda]\\\hline
|U|\text{ ceiling}&.375&.325&.275&.25.
\end{array}
\]
Each is intersected with an independently acquired first amplitude
bound. A separate support/concentration argument proves these amplitude
profiles before the final local score comparison uses them. The exact
central profile replay checks 248 signed configurations and 210,130
strict rational gates, with 14 malformed or false-object controls.

The analytic statement proved here is Theorem~\ref{thm:triangle}.
It does not, by itself, prove attainment, a Fourier remainder, a numerical
cutoff, or acquisition of any numerical resource table. Those edges
remain explicit in the project dependency map. The proof uses no
floating-point positivity claim and no limiting assertion about varying
marginal laws.


% Source: signed_cubic_repair.tex
\section{The sign of the cubic correction}

The cubic correction may have either sign, so its multiplier must be
bounded in the corresponding direction. The following argument uses the
stated support, moment and score bounds, supplied in the finite application
at $42\le n\le48$. Every conditional moment concerns the same iid marginal.

\subsection{Conditional centering and the exact cubic bound}
Let $B$ be Bernoulli with parameter $0<p<1/2$, put $q=1-p$ and $v=pq$,
and write $Y=B+U$, with $\E(U\mid B)=0$.
Let $U_0,U_1$ denote the conditional laws of $U$, and assume
\[
 U_0\le a_0<p,\qquad U_1\ge-a_1>-(1-p),\qquad
 a_0,a_1\ge0,\qquad d=a_0+a_1.
\]
Put $s_j=\E U_j^2$, $s=qs_0+ps_1$, $\lambda=ns$, and
\[
 X=\frac{Y-p}{\sqrt{v+s}},\quad
 \rho=\E|X|^3,\quad
 \rho_B=\frac{1-2v}{\sqrt v}.
\]
Finite third moments are assumed. The support conditions fix the sign of
$Y-p$ in each class, which is necessary for the following expansion.

\begin{lemma}[No positivity assumption on the cubic bracket]\label{cubic}
With the preceding notation,
\begin{equation}\label{bracket}
 (v+s)^{3/2}(\rho-\rho_B)\ge sQ,\qquad
 Q=6p-3p^2-\frac32-d-\frac{3(1-2v)s}{8v}.
\end{equation}
\end{lemma}
\begin{proof}
Expanding within each sign class and using conditional centering gives
\[
 \E|Y-p|^3=v(1-2v)+3pq(s_0+s_1)
                         +p\E U_1^3-q\E U_0^3.
\]
The one-sided support bounds imply
$\E U_1^3\ge-a_1s_1$ and $\E U_0^3\le a_0s_0$.
Moreover $3pq(s_0+s_1)\ge3p(qs_0+ps_1)=3ps$, because $q\ge p$.
Consequently
\[
                         \E|Y-p|^3\ge v(1-2v)+(3p-d)s.
\]
For $z\ge0$, Taylor's formula gives
$(1+z)^{3/2}\le1+3z/2+3z^2/8$: the third derivative is negative.
Subtract
\[
 \rho_B(v+s)^{3/2}
 =v(1-2v)(1+s/v)^{3/2}
 \le v(1-2v)+\tfrac32(1-2v)s+\frac{3(1-2v)s^2}{8v}
\]
from the preceding moment lower bound and use
$3p-\tfrac32(1-2v)=6p-3p^2-\tfrac32$.
\end{proof}

\subsection{The coefficient must be bounded in the correct direction}
Suppose $n\ge N$, $p\in[p_-,p_+]\subset(0,1/2)$,
$0\le d\le d_0$, and $0\le\lambda\le\Lambda$.
Choose $0<v_*\le p_-(1-p_-)$ and put
\[
 V^*=p_+(1-p_+)+\Lambda/N,\qquad
 Q_*=6p_--3p_-^2-\frac32-d_0
                     -\frac{3(1-2v_*)\Lambda}{8v_*N}.
\]
Then $Q\ge Q_*$: the first term increases with $p$, and
$(1-2v)/v$ decreases with $v$. In particular, $Q_*$ may be negative.
Already with $d=s=0$ its sign changes at $p=1-1/\sqrt2$.

Let $C$ be the score multiplying the moment correction, and assume
$0<c_*\le C\le c^*$. Choose positive multipliers satisfying
\[
        0<m_-\le\frac{c_*}{(V^*)^{3/2}},\qquad
        m_+\ge\frac{c^*}{v_*^{3/2}},
\]
and define the signed gain
\begin{equation}\label{gain}
             G_*:=\min\{m_-Q_*,m_+Q_*\}.
\end{equation}

\begin{proposition}[Signed gain and nonnegative repair]\label{repair}
Let $\Phi$ denote the standard normal distribution function. For an actual
count displacement $\delta$, set
\[
 \mathcal R=\sqrt n\left\{
 \Phi\!\left(\frac{\delta}{\sqrt{nv}}\right)
 -\Phi\!\left(\frac{\delta}{\sqrt{n(v+s)}}\right)\right\}
                         +C(\rho_B-\rho).
\]
If $D_+\ge\max(\delta,0)$ and $P_*\ge .4/(2v_*^{3/2})$, then
\begin{equation}\label{repair_bound}
 \mathcal R\le(P_*D_+-G_*)\frac{\lambda}{n}
 \le \kappa\frac{\lambda}{N},\qquad
 \kappa=\max\{0,P_*D_+-G_*\}.
\end{equation}
\end{proposition}
\begin{proof}
Write $A=C/(v+s)^{3/2}$. The hypotheses give
$0<m_-\le A\le m_+$. By Lemma~\ref{cubic},
$C(\rho-\rho_B)\ge sAQ\ge sAQ_*$.
If $Q_*\ge0$, $AQ_*\ge m_-Q_*$; if $Q_*<0$,
$AQ_*\ge m_+Q_*$. This proves $C(\rho-\rho_B)\ge sG_*$
in both cases. One must use the \emph{upper} multiplier in the negative
case.

The normal part is nonpositive when $\delta\le0$.
For $\delta>0$, differentiation with respect to an intermediate variance
$t\in[0,s]$ gives
\[
 \frac{d}{dt}\left[-\sqrt n\,
 \Phi\!\left(\frac{\delta}{\sqrt{n(v+t)}}\right)\right]
 =\frac{\delta}{2(v+t)^{3/2}}
   \phi\!\left(\frac{\delta}{\sqrt{n(v+t)}}\right)
 \le\frac{.4\delta}{2v_*^{3/2}}\le P_*D_+.
\]
Integrate over $[0,s]$, and combine with the moment bound.
Finally $s=\lambda/n$ and $\kappa\ge0$, so
$(P_*D_+-G_*)\lambda/n\le\kappa\lambda/n\le\kappa\lambda/N$.
The clamp at zero is needed before replacing $n$ by its lower bound.
\end{proof}

\subsection{Exact checks and the finite application}
The new LEFT42 normalization uses
\[
                    N=42,\qquad c_*=.4,\qquad c^*=.4226.
\]
These are hypotheses supplied by the surrounding finite-index proof.
The rational normalization checker verifies multiplier inequalities by
\[
 (m_-)^2(V^*)^3<c_*^2,\qquad
 (m_+)^2v_*^3>(c^*)^2,
\]
\emph{together with positivity of both multipliers}. Squaring without
those positivity checks would accept a false negative coefficient.
The normal multiplier is checked similarly using
$4P_*^2v_*^3>.4^2$ and $P_*>0$.

The checker recomputes the minimum in \eqref{gain} from the same input
row, then verifies a rational upper enclosure of the nonnegative
coefficient in \eqref{repair_bound}. Altering the sign, the score cap,
or the acquired probability/noise row is not a licensed change of
hypotheses. Exact arithmetic verifies these finite inequalities;
Lemma~\ref{cubic} and Proposition~\ref{repair} explain why they imply the
probabilistic correction.


% Source: local42_improvements.tex
\section{Keeping phase and resource information in the local comparison}

We refine the local comparison of Section~\ref{shortloc:section} by keeping
the phase and the joint noise/displacement constraint throughout each
estimate. The nearest-count and whole-count identities, quartic minorant
and Poisson comparison were proved there; the conditional support inputs
come from Section~\ref{conf:re:confinement}. Each refinement below states
which of those acquired bounds it uses.

\subsection{Notation and the information that must be retained}
Let $Y=B+U$, where $\Prb(B=1)=p$, $\E(U\mid B)=0$,
$v=p(1-p)$, $s=\E U^2$, $\lambda=ns$, and
$\ell=(v+s)^{-1/2}$. The actual threshold has the form
\[
 x=np+H/\ell=k+u,\qquad |u|\le\tfrac12,
 \qquad \delta=k-np=H/\ell-u.
\]
The integer $k$ and the value $\delta$ stay fixed when integrating a function
of a new variable $t$ below. Write $\pi_k=\Prb(\operatorname{Bin}(n,p)=k)$ and
\[
 f(t)=\frac{\Phi((\delta+t)/\sqrt{n(v+s)})-
                   \Phi(\delta/\sqrt{n(v+s)})}{\pi_k}.
\]
The conditional errors satisfy $|U_j|\le A$, $U_0\le a_0$,
$U_1\ge-a_1$, with $a_0,a_1\ge0$ and $d=a_0+a_1$.
All quantities in an application refer to the same iid marginal.

\begin{lemma}[A joint budget on a noise interval]\label{lem:joint}
Suppose the acquired same-law resource inequality is
\[
        \frac d{d_*}+\frac\lambda\Lambda\le1,
        \qquad d_*>0,\quad\Lambda>0.
\]
On a closed noise interval $L\le\lambda\le U\le\Lambda$ one may use
\[
                d\le d_L:=d_*(1-L/\Lambda).
\]
Every conditional support, cubic-moment, tail-distance and characteristic
function bound whose hypothesis is $d\le d_0$ may therefore be applied with
$d_0=d_L$ on that interval.
\end{lemma}
\begin{proof}
Rearrange the displayed resource inequality and use $\lambda\ge L$.
The substitution is valid separately on every noise interval. Neither the
pair of inequalities $d\le d_*$, $\lambda\le\Lambda$ nor independently
maximized numerical coordinates imply the joint resource inequality.
Its acquisition is a separate obligation.
\end{proof}

For example, the whole-count kernel uses $c=1-d_0$. Its moment coefficients
are recomputed as
\[
 \alpha=6m_0\ell_-^2(1-d_0),\qquad
 \beta=\min\{3(4.16)\ell_-(1-p_+-d_0),
                  3(6.16)\ell_-(p_--d_0)\},
\]
where $m_0,\ell_-$ are the previously justified moment and scale lower
bounds. The ordinary cubic correction is recomputed with the same $d_0$.
Likewise, a noise amplitude bound $A$ gives a conditional support span at
most $A+d_0$; the characteristic-function rate is recomputed from that span.
These changes preserve the original identities and their sign conditions.

\subsection{A binomial mass bound on the actual phase cell}
Suppose the finite Stirling argument has supplied positive constants
$E,L_0,R$ on a fixed segment containing every $k/n$ under consideration, with
\[
 \sqrt n\,\pi_k=\frac{\phi(0)}{\sqrt v}
       \exp\{-\delta^2/(2nv)+\zeta\},\qquad
 n|\zeta|\le Z_N(|\delta|),\qquad
 Z_N(t)=\frac{L_0t}{2}+R+\frac{Et^3}{6N},
 \quad n\ge N.
\]
The constants may have been computed on a much wider global phase range.
Let $v_-\le v\le v_+$ and let $q_+\ge\sqrt{v_+}$.
\begin{lemma}[Mass localization]\label{lem:mass}
On any smaller closed cell where $|\delta|\le D$,
\[
 \sqrt n\,\pi_k\ge
 \frac{.39886}{q_+}
 \left[1-\frac{D^2/(2v_-)+Z_N(D)}N\right].
\]
The maximum of this expression and any previously proved mass lower bound
is also a valid mass lower bound. There is no requirement $D\ge1$.
\end{lemma}
\begin{proof}
Set $A_*=D^2/(2v_-)+Z_N(D)\ge0$. The logarithmic formula gives
$\sqrt n\pi_k\ge(.39886/q_+)e^{-A_*/n}$.
Now $e^{-z}\ge1-z$ and $n\ge N$. If the bracket is negative the displayed
inequality still holds, and the previously positive bound can be retained.
The global Stirling segment, rather than the smaller value $D$, continues
to justify $E,L_0,R$.
\end{proof}

\subsection{A sharper elementary harmonic constant}
\begin{lemma}\label{lem:cos}
For every $x\ge0$,
\[
               1-\cos(2\pi x)\le\frac{4553}{1000}x.
\]
\end{lemma}
\begin{proof}
Put $K=4553/1000$. For $x\ge2/K$ the assertion follows from
$1-\cos\le2$. On $[0,2/K]$ take the degree-fourteen alternating upper
Taylor polynomial $Q(x)$ for $1-\cos(2\pi x)$. Its remainder is nonpositive:
the first omitted term is negative, and the remaining term magnitudes
decrease, since $(4\pi/K)^2<16\cdot15$.

Use the directed endpoints
\[
 \pi_-=3.141592653589793<\pi<3.141592653589794=\pi_+.
\]
In each positive coefficient of $Q$ replace $\pi$ by $\pi_+$; in each
negative coefficient replace it by $\pi_-$. This only increases $Q$.
The exact checker expands $P(x)=K-Q(x)/x$, a polynomial of degree thirteen,
on each of the64 intervals $[2j/(64K),2(j+1)/(64K)]$. If
$P(x)=\sum_{h=0}^{13}p_hx^h$, its Bernstein coefficients on $[a,a+w]$ are
\[
 b_k=\sum_{i=0}^k\frac{\binom ki}{\binom{13}i}
        \sum_{h=i}^{13}p_h\binom hi a^{h-i}w^i,
               \qquad k=0,\ldots,13.
\]
All896 rational coefficients are positive; the smallest is greater than
$.00012312$. Thus $P>0$ on the full interval. The value at zero follows by
continuity (or directly).

For completeness the stated $\pi$ endpoints are checked using alternating
series for $\arctan(1/5)$ (32 terms) and $\arctan(1/239)$ (12 terms).
The tangent addition formula gives
$\tan(4\arctan(1/5)-\arctan(1/239))=1$.
The angle lies between zero and $4/5<\pi/2$, so it equals $\pi/4$.
The positivity follows already from
$4(1/5-1/(3\cdot5^3))-1/239>0$; the upper bound follows from
$\arctan z<z$. This proves Machin's identity on the correct branch.
The full rational recipe is \texttt{check\_harmonic4553.py}; no sampled
positivity is used.
\end{proof}

If the original characteristic-function argument gives
$|\E e^{2\pi iW}|\le e^{-rt}$, Lemma~\ref{lem:cos} upgrades its fractional
part estimate from $(5/23)(1-e^{-rt})$ to
\[
                 \frac{1000}{4553}(1-e^{-rt}).
\]
Consequently every occurrence of $5/23$ and $5/46$ in that harmonic gain is
replaced by $1000/4553$ and $500/4553$, respectively. The existing degree-six
Taylor gain and its monotonicity argument remain unchanged.

\subsection{The Gaussian derivative knows the sign of the threshold}
\begin{lemma}[Positive phase, with a fixed actual count]\label{lem:mid}
Assume $u\in[a,b]\subset[0,1/2]$, $H/\ell\le h_*$, $|\delta|\le D$,
$v\ge v_->0$, $\lambda\le U$, and the Stirling remainder bound above.
Set
\[
 G=\max\left(0,\max_{z\in[a,b]}(2h_*z-z^2)\right),\qquad
 \eta=\frac{U+G}{2v_-}+Z_N(D).
\]
Then $f(u)\ge u(1-\eta/N)$ for every $n\ge N$.
The inner maximum is attained at the projection of $h_*$ onto $[a,b]$.
If $H\le0$, one may take $h_*=0$, hence $G=0$.
\end{lemma}
\begin{proof}
For a fixed actual law and threshold, differentiate while keeping
$\delta=k-np$ fixed:
\[
 \log f'(t)=-\tfrac12\log(1+s/v)
       +\frac{\delta^2}{2nv}-\frac{(\delta+t)^2}{2n(v+s)}-\zeta.
\]
For $0\le t\le u$, use $\log(1+s/v)\le s/v$ and $v+s\ge v$ to obtain
\[
 \log f'(t)\ge-\frac{\lambda+2\delta t+t^2}{2nv}-\zeta.
\]
The convex function $t\mapsto2\delta t+t^2$ attains its maximum on
$[0,u]$ at an endpoint. Since the actual relation is $\delta=H/\ell-u$,
\[
 2\delta t+t^2\le\max(0,2(H/\ell)u-u^2)\le G.
\]
Thus $\log f'(t)\ge-\eta/n$, and
$f'(t)\ge e^{-\eta/n}\ge1-\eta/N$. Integrate from zero to the actual $u$.
This argument does not change $\delta$ when $t$ varies and does not reflect
the distribution.
\end{proof}

In the code the signed $H$ interval supplies the bound $D$. If $H\ge0$,
one can use $h_*=H_+/\ell_-$ and $D=\max(b,h_*-a)$; if $H\le0$,
use $D=H_-/\ell_-+b$. The new lower derivative constant is maximized with
the older valid constant. Every resulting bound is checked positive.

The normal variance and cubic-moment repair is localized separately.
On a phase cell, use its actual upper bound for $\delta_+$, rather than
$|\delta|$. If the resulting coefficient is $\kappa$, replace it by
$\max(0,\kappa)$ before using $\kappa/n\le\kappa/N$.
That sign condition is essential for an all-index conclusion.

\subsection{Coverage and a conservative cache}
The selected resource table has $47$ probability cells covering $[.38,.474]$.
Zero noise is handled by the Bernoulli theorem. TINY covers all phases for
$0<\lambda\le\min(.05,\Lambda)$. Above $.05$, the LOW phase radii are
$.09,.14,.19,.29,.5$ on noise bands separated by $.06,.08,.10,.12$.
MID covers the complementary positive and negative phase intervals through
$.12$. All intervals are closed at shared boundaries and all actual-event
conventions are inherited from the exact count identity.

The LOW implementation caches a curvature bound at fixed probability,
$H$ and phase intervals. The cache key in the selected version omits $c$.
This is conservative here: noise bands are traversed in increasing order,
so Lemma~\ref{lem:joint} makes $c=1-d_L$ increase. Each denominator in the
cached terms $(c-u)^{-2}$ and $(c+\min(u,c/3))^{-2}$ increases with $c$
on the checked positive domain. The numerators are nonnegative and
independent of $c$. Thus an older cache entry remains an upper bound.
Other $c$-dependent quantities are recomputed in each noise band.

Each table is applied after acquisition has supplied its joint resource
inequality and amplitude profile for the same marginal. The complete
composition checks these inputs as well as the closed table coverage.


% Source: local42_signed_gaussian_v2.tex
\section{The signed Gaussian correction in nearest-count comparison}

Let $n\ge N$, $v=p(1-p)>0$, $s\ge0$, $\lambda=ns$, and let the
actual nearest count and phase satisfy
\[
 \delta=k-np=h-u,\qquad
 f(t)=\frac{\Phi((\delta+t)/\sqrt{n(v+s)})-
                   \Phi(\delta/\sqrt{n(v+s)})}{\pi_k}.
\]
The actual $k$ and $\delta$ are held fixed when $t$ varies. Assume the
proved finite Stirling formula on this count range is
\[
 \sqrt n\,\pi_k=\frac{\phi(0)}{\sqrt v}
       \exp\{-\delta^2/(2nv)+\zeta\},\qquad
 n|\zeta|\le Z_N(D),\quad |\delta|\le D.
\]
Write $v\ge v_->0$ and $Z=Z_N(D)\ge0$.

\begin{lemma}[Negative phase]
Suppose $u=-r$, $0\le a\le r\le b$, and $h\le h_*$. Define
\[
 t_*=\max(0,\min(b,h_*+b)),\qquad G_-=2(h_*+b)t_*-t_*^2,\qquad
 \eta_-=G_-/(2v_-)+Z.
\]
If $D^2\le Nv_-$ and $\eta_-<N$, then
\[
 f'(-t)\le\frac1{1-\eta_-/N}\quad(0\le t\le r),\qquad
 u-f(u)\le\frac{b\eta_-}{n(1-\eta_-/N)}.
\]
\end{lemma}
\begin{proof}
Direct differentiation and the Stirling formula give
\[
 \log f'(-t)=-\tfrac12\log(1+s/v)
 +\frac{\delta^2s}{2nv(v+s)}
 +\frac{2\delta t-t^2}{2n(v+s)}-\zeta.
\]
Since $\log(1+x)\ge x/(1+x)$ for $x\ge0$, the first two terms are at most
\[
 \frac{s}{2(v+s)}\left(\frac{\delta^2}{nv}-1\right)\le0.
\]
This cancellation removes the variance term from the upper derivative
bound. Because $\delta=h+r$, $r\le b$ and $0\le t\le r\le b$,
\[
 2\delta t-t^2\le 2(h_*+b)t-t^2\le G_-.
\]
The last maximum is the concave quadratic maximum on $[0,b]$, achieved
at the displayed projection $t_*$. It is nonnegative because $t=0$ is
available. Therefore $\log f'(-t)\le\eta_-/n$.
Using $e^x\le(1-x)^{-1}$ for $0\le x<1$ gives
\[
 f'(-t)-1\le\frac{\eta_-/n}{1-\eta_-/n}
             \le\frac{\eta_-}{n(1-\eta_-/N)}.
\]
Integrate $u-f(u)=\int_0^r(f'(-t)-1)\,dt$.
\end{proof}

Equivalently, $G_-=b^2+2h_*b$ for $h_*\ge0$, and
$G_-=\max(0,h_*+b)^2$ for $h_*<0$. This sharper vertex is used in the
selected LEFT42 LOW implementation.

\begin{lemma}[Positive phase]
Suppose $u\in[a,b]\subset[0,1/2]$, $h\le h_*$, and define
\[
 G_+=\max\left(0,\max_{z\in[a,b]}(2h_*z-z^2)\right).
\]
Then
\[
 u-f(u)\le\frac b n
       \left(\frac\lambda{2v_-}+\frac{G_+}{2v_-}+Z\right).
\]
The inner maximum is at the projection of $h_*$ onto $[a,b]$.
\end{lemma}
\begin{proof}
The exact derivative identity and $\log(1+s/v)\le s/v$ give
\[
 \log f'(t)\ge-\frac{\lambda+2\delta t+t^2}{2nv}-\zeta.
\]
For fixed actual $u$, the function $2\delta t+t^2$ is convex on $[0,u]$.
Its maximum is at an endpoint, and $\delta=h-u$ gives
\[
 2\delta t+t^2\le\max(0,2hu-u^2)\le G_+.
\]
Use $e^{-x}\ge1-x$ and integrate from zero to $u$.
\end{proof}

\paragraph{Whole-count insertion and interval coverage.}
The whole-count identity is split as
$Q/\pi_k-f(u)=(Q/\pi_k-u)+(u-f(u))$.
Thus one needs an upper bound on the signed correction $u-f(u)$;
an absolute-value estimate is unnecessary. After division by $\lambda$,
the new coefficients replacing the old Gaussian correction are
\[
 \begin{array}{ll}
 u\ge0:&\displaystyle
 \frac b{2v_-}+\frac{b\{G_+/(2v_-)+Z\}}\lambda,\\[6pt]
 u\le0:&\displaystyle
 \frac{b\eta_-}{(1-\eta_-/N)\lambda}.
 \end{array}
\]
On each fixed phase cell these have the form $A+B/\lambda$ with
$A,B\ge0$. They therefore preserve the original convex endpoint check on
each closed noise interval. Every phase cell is confined to one half-line;
zero is a shared closed endpoint. The count displacement bound $D$ comes
from the same actual $h-u$ cell, and the Stirling constants continue to
come from the already proved global count segment.

The other all-index steps are unchanged: the event bracket is checked
negative before substituting a binomial mass lower bound, and the entire
inverse-$n$ correction is checked positive before replacing $n$ by $N$.
The inverse correction is $I(\lambda)=K+A+B/\lambda$, hence decreases with $\lambda$; positivity at both endpoints (in particular the upper endpoint) licenses the all-index step throughout the interval. The final score is convex in $\lambda$: its event fraction is a positive constant divided by an affine positive function, its $B/\lambda$ coefficient is nonnegative, and its remaining terms are affine or constant. These are independent gates in the exact kernel. The formulas apply to
both signs of $h$ without reflecting the law or changing the count during
integration. They also supply an improved upper derivative for the
negative-phase MID inequalities.
\paragraph{Full-phase MID below the whole-count threshold.}
For $\lambda\in[.05,.08]$, the selected shifted-quartic MID comparisons
are checked on every closed phase interval of length $.001$ in
$[0,1/2]$, for both phase signs. The shifted-quartic polynomial and
nearest-count identity are defined at phase zero, so the two closed covers
include $u=0$ directly. This permits the whole-count LOW comparison to
start at $\lambda=.08$. The four MID bands use phase starts
$0,0,.19,.27$ at noise bands $[.05,.06]$, $[.06,.08]$,
$[.08,.10]$, and $[.10,.12]$, respectively.


% Source: paired_extension_with_signed_k.tex
\section{A larger domain for paired Fourier cancellation}

The signed Gaussian transfer lemma in the cutoff49 formula atlas already
allows $\mu_4\le11.1$ and $|\gamma|\le1.85$. Its specialized positive
coefficient rectangle was used only for $\mu_4\le2$, $|\gamma|\le1$.
The following extension changes that specialization. Bounded support and
all numerical resources remain hypotheses to be acquired separately.

\begin{proposition}[The actual absolute third moment]
Suppose $0<l\le p\le u<1/2$, $0\le\theta\le\Theta<1$, and the actual
conditional cubic identity gives
\[
 \rho=\E|X|^3\ge(1-\theta)^{3/2}\rho_B(p)+3a\theta,
 \qquad a\ge0,\quad
 \rho_B(p)=\frac{1-2p(1-p)}{\sqrt{p(1-p)}}.
\]
Here $a$ can be the minimum of the two acquired lower support endpoints.
With outward bounds $r_-\le\rho_B(u)$ and $r_+\ge\rho_B(l)$, a valid
lower bound is
\[
 r_*:=\max\{1,r_--(3r_+/2-3a)_+\Theta\}\le\rho.
\]
\end{proposition}
\begin{proof}
The function $\rho_B$ decreases on $(0,1/2)$ and convexity gives
$(1-\theta)^{3/2}\ge1-3\theta/2$. Therefore
\[
 \rho\ge\rho_B(p)-(3\rho_B(p)/2-3a)\theta
 \ge r_--(3r_+/2-3a)_+\Theta.
\]
The bound $\rho\ge1$ follows from $\E X^2=1$ and monotonicity of the
absolute moment norms. No separate requirement $3a\ge3r_+/2$ is needed.
\end{proof}

\begin{proposition}[Positive cancellation rectangle beyond $M=2$]
Let $\E X=0$, $\E X^2=1$, $|X|\le R$ almost surely, and let
\[
 1\le\mu_4\le M<3,\qquad
 0\le\gamma_l\le|\gamma|\le g,\qquad g^2<2,\quad R>2g,
 \quad M-1-\gamma_l^2\ge0.
\]
The paired formulas of the formula atlas remain valid, with exactly one
change to their damping constant:
\[
                 b_0=\frac{45-M}{96}(1-2/N),\qquad N>2.
\]
In particular, after simultaneous reflection of both odd Fourier
coefficients, the valid rectangle for
$a_0=(3-\mu_4)/24$, $z=(10\gamma-\E X^5)/120$,
$c=a_0\gamma/6$ is
\begin{gather*}
 a_*=(2-\gamma_l^2)/24,\qquad
 c_l=(3-M)\gamma_l/144,\quad
 c_u=(2-\gamma_l^2)g/144,\\
 t_* =\min\{g,(10-2M)/(2(R-g))\},\\
 120Z=\min\{E_5+10g,R(M-1)+(10-2M)t_*-(R-g)t_*^2\},\\
 z_l=\max\{-Z,[8\gamma_l-\gamma_l^3
        -(R+2\gamma_l)(M-1-\gamma_l^2)]/120\},\quad z_u=Z,
\end{gather*}
where $E_5\ge\E|X|^5$. If $3g^2\le2$, $c_u$ may also be decreased
to $\min(c_u,(2g-g^3)/144)$.
The same sixth-moment bound remains valid:
\[
 M_6=\min\{RE_5,1+3g^2+g^4+
 (R^2+2Rg+3g^2+2)(M-1-\gamma_l^2)\}.
\]
\end{proposition}
\begin{proof}
Pearson's identity gives $\mu_4\ge1+\gamma^2$, so $0\le a_0\le a_*$.
Both stated bounds on $c$ follow by multiplying nonnegative quantities.
The optional refinement uses the monotonicity of $2t-t^3$ on
$[0,g]$ when $3g^2\le2$.

Set $Q=X^2-\gamma X-1$. The exact identities
\[
 \E Q^2=\mu_4-1-\gamma^2,\qquad
 \E X^5=\gamma^3+2\gamma+\E[Q^2(X+2\gamma)]
\]
give the upper absolute fifth-coefficient bound by maximizing
\[
 R(M-1)+(10-2M)t-(R-g)t^2,\qquad0\le t\le g.
\]
The larger of the two Pearson sign bounds is still this one: $M<3<5$,
and $R>2g$ permits replacing $\mu_4$ by its upper bound. The quadratic
has its maximum at the displayed clipped vertex. The numerator in the
lower bound for $z$ has derivative
$10-2M+2Rt+3t^2>0$, so its lower endpoint is obtained at $\gamma_l$.
The sixth-moment identity
\[
 \mu_6=1+3\gamma^2+\gamma^4+
 \E[Q^2\{(X+\gamma)^2+2\gamma^2+2\}]
\]
gives the stated $M_6$.

The signed Gaussian transfer is applicable because $M<3<11.1$ and
$g<\sqrt2<1.85$. Its norm is convex in $(z,c)$, so the rectangle is
checked at its four corners; all Gaussian and Mills estimates remain
unchanged. Its damping hypothesis is
$b_0\le(45-M)(1-2/N)/96$. The older value
$43(1-2/N)/96$ is unavailable when $M>2$ and must be replaced as stated.
\end{proof}

\paragraph{A necessary implementation seam.}
The same new $r_*$ must be used in the four Fourier corners, the signed
phase-curvature correction, all conserved budgets, and exported quadratic
resources. The latter correction remains
\[
 q(C,\rho)=C\rho+\tfrac12\min\{\phi_{\rm low}g_0,.4g_0\},
 \qquad g_0=3.16\rho-3-d.
\]
Recomputing the old $\rho_B$ lower endpoint in any of these locations would
invalidate the rectangle association. The new adapter returns $r_*$ as
part of its paired result and the corrected callback uses that value.

\subsection*{The same lower bound improves signed fourth-moment errors}
Let $V=\E(X\mid\text{sign class})$ and $W=X-V$. In the actual central-gap
domain there is no zero atom. Write $\varepsilon=\operatorname{sgn}X$,
$x=|X|$, $a=|V|$, $\beta=\E[X|X|]$, and
\[
 K_1(X)=\E[X|X|^3]-3\beta-\tfrac32\rho\gamma.
\]
Conditional Taylor expansion gives the exact identity
\[
 K_1(X)-K_1(V)=\E W^2\left\{
 \varepsilon(w_4-3-\tfrac32\rho w_3)
                  -\tfrac32\gamma_V w_3\right\},
\]
where $w_3=x+2a$, $w_4=x^2+2ax+3a^2$ and $\gamma_V=\E V^3$.
For the product term use
$\rho\gamma-\rho_V\gamma_V
=\rho(\gamma-\gamma_V)+\gamma_V(\rho-\rho_V)$.
Every other term follows by expanding the signed or unsigned second,
third and fourth power about the conditional mean; the linear remainder
has conditional expectation zero.

Define
\[
 F(x,a;t)=x^2+2ax+3a^2-3-\tfrac32t(x+2a).
\]
The positive-class coefficient is $F(x,a;\rho+\gamma_V)$; the
negative-class coefficient is $-F(x,a;\rho-\gamma_V)$.
For $x,a>0$, $F$ decreases with $t$. Thus any proved lower bound
$r_*\le\rho$ may be used in the upper endpoint for the positive
coefficient and lower endpoint for the negative coefficient. The proof
does not need the stronger assertion $\rho\ge\rho_B$.

On each actual conditional mean/support rectangle, $F$ is convex with
Hessian $\left(\begin{smallmatrix}2&2\\2&6\end{smallmatrix}\right)$.
Its maximum is at a corner. Its minimum is attained at a corner, at an
edge stationary point, or at $(x,a)=(3t/8,3t/8)$ if that point is inside.
The edge points are $a=t/2-x/3$ at fixed $x$ and
$x=3t/4-a$ at fixed $a$, clipped to their respective edges. These formulas
allow negative $t$ as well.

If $L$ and $U$ bound the two signed coefficients together, then
\[
 \Theta\min(0,L)\le K_1(X)-K_1(V)\le\Theta\max(0,U),
\]
since $\E W^2\le\Theta$. The enclosing $p$ and conditional-mean intervals
are the same ones used by the formula atlas. Its outer condition
$1-\Theta>2/3$, which licenses the separate binary $K$ monotonicity
calculation, is retained. Intersect this signed bound with the existing
unsigned estimate and reject an empty intersection.

The implementation in \texttt{paired\_k\_extension} uses the proved
$r_*$ lower endpoint in this signed rectangle calculation. It retains
the stated moment domain, the outer condition $1-\Theta>2/3$, and
all intersection checks. The acquisition step must supply these premises
before the improved fourth-moment bound is used.

